\section{Construction order by order}
\label{ch:order-by-order}

This chapter carries out the Hubbard--Stratonovich transformation of closed string field theory one order in $\gs$ at a time, from the action of the theory to the octic order.
Its core is the full order-by-order section of the paper in its earlier version: the action, then the quartic, quintic, sextic, septic and octic orders.
The paper's sections on the quartic to the septic order are the condensed version of this chapter, and the octic order, which the paper only summarizes, is kept here in full.
Three sections of a companion note on the quintic order follow: the completion of the square with the dressed kernel, the orders at which the corrected saddle enters, and the proof that the kernel condition is necessary as well as sufficient.
A section of an older note on the two presentations of the $3$-string field at the sextic order comes next.
Two sections exist only here: the four-point amplitude of the quartic order in full, and the gauge laws and the products of the sextic order in both presentations.

\medskip

The first two blocks are the introduction of the paper's order-by-order section and its recall of the non-polynomial action, in the paper's notation, which the whole chapter uses except inside the blocks of the two older notes (each preceded by a translation box).

\subsection{Hubbard–Stratonovich transformation order by order}
\label{p:sec:hs}

A Hubbard--Stratonovich transformation~\cite{Stratonovich:1957:MethodCalculatingQuantum,Hubbard:1959:CalculationPartitionFunctions} trades an interaction that is the square of some composite, such as $\lambda f(\phi)^2$, for a Gaussian integral over an auxiliary field coupled linearly to that composite.
For example, the quartic vertex $\lambda \phi^4$ becomes a cubic coupling $\sigma \phi^2$ plus a quadratic term for the auxiliary field, and integrating the auxiliary field out gives back the original interaction exactly (see \cref{p:app:scalar-hs}).
In closed string field theory the interaction at order $\gs^{n-2}$ is the vertex $\vertex{n}$ evaluated on $n$ copies of the string field.
The natural composites are thus the tensor powers $\Psi^k := \Psi^{\otimes k}$, and the auxiliary fields replacing them are $k$-string fields $\Sigma_k \in \Hn{k}$.
This section carries out the transformation order by order in $\gs$, introducing one auxiliary field at a time.
The purpose is to describe the different ingredients which are needed for rewriting the action at all orders (\cref{p:sec:all-orders}).

\Cref{p:sec:hs:action} recalls the action of closed string field theory and fixes the notations.
The quartic order (\cref{p:sec:hs:quartic}) introduces $\Sigma_2$ and is worked out in full.
At the quintic order (\cref{p:sec:hs:quintic}), $\Sigma_2$ alone is still sufficient, but its saddle-point solution and its gauge transformation receive corrections.
At the sextic order (\cref{p:sec:hs:sextic}), a $3$-string field $\Sigma_3$ is introduced, together with the reduced sextic vertex: the sextic vertex minus the exchange diagram induced by $\Sigma_2$.
The septic and octic orders (\cref{p:sec:hs:septic,p:sec:hs:octic}) make the general pattern visible.

We treat the quartic order in more detail than its simplicity requires, because every ingredient of the construction already appears there.
We start by describing the equation of motion which identifies the saddle-point of the auxiliary field with the tensor power $\Psi^2$, and the kernel of the quartic vertex with the quotient on which the auxiliary field lives.
We then show in detail how to eliminate the auxiliary field by completing a square.
We describe different properties of the field: its gauge transformation, its reality condition and its spacetime content.
Last come the Feynman rules through which the exchange of the auxiliary field reproduces the quartic vertex.
Each of these points returns, with corrections in $\gs$, at every higher order.


\subsubsection{The non-polynomial action}
\label{p:sec:hs:action}

The classical action of closed string field theory~\cite{Zwiebach:1993:ClosedStringField} is a sum of a kinetic term and of one interaction per number of strings:
\begin{equation}
  \label{p:eq:hs1:action}
  S
    = \frac{1}{2} \, \pair{\Psi}{Q_B \Psi}
      + \sum_{n \ge 3} \frac{\gs^{n-2}}{n!} \, \vertex{n}(\Psi^n),
\end{equation}
where $\gs$ is the string coupling and $\Psi^n := \Psi^{\otimes n}$.

The string field $\Psi \in \Hs$ is an element of the space of states of the matter--ghost CFT satisfying the level-matching conditions
\begin{equation}
  \label{p:eq:hs1:level-matching}
  b_0^- \ket{\Psi}
    = 0,
  \qquad
  L_0^- \ket{\Psi}
    = 0,
\end{equation}
and carrying ghost number
\begin{equation}
  \label{p:eq:hs1:ghost-psi}
  \Ngh(\Psi)
    = 2.
\end{equation}
The pairing is the BPZ inner product with one insertion of $c_0^-$:
\begin{equation}
  \label{p:eq:hs1:pairing}
  \pair{A}{B}
    := \braket{A | c_0^- | B}.
\end{equation}

The maps $\vertex{n} \colon \Hn{n} \to \C$ are the genus-zero string vertices.
They are obtained by integrating the correlation function of the $n$ states over the vertex region $\mathcal V_{0,n} \subset \mc M_{0,n}$ of the moduli space of the $n$-punctured sphere, with the appropriate antighost insertions.
The string products $\ell_n \colon \Hn{n} \to \Hs$ are the maps they define through the pairing:
\begin{equation}
  \label{p:eq:hs1:vertex-product}
  \vertex{n}(A_1, \ldots, A_n)
    =: \pair{A_1}{\ell_{n-1}(A_2, \ldots, A_n)},
  \qquad
  \ell_1
    := Q_B.
\end{equation}
We follow the conventions of~\cite{Erbin:2021:StringFieldTheory} for the vertices and the products.
They are collected in \cref{p:app:conventions}, together with the normalization of the sphere correlators, the ghost numbers and the reality condition.

Here $\Hn{n}$ denotes the space of $n$-string states whose components are graded-symmetric under permutations of the strands.
For the Grassmann-even strands of $\Psi^n$ and $\Sigma_n$, it is the $n$-th tensor power of $\Hs$ symmetrized with the Grassmann signs~\cite{Erler:2014:NSNSSectorClosed, Erbin:2020:ClassicalAlgebraicStructures}.
The odd strands of the outputs of the products are treated in \cref{p:app:conventions}.
The vertices are graded-symmetric in their arguments.
Ghost-number conservation on the sphere requires the insertions to total $\Ngh = 6$, which fixes
\begin{equation}
  \label{p:eq:hs1:ghost-vertices}
  \Ngh(\vertex{n})
    = 6 - 2n,
  \qquad
  \Ngh(\ell_n)
    = 3 - 2n.
\end{equation}
The vertices and products are extended by multilinearity to tensor arguments, by expanding a tensor on a basis of $\Hn{n}$ built from a basis of $\Hs$.
We write the tensors in the slots they fill, as in $\vertex{4}(\Sigma_2, \Psi^2)$ for $\Sigma_2, \Psi^2 \in \Hn{2}$.

We write the vertices, maps $\Hn{n} \to \C$, with a capital letter, $\vertex{n}$, and the string products with a lowercase one.
A product with $m$ inputs and $n$ outputs, $\Hn{m} \to \Hn{n}$, carries both numbers as subscripts, $\sprod{m}{n}$, except when they equal $0$ or $1$: the products $\ell_n = \sprod{n}{1}$ have one output, the vertices $\vertex{n} = \sprod{n}{0}$ have none, while the surface states $\bra{\mc V_n}$ have no input.
An object that changes from one order in $\gs$ to the next carries the highest power of $\gs$ kept in the action as a parenthesized superscript.
A field living in $\Hn{n}$ carries the index $n$.

The construction of this section rests on reading a vertex with some of its slots as outputs.
The simplest instance is the quartic vertex read as a bilinear form on $2$-string states.
We define the product $\sprod{2}{2} \colon \Hn{2} \to \Hn{2}$ by
\begin{equation}
  \label{p:eq:hs1:p22}
  \vertex{4}(A_1, A_2, A_3, A_4)
    =: \pair[2]{A_1 \otimes A_2}{\sprod{2}{2}(A_3, A_4)},
  \qquad
  \pair[2]{\cdot}{\cdot}
    := \pair{\cdot}{\cdot}^{\otimes 2},
\end{equation}
where the $2$-string pairing has one $c_0^-$ insertion per strand.
The product $\sprod{2}{2}$ raises the ghost number by two, $\Ngh(\sprod{2}{2}) = 2$.
Since $\vertex{4}$ is symmetric, the product $\sprod{2}{2}$ is symmetric with respect to the $2$-string pairing:
\begin{equation}
  \label{p:eq:hs1:p22-symmetric}
  \pair[2]{A}{\sprod{2}{2}(B)}
    = \pair[2]{\sprod{2}{2}(A)}{B}.
\end{equation}
In the same way, $\ell_3 = \sprod{3}{1}$ is the quartic vertex read with one output.
The general case $\sprod{m}{n}$, its ghost number $3n - 2m$ for $n \ge 1$ (a vertex $\sprod{m}{0}$ keeps the ghost number $6 - 2m$ of~\eqref{p:eq:hs1:ghost-vertices}), and the transposition identity relating $\sprod{m}{n}$ to $\sprod{n}{m}$ are given in \cref{p:app:conventions}.

We denote by $S_n$ the term of order $\gs^{n-2}$ in~\eqref{p:eq:hs1:action}, $S_2$ being the kinetic term, and by
\begin{equation}
  \label{p:eq:hs1:truncation}
  S
    = \sum_{n \ge 2} S_n,
  \qquad
  S_{\le N}
    := \sum_{n = 2}^N S_n,
\end{equation}
the truncation of the action at order $\gs^{N-2}$.
The construction proceeds truncation by truncation: at each order we look for an action, cubic in the fields, which reproduces $S_{\le N}$ once the auxiliary fields are eliminated.

We write the variation of the action through the pairing, $\delta S = \pair{\delta \Psi}{\mathcal E_\Psi}$, so that the equation of motion of the original theory is $\mathcal E_\Psi = 0$, with:
\begin{equation}
  \label{p:eq:hs1:eom-psi}
  \mathcal E_\Psi
    := Q_B \Psi
      + \sum_{n \ge 2} \frac{\gs^{n-1}}{n!} \, \ell_n(\Psi^n),
\end{equation}
the term of order $\gs^{n-1}$ being the variation of $S_{n+1}$.

The action~\eqref{p:eq:hs1:action} is invariant, order by order in $\gs$, under the gauge transformation~\cite{Zwiebach:1993:ClosedStringField}
\begin{equation}
  \label{p:eq:hs1:gauge-psi}
  \delta_\Lambda \Psi
    = Q_B \Lambda
      + \sum_{n \ge 2} \frac{\gs^{n-1}}{(n-1)!} \, \ell_n(\Psi^{n-1}, \Lambda),
\end{equation}
with $\Lambda \in \Hs$ of ghost number one.
This invariance is the content of the $\Loo$ relations satisfied by the products $\ell_n$, which we take as given with the vertices.

The string field satisfies the reality condition
\begin{equation}
  \label{p:eq:hs1:reality-psi}
  \Psi^\star
    = - \Psi,
\end{equation}
where $\star$ is the star conjugation, the composition of the BPZ and Hermitian conjugations, $A^\star := \adj{\transp{A}}$, as recalled in \cref{p:app:conventions}.


The quartic order introduces the $2$-string field $\Sigma_2$ by completing a square, derives the equations of motion and the quotient by the kernel of $\sprod{2}{2}$, eliminates $\Sigma_2$, gives the gauge transformation and the reality condition, writes $\Sigma_2$ in momentum space and reads off the Feynman rules.
It states the following results:
\begin{itemize}
  \item the cubic action and its equation of motion, which gives $\Sigma_2 = \Psi^2$ on the quotient when $\sprod{2}{2}$ is invertible there: \status{Verified}\record{001-cubic-sft-overview-note (R1)};
  \item $\Sigma_2$ is purely auxiliary if and only if $\sprod{2}{2}$ is injective on the $2$-string states: \status{Verified}\record{004-i3-m2-invertibility (R9)};
  \item integrating $\Sigma_2$ out gives back the quartic action: \status{mixed}\record{005-s6-integrate-out (R10)};
  \item the gauge transformation of $\Sigma_2$ makes the variations of the cubic and of the original action equal, exactly and off-shell: \status{Proven}\footnote{Ledger entry \recordref{R11}, with no derivation record of its own.};
  \item the $\Sigma_2$ exchange collapses to the quartic contact term with coefficient one: \status{Conditional(PH1) -- Verified}\record{007-a1-amplitude-telescoping (R14)}.
\end{itemize}

\subsubsection{Quartic order, \texorpdfstring{$O(\gs^2)$}{O(gs**2)}}
\label{p:sec:hs:quartic}


\subparagraph{The auxiliary field.}

The quartic term of the action reads:
\begin{equation}
  S_4
    = \frac{\gs^2}{4!} \, \vertex{4}(\Psi^2, \Psi^2).
\end{equation}
It is the square of the $2$-string state $\Psi^2$ with respect to the bilinear form $\vertex{4}$, in the same way as $\lambda \phi^4$ is the square of $\phi^2$.
A Hubbard--Stratonovich transformation trades a square for a linear coupling to a new field, at the cost of a quadratic term for that field.
We introduce an auxiliary $2$-string field $\Sigma_2 \in \Hn{2}$ and add to the action the square of its fluctuation from $\Psi^2$:
\begin{equation}
  \label{p:eq:hs4:DeltaS4}
  \Delta S_4
    := - \frac{\gs^2}{4!} \, \vertex{4}(\Sigma_2 - \Psi^2, \Sigma_2 - \Psi^2).
\end{equation}
This term vanishes when $\Sigma_2 = \Psi^2$.
We put a tilde on every action that contains a term $\Delta S_n$ vanishing on-shell, and write $\caction[n] := S_n + \Delta S_n$.
Expanding the square and using the symmetry of $\vertex{4}$ for the cross term, we find the quartic term of the action:
\begin{equation}
  \label{p:eq:hs4:S4tilde}
  \caction[4]
    := S_4 + \Delta S_4
    = - \frac{\gs^2}{4!} \, \vertex{4}(\Sigma_2^2)
      + \frac{2 \gs^2}{4!} \, \vertex{4}(\Sigma_2, \Psi^2),
\end{equation}
where $\vertex{4}(\Sigma_2^2) := \vertex{4}(\Sigma_2, \Sigma_2)$.
The term $\vertex{4}(\Psi^4)$ has disappeared: the first term is quadratic in $\Sigma_2$, and the second one is cubic in the fields, with one $\Sigma_2$ and two $\Psi$.
Together with the kinetic term and the cubic vertex of~\eqref{p:eq:hs1:action}, which are untouched, the action
\begin{equation}
  \label{p:eq:hs4:action-tilde}
  \caction[\le 4]
    := S_2 + S_3 + \caction[4]
    = S_{\le 4} + \Delta S_4
\end{equation}
is cubic in $(\Psi, \Sigma_2)$.

We could redistribute the powers of $\gs$ in~\eqref{p:eq:hs4:S4tilde} by absorbing $\gs$ in the auxiliary field, $\Sigma_2 \to \Sigma_2 / \gs$.
The quadratic term then carries no coupling and the mixed term a single power of $\gs$, like the cubic vertex $S_3$, which makes its cubic nature manifest.
We nevertheless keep $\gs$ explicit, so that the solution for $\Sigma_2$ is $\Psi^2$ without a coupling in front and the two actions can be compared order by order in $\gs$ term by term.

\subparagraph{Equations of motion.}

We vary the cubic action~\eqref{p:eq:hs4:action-tilde} with respect to both fields and write each term through the pairing:
\begin{equation}
  \label{p:eq:hs4:delta-S}
  \begin{aligned}
    \delta \caction[\le 4]
      = \pair{\delta \Psi}{Q_B \Psi}
        &
        + \frac{\gs}{2} \, \pair{\delta \Psi}{\ell_2(\Psi^2)}
        + \frac{\gs^2}{3!} \, \pair{\delta \Psi}{\ell_3(\Psi, \Sigma_2)}
        \\ &
        - \frac{2 \gs^2}{4!} \, \pair[2]{\delta \Sigma_2}{\sprod{2}{2}(\Sigma_2)}
        + \frac{2 \gs^2}{4!} \, \pair[2]{\delta \Sigma_2}{\sprod{2}{2}(\Psi^2)},
  \end{aligned}
\end{equation}
where $\ell_3(\Psi, \Sigma_2)$ is the product $\ell_3$ with $\Sigma_2$ filling two of its three slots.
Since the pairing is non-degenerate, the equations of motion are:
\begin{subequations}
  \label{p:eq:hs4:eom}
  \begin{align}
    \label{p:eq:hs4:eom-psi}
    \mathcal E^{(2)}_\Psi
      &
      := Q_B \Psi
        + \frac{\gs}{2} \, \ell_2(\Psi^2)
        + \frac{\gs^2}{3!} \, \ell_3(\Psi, \Sigma_2)
      = 0,
    \\
    \label{p:eq:hs4:eom-sigma}
    \mathcal E^{(2)}_{\Sigma_2}
      &
      := \sprod{2}{2}(\Psi^2) - \sprod{2}{2}(\Sigma_2)
      = 0.
  \end{align}
\end{subequations}
The equation for $\Sigma_2$ contains no $Q_B$: it is algebraic, and it says that $\Sigma_2$ and $\Psi^2$ have the same image under $\sprod{2}{2}$.

This equation determines $\Sigma_2$ up to the kernel of $\sprod{2}{2}$, which is the kernel of the quartic vertex seen as a bilinear form on $2$-string states:
\begin{equation}
  \label{p:eq:hs4:kernel}
  N_2 \in \ker \sprod{2}{2}
  \quad \Longleftrightarrow \quad
  \forall A \in \Hn{2}:
  \quad
  \vertex{4}(N_2, A)
    = 0.
\end{equation}
The two statements are equivalent since the $2$-string pairing is non-degenerate.
This kernel can contain entangled $2$-string states, so it is in general larger than the space spanned by the products $A \otimes B$ with $A$ a single-string state annihilated by $\vertex{4}(A, \cdot, \cdot, \cdot)$.
A kernel direction $N_2$ appears neither in the quadratic term $\vertex{4}(\Sigma_2^2)$ nor in the source term $\vertex{4}(\Sigma_2, \Psi^2)$ of~\eqref{p:eq:hs4:S4tilde}.
Shifting $\Sigma_2$ by $N_2$ leaves the action invariant, with zero action for $N_2$ itself, so such a direction drops out of every equation and of every amplitude at this order.

We therefore take the auxiliary field on the quotient:\footnotemark{}
\footnotetext{
  In the same way, the auxiliary field associated with stubs lives outside the kernel of $L_0$ or $L_0^+$~\cite{Erbin:2024:OpenStringStub}.
}
\begin{equation}
  \label{p:eq:hs4:quotient}
  \Sigma_2
    \in \Hn{2} / \ker \sprod{2}{2},
  \qquad
  \Ngh(\Sigma_2)
    = 4.
\end{equation}
We restrict each strand to ghost number two, like the strands of $\Psi^2$.
On this quotient, $\sprod{2}{2}$ is injective by construction, and the equation of motion~\eqref{p:eq:hs4:eom-sigma} has the unique solution:
\begin{equation}
  \label{p:eq:hs4:saddle}
  \Sigma_2^{(2)*}
    := \Psi^2
    = \quad
      \begin{tikzpicture}[diagram]
        \draw (0, 0.1) -- (2, 0.1) pic {ext} node[above right] {$\Psi$};
        \draw (0, -0.1) -- (2, -0.1) pic {ext} node[below right] {$\Psi$};
      \end{tikzpicture}
\end{equation}
We represent the solution graphically by two strands, one for each $\Psi$ of the solution.
This description will be helpful at higher orders.

Taking the auxiliary field in \eqref{p:eq:hs4:quotient} is also consistent with the form of the action: single-string states that never enter the quartic vertex have no quartic term and need no auxiliary field.
This is illustrated in \cref{p:app:scalar-hs} for the $\phi^4$ theory.
What is problematic is the consistency with composite vertices at higher orders.
As a consequence, it is still important to understand if there is a kernel and what are its properties, which is the topic developed below.
The $\phi^4$ model of \cref{p:app:scalar-hs} with a degenerate coupling matrix illustrates this situation, and also the multiplier form, which we use on a kernel direction when a higher coupling sources it.

The quotient needs no inner product, and the $2$-string pairing cannot provide one: it vanishes identically between two fields of ghost number four, so it defines no orthogonal complement of the kernel inside the space where $\Sigma_2$ lives.
A complement of the kernel appears only in the multiplier form of \cref{p:sec:all-orders}.

The solution~\eqref{p:eq:hs4:saddle} shows that $\Sigma_2$ is an auxiliary field: it is fixed entirely by $\Psi$ through an algebraic equation and carries no degree of freedom of its own.
It also shows that $\Delta S_4$ vanishes on-shell, since $\Delta S_4$ is the square of $\Sigma_2 - \Psi^2$.
Eliminating $\Sigma_2$ therefore gives back the truncated action exactly.
Substituting the solution into~\eqref{p:eq:hs4:S4tilde}, we get:
\begin{equation}
  \label{p:eq:hs4:eliminate}
  \caction[4] \big|_{\Sigma_2 = \Psi^2}
    = \Big( - \frac{1}{4!} + \frac{2}{4!} \Big) \, \gs^2 \, \vertex{4}(\Psi^4)
    = S_4.
\end{equation}
This is the completed square~\eqref{p:eq:hs4:action-tilde} read backwards, with no extra term.
Since the cubic action is quadratic in $\Sigma_2$, with the quadratic form $\sprod{2}{2}$ independent of $\Psi$, the same holds in the path integral: the Gaussian integral over $\Sigma_2$ is formally exact whatever the signature of the quartic form.
From the quintic order on, the kinetic operators of the auxiliary fields depend on $\Psi$, and so do their determinants; the eliminations are exact at tree level, to which we restrict ourselves from then on.
It contributes only a field-independent normalization, so that, formally, all the Green's functions of $\Psi$ computed from $\caction[\le 4]$ coincide with those of $S_{\le 4}$.

Physically, the auxiliary field is a contact field: its equation of motion is built from an invertible operator with no zeros.
Whether that inverse is singular at some momentum depends on the quartic form itself, and this question enters the invertibility assumption stated below with the Feynman rules.
In particular, the latter is not the usual kinetic operator $L_0^+$, so it has no pole at the mass of any string state and describes no propagating state.
Its ``propagator'' is the inverse of a vertex rather than of $L_0^+$, as we make explicit below.
The auxiliary field changes the arrangement of the interaction: the $4$-string contact term of closed string field theory is now produced by two cubic couplings joined by a $2$-string line.


\subparagraph{Gauge transformation.}

Under the gauge transformation~\eqref{p:eq:hs1:gauge-psi}, the truncated action is invariant up to the order of the truncation, $\delta_\Lambda S_{\le 4} = O(\gs^3)$, the remainder being cancelled by the quintic vertex.
We look for a transformation of $\Sigma_2$ under which the cubic action is invariant to the same order.
Since $\Sigma_2 = \Psi^2$ on-shell, the natural candidate is the transformation of $\Psi^2$ itself:
\begin{equation}
  \label{p:eq:hs4:gauge-sigma}
  \delta_\Lambda \Sigma_2
    = 2 \, \delta_\Lambda \Psi \otimes \Psi.
\end{equation}
We do not symmetrize it, since it only appears inside a symmetric vertex.
Inserting~\eqref{p:eq:hs4:gauge-sigma} into the variation~\eqref{p:eq:hs4:delta-S}, we find that the two terms in $\delta \Sigma_2$ become:
\begin{equation}
  \begin{gathered}
    - \frac{\gs^2}{3!} \, \pair[2]{\delta \Psi \otimes \Psi}{\sprod{2}{2}(\Sigma_2)}
    + \frac{\gs^2}{3!} \, \pair[2]{\delta \Psi \otimes \Psi}{\sprod{2}{2}(\Psi^2)}
      \\
      = - \frac{\gs^2}{3!} \, \pair{\delta \Psi}{\ell_3(\Psi, \Sigma_2)}
        + \frac{\gs^2}{3!} \, \vertex{4}(\delta \Psi, \Psi^3),
  \end{gathered}
\end{equation}
where we used that both sides are the same quartic vertex read with different outputs:
\begin{equation}
  \label{p:eq:hs4:same-vertex}
  \pair[2]{\delta \Psi \otimes \Psi}{\sprod{2}{2}(A)}
    = \vertex{4}(\delta \Psi, \Psi, A)
    = \pair{\delta \Psi}{\ell_3(\Psi, A)}.
\end{equation}
The first term cancels the term $\pair{\delta \Psi}{\ell_3(\Psi, \Sigma_2)}$ of~\eqref{p:eq:hs4:delta-S}, and the second one is the variation of $S_4$.
We obtain:
\begin{equation}
  \label{p:eq:hs4:delta-S-equal}
  \delta \caction[\le 4]
    = \pair{\delta \Psi}{Q_B \Psi}
      + \frac{\gs}{2} \, \pair{\delta \Psi}{\ell_2(\Psi^2)}
      + \frac{\gs^2}{3!} \, \vertex{4}(\delta \Psi, \Psi^3)
    = \delta S_{\le 4}.
\end{equation}
This equality holds identically, off-shell, for any variation $\delta \Psi$ and not only for a gauge transformation.
The cubic action with the induced transformation~\eqref{p:eq:hs4:gauge-sigma} is thus invariant to the same order $O(\gs^3)$ as the truncated action.
We did not use the solution~\eqref{p:eq:hs4:saddle} itself, only the transformation it suggests.

We give another approach which will be more convenient at higher orders.
This equality follows more directly from the completed square: by~\eqref{p:eq:hs4:action-tilde}, $\delta \caction[\le 4] = \delta S_{\le 4} + \delta \Delta S_4$, and
\begin{equation}
  \label{p:eq:hs4:delta-DeltaS}
  \delta \Delta S_4
    = - \frac{2 \gs^2}{4!} \, \vertex{4}\big(\delta \Sigma_2 - \delta (\Psi^2), \Sigma_2 - \Psi^2\big)
    = 0,
\end{equation}
since $\delta (\Psi^2) = \delta \Psi \otimes \Psi + \Psi \otimes \delta \Psi$ equals $2 \, \delta \Psi \otimes \Psi$ inside the symmetric vertex.
The transformation law~\eqref{p:eq:hs4:gauge-sigma} is the one under which the fluctuation $\sigma_2^{(2)} := \Sigma_2 - \Psi^2$ is invariant.
The transformation of $\Sigma_2$ is induced by that of $\Psi$, and no independent parameter enters the on-shell matching.

Before the quotient is taken, the cubic action also has the shift symmetry along the kernel directions, which have zero action, and the transformations proportional to its equations of motion that every action possesses.
We do not look for others here.
Whether the cubic action can be seen as a gauge-fixed form of a theory with an independent $2$-string parameter $\Lambda_2 \in \Hn{2}$ is a different question, which we take up in \cref{p:sec:gauge}.
At the next order, the transformation law~\eqref{p:eq:hs4:gauge-sigma} receives corrections in $\gs$, as does the solution~\eqref{p:eq:hs4:saddle}.


\subparagraph{Reality.}

The string field satisfies the reality condition~\eqref{p:eq:hs1:reality-psi}, $\Psi^\star = - \Psi$.
The star conjugation acts strand by strand on tensors.
We assume, as in \cref{p:app:conventions}, that the vertices are real:
\begin{equation}
  \label{p:eq:hs4:vertex-real}
  \vertex{n}(A_1, \ldots, A_n)^*
    = (-1)^n \, \vertex{n}(A_1^\star, \ldots, A_n^\star)
  \qquad
  \text{for even arguments}.
\end{equation}
Together with $\Psi^\star = - \Psi$, this property makes the action~\eqref{p:eq:hs1:action} real.
The solution~\eqref{p:eq:hs4:saddle} then shows that the natural reality condition on the auxiliary field is:
\begin{equation}
  \label{p:eq:hs4:reality}
  \Sigma_2^\star
    = \Sigma_2.
\end{equation}
Indeed, $(\Psi \otimes \Psi)^\star = \Psi^\star \otimes \Psi^\star = \Psi \otimes \Psi$ for a string field satisfying its reality condition, and the reality of the vertices makes the cubic action real for real $\Psi$ and $\Sigma_2$.
We derive~\eqref{p:eq:hs4:reality} and the reality conditions of the higher auxiliary fields in \cref{p:app:reality}.


\subparagraph{Spacetime fields.}

The string field $\Psi$ expands on a basis $\ket{\phi_\alpha(k)}$ of single-string states of momentum $k$ with coefficients $\psi_\alpha(k)$, the spacetime fields in momentum space.
The auxiliary field expands in the same way on $2$-string states.
Since a general element of $\Hn{2}$ is an entangled state, its coefficient is a function of two momenta:
\begin{equation}
  \label{p:eq:hs4:spacetime}
  \Sigma_2
    = \sum_{\alpha_1, \alpha_2} \int \frac{\dd^d k_1}{(2\pi)^d} \, \frac{\dd^d k_2}{(2\pi)^d} \,
      \sigma_{\alpha_1 \alpha_2}(k_1, k_2) \,
      \ket{\phi_{\alpha_1}(k_1)} \otimes \ket{\phi_{\alpha_2}(k_2)},
\end{equation}
where $d$ is the spacetime dimension.
The coefficient is symmetric under the joint exchange of the two strands, $(\alpha_1, k_1) \leftrightarrow (\alpha_2, k_2)$.
The components $\sigma_{\alpha_1 \alpha_2}(k_1, k_2)$ are bilocal fields in momentum space, and thus also in position space, where they are functions of two spacetime points.

Every vertex conserves only the total momentum of its strands: its single momentum delta function comes from the zero mode of the matter CFT.
We introduce the total and relative momenta:
\begin{equation}
  \label{p:eq:hs4:momenta}
  P
    := k_1 + k_2,
  \qquad
  q
    := \frac{1}{2} \, (k_1 - k_2).
\end{equation}
The quadratic term is thus diagonal in $P$:
\begin{equation}
  \label{p:eq:hs4:diagonal}
  \vertex{4}(\Sigma_2, \Sigma_2')
    = \int \frac{\dd^d P}{(2\pi)^d}
      \int \frac{\dd^d q}{(2\pi)^d} \frac{\dd^d q'}{(2\pi)^d} \;
      \sigma(- P, q') \, V_4(P; q', q) \, \sigma'(P, q),
\end{equation}
where we suppressed the oscillator labels and where $V_4(P; q', q) := V_4\big(- \tfrac{P}{2} + q', - \tfrac{P}{2} - q', \tfrac{P}{2} + q, \tfrac{P}{2} - q\big)$ is the momentum-space quartic vertex on two pairs of states with total momenta $- P$ and $P$.
At fixed total momentum, the quadratic term is thus an integral operator in the relative momentum.
We define such expressions using the Sen's generalized Wick rotation~\cite{Pius:2016:CutkoskyRulesSuperstring,deLacroix:2017:ClosedSuperstringField,Erbin:2021:StringFieldTheory,Bautista:2019:QuantumGravityTimelike}.
The theory is defined with Euclidean momenta, and the Lorentzian theory is obtained from it by a Wick rotation afterwards.
We keep the total momentum $P$ as a real label.
The operator acts on the square-integrable functions of the relative momentum at fixed total momentum.
Throughout, a kernel of $\sprod{2}{2}$ at fixed total momentum is meant in this sense: its elements are square-integrable functions of the relative momentum, valued in the oscillator labels.
We now restrict each strand to oscillator level at most $L$, as in level truncation.
At level $L$, $V_4(P; q', q)$ is then the integral over the compact vertex region of an exponential of a real quadratic form in $(P, q, q')$, the same as at level $0$, times a polynomial in the momenta whose degree grows with the level.
Since the region avoids the collisions of the punctures, $V_4(P; q', q)$ is an entire function of the total momentum and of both relative momenta, rather than a delta function times a polynomial.
The kinetic term of $\Sigma_2$ is therefore not local in the relative coordinate.
The Gaussian decay in $q$ and $q'$ produced by the stubs~\cite{Sonoda:1990:ClosedStringField}, when its rate is positive, dominates any polynomial.
  As a consequence, $\sprod{2}{2}$ is a Hilbert--Schmidt, hence compact, operator at every level, under this condition, which makes the level-$0$ kernel square integrable (\cref{p:sec:worldsheet}).

The $2$-string states carry no canonical positive inner product above level $0$, but counting the null vectors of $\sprod{2}{2}$ is nevertheless a self-adjoint problem.
  Reflecting the relative momentum along the direction of $P$, together with the matching reflection of the oscillator labels, defines a fixed involution.
  At real spacelike $P$, $\sprod{2}{2}$ at level $L$ composed with this involution is a real symmetric compact operator.
  For timelike $P$ in its rest frame, $\sprod{2}{2}$ itself is self-adjoint.
  At every real $P^2$, of either sign, the null vectors can therefore be counted with the spectral theorem.
  A null vector is defined by~\eqref{p:eq:hs4:kernel}, whatever basis or inner product is used.

The source term $\vertex{4}(\Sigma_2, \Psi^2)$ is of the same form, with $\sigma'(P, q)$ replaced by the product $\psi_{\alpha_1}(\tfrac{P}{2} + q) \, \psi_{\alpha_2}(\tfrac{P}{2} - q)$ of two components of $\Psi$.
The component of $\Sigma_2$ at total momentum $P$ is sourced by every pair of components of $\Psi$ whose momenta add up to $P$.

A translation-invariant configuration of $\Sigma_2$, such as a translation-invariant classical solution or a vacuum expectation value, is proportional to $\delta^d(k_1 + k_2)$ and depends on the relative momentum only~\cite{Weinberg:2005:QuantumTheoryFields-2}.
It is then described by a single momentum, like a local field, though with the non-local kinetic term just described.
This is the $n = 2$ case of a general statement: the $n$-string auxiliary field introduced at higher orders is a function of one total momentum and $n - 1$ relative momenta, and its quadratic term is diagonal in the total momentum.
Its translation-invariant configurations depend on the $n - 1$ relative momenta only, that is, on $n$ points modulo an overall translation.
We discuss the form of the kernel $V_4(P; q', q)$ and its worldsheet interpretation in \cref{p:sec:worldsheet}.


\subparagraph{Feynman rules.}

We read the Feynman rules off the action $\caction[\le 4]$, the path-integral weight being $\e^{\caction}$.
A vertex factor is the derivative of the action with respect to the fields on its legs, and a propagator is minus the inverse of the Hessian of the action.
A tree diagram then gives the corresponding derivative of the action with the internal fields eliminated.
The auxiliary sector of~\eqref{p:eq:hs4:S4tilde} contributes a propagator and a cubic vertex:
\begin{equation}
  \label{p:eq:hs4:feynman-rules}
  \frac{12}{\gs^2} \, \sprod{2}{2}^{-1},
  \qquad
  \frac{\gs^2}{3!} \, \vertex{4}(\Sigma_2, \Psi^2).
\end{equation}



The next block closes the paper's paragraph on the Feynman rules: the four-point function with identical external states, the collapse of the $\Sigma_2$ line, and the assumption on the intermediate $2$-string states under which the collapse holds for any external states.


We now consider the four-point function with four identical external states $\Psi$ at order $\gs^2$.
The auxiliary sector contributes three channels, $s$, $t$ and $u$, each with two cubic vertices and one $\Sigma_2$ line.
Since the external states are identical, the three channels are equal and we get:
\begin{equation}
  \label{p:eq:hs4:telescoping}
  \begin{aligned}
    3 \times
    \Big( \frac{\gs^2}{3!} \, \vertex{4}(\Psi^2, \cdot) \Big)
    \circ
    \Big( \frac{12}{\gs^2} \, \sprod{2}{2}^{-1} \Big)
    \circ
    \Big( \frac{\gs^2}{3!} \, \vertex{4}(\cdot, \Psi^2) \Big)
      &
      = \gs^2 \, \vertex{4}\Big( \Psi^2, \sprod{2}{2}^{-1} \circ \sprod{2}{2}(\Psi^2) \Big)
      \\ &
      = \gs^2 \, \vertex{4}(\Psi^4),
  \end{aligned}
\end{equation}
where the second vertex is written through the pairing, $\vertex{4}(\cdot, \Psi^2) = \pair[2]{\cdot}{\sprod{2}{2}(\Psi^2)}$, and where $\sprod{2}{2}^{-1} \circ \sprod{2}{2}$ is the identity on the quotient.
We call \emph{collapse} the fact that a $\Sigma_2$ line ending on a source vertex~\eqref{p:eq:hs4:feynman-rules} reduces to a contact term, $\sprod{2}{2}^{-1} \circ \sprod{2}{2}(\Psi^2) = \Psi^2$.
The exchange~\eqref{p:eq:hs4:telescoping} is the diagrammatic form of the elimination~\eqref{p:eq:hs4:eliminate}.
We give the worldsheet interpretation of the exchange in \cref{p:sec:worldsheet}.

The Feynman rules of the $\Psi$ sector, the propagator $- b_0^+ / L_0^+$ in Siegel gauge and the cubic vertex $\vertex{3}$, are those of closed string field theory.
The $\Sigma_2$ exchange collapses as in~\eqref{p:eq:hs4:telescoping} for any external states: the two pair states filling the vertices replace $\Psi^2$.
Hence every tree-level Green's function of $\Psi$ computed from the cubic action equals its counterpart in the quartic-truncated theory, for the same vertex family and off-shell.
Besides the invertibility of $\sprod{2}{2}$ on the quotient assumed above, we need one more assumption.
The composition of vertex, propagator and vertex is a sum over a complete set of intermediate $2$-string states, which runs over the oscillator levels of the two strands and over the relative momentum of the pair.
We assume that this sum converges on the $2$-string states the external sources produce.


\paragraph{The four-point amplitude in full}
\label{sec:order:amplitude}

\status{mixed}\record{023-a4pt-virasoro-shapiro-chain (R29)}
The statements below have different statuses: the enumeration of the diagrams, the prefactors and the collapse in a finite-dimensional model \status{Proven}; the equality of the off-shell amplitudes \status{Conditional(PH1, PH3) -- Proven}; the on-shell Virasoro--Shapiro amplitude \status{Conditional(PH1, PH3, PH4) -- Proven}; the on-shell limit at level~$0$ \status{Conditional(PH4, PH25, PH29, PH30) -- Proven}.

The block above shows the collapse for four identical external states.
We now compute the whole four-point amplitude at order $\gs^2$ on four distinct external states, with every prefactor, and follow it to the on-shell amplitude.
The computation is at tree level, for a fixed consistent family of vertices: the argument does not depend on which family.
It uses three assumptions, which we state once here.
The first is the one made in the quartic order: $\sprod{2}{2}$ descends to an injective map on the quotient of the $2$-string states by its kernel, invertible there at each truncation level, and the kernels of the level-truncated forms stabilize as the level grows, so that the quotient and the propagator of $\Sigma_2$ have a limit; at the external momenta used, including the on-shell ones, the truncated determinant on the quotient has no zero that persists with the level.
The second is the one stated at the end of the block above: the sum over intermediate $2$-string states, which runs over the levels of the two strands and over the relative momentum of the pair, is complete and converges on the states $\sprod{2}{2}(P)$ produced by an external pair $P$.
The third is needed only on-shell: the three regions of the moduli space $\mathcal M_{0, 4}$ covered by the Feynman diagrams with one propagator, together with the region of the quartic vertex, cover $\mathcal M_{0, 4}$ exactly once.
No other input enters, beyond the defining properties of the vertex family, among which the total symmetry of $\vertex{4}$ under permutations of its four slots.

\subparagraph{Normalization.}
The path-integral weight is $\e^{S}$, so that every propagator is minus the inverse of the kinetic form.
For four distinct external states $\Psi_1, \ldots, \Psi_4$ we define the amplitude as the connected four-point function with amputated external legs,
\begin{equation}
  \label{eq:ch2:amp:normalization}
  A_4(\Psi_1, \ldots, \Psi_4)
    := \partial_{t_1} \partial_{t_2} \partial_{t_3} \partial_{t_4} \log Z[J(t)] \Big|_{t = 0},
  \qquad
  J(t)
    = \vertex{2} \Big( \sum_i t_i \Psi_i, \cdot \Big),
\end{equation}
where $\vertex{2}$ is the kinetic form of $\Psi$, so that the source $J$ cancels the external propagators.
The coefficient of $t_1 t_2 t_3 t_4$ selects the part multilinear in four distinct states.
In this normalization every tree diagram on distinct external states has symmetry factor one, and the factorials of the action are absorbed by the sums over the assignments of the legs; the rational factors that remain are shown explicitly below.
On-shell, $Q_B \Psi_i = 0$, the amplitude~\eqref{eq:ch2:amp:normalization} is the tree-level S-matrix element up to the normalization of the external states, which is common to the cubic theory and to closed string field theory.

\subparagraph{Feynman rules at order $\gs^2$.}
The rules read off the cubic action through the quartic order are:
\begin{itemize}
  \item the propagator $\mathsf P$ of $\Psi$, of order $\gs^0$, the same as in closed string field theory;
  \item the propagator of $\Sigma_2$, $\tfrac{12}{\gs^2} \, \sprod{2}{2}^{-1}$, of order $\gs^{-2}$, from the kinetic term $- \tfrac{\gs^2}{4!} \vertex{4}(\Sigma_2, \Sigma_2)$;
  \item the cubic vertex, from $\tfrac{\gs}{3!} \vertex{3}(\Psi^3)$, which on distinct external states reads $\gs \, \vertex{3}(\Psi_i, \Psi_j, \cdot)$;
  \item the source vertex, from $\tfrac{2 \gs^2}{4!} \vertex{4}(\Sigma_2, \Psi^2)$, which on distinct external states reads $\tfrac{\gs^2}{6} \, \vertex{4}(\cdot, \Psi_i \Psi_j)$ for $i \ne j$, the factor two coming from the two assignments of the external states to the slots.
\end{itemize}
There is no self-coupling of $\Sigma_2$ at this order: $\Sigma_2$ enters the quartic-order action quadratically and linearly only.
The propagator of $\Sigma_2$ is defined with respect to the identity on the symmetric $2$-string states.

\subparagraph{The diagrams.}
At order $\gs^2$ the connected four-point amplitude receives contributions from exactly four diagrams: the three trees with two cubic vertices joined by a $\Psi$ line, in the channels $s$, $t$ and $u$, and the tree with two source vertices joined by a $\Sigma_2$ line, in its three matchings of the external states.
No loop, tadpole or mixed diagram contributes.
This follows from counting.
Let $n_3$ be the number of cubic vertices, $n_B$ the number of source vertices, $I_\Psi$ and $I_\Sigma$ the numbers of internal lines of $\Psi$ and $\Sigma_2$, and $E = 4$ the number of external states.
Since $\Sigma_2$ has no external leg and each $\Sigma_2$ line ends on two source vertices, $n_B = 2 I_\Sigma$.
The order in $\gs$ of a diagram is then
\begin{equation}
  n_3 + 2 n_B - 2 I_\Sigma
    = n_3 + 2 I_\Sigma
    = 2,
\end{equation}
so that $(n_3, I_\Sigma)$ is $(2, 0)$ or $(0, 1)$, for any connected diagram, loops included.
Counting the legs of $\Psi$, $3 n_3 + 2 n_B = 4 + 2 I_\Psi$, fixes $I_\Psi = 1$ in the first case and $I_\Psi = 0$ in the second, and the number of loops
\begin{equation}
  L
    = (I_\Psi + I_\Sigma) - (n_3 + n_B) + 1
\end{equation}
vanishes in both cases: tree topology is forced by the order and the number of external states.
The first case gives two cubic vertices and one $\Psi$ line in three matchings, the channels $s$, $t$, $u$; the second gives two source vertices and one $\Sigma_2$ line in three matchings.
A tadpole of $\Psi$ would need a loop at this order, a tadpole of $\Sigma_2$ vanishes since the Gaussian of $\Sigma_2$ has zero mean, the disconnected pieces are removed by the logarithm in~\eqref{eq:ch2:amp:normalization}, and the Gaussian determinant of $\Sigma_2$ does not depend on the fields and cancels in the normalized correlator.

\subparagraph{The assembled amplitude.}
With $\{ij | kl\}$ running over the three matchings $\{12 | 34\}$, $\{13 | 24\}$, $\{14 | 23\}$, the amplitude is
\begin{equation}
  \label{eq:ch2:amp:assembled}
  A_4^{\text{cubic}}
    = \sum_{\{ij | kl\}} \gs \, \vertex{3}(\Psi_i, \Psi_j, \cdot) \; \mathsf P \; \gs \, \vertex{3}(\cdot, \Psi_k, \Psi_l)
      + \sum_{\{ij | kl\}} \frac{\gs^2}{6} \, \vertex{4}(\cdot, \Psi_i \Psi_j) \; \frac{12}{\gs^2} \, \sprod{2}{2}^{-1} \; \frac{\gs^2}{6} \, \vertex{4}(\cdot, \Psi_k \Psi_l),
\end{equation}
with no other rational factor.
The first sum is identical, diagram by diagram, in the cubic theory and in the quartic-truncated action of closed string field theory, $S_{\le 4} = \tfrac12 \pair{\Psi}{Q_B \Psi} + \tfrac{\gs}{3!} \vertex{3}(\Psi^3) + \tfrac{\gs^2}{4!} \vertex{4}(\Psi^4)$, since the two theories share every rule of the $\Psi$ sector.
The whole comparison is thus between the second sum and the contact term $\gs^2 \, \vertex{4}(\Psi_1, \Psi_2, \Psi_3, \Psi_4)$ of $S_{\le 4}$.

\subparagraph{The collapse, step by step.}
Take one matching and write $P$ and $Q$ for the symmetrized external pairs $\Psi_i \Psi_j$ and $\Psi_k \Psi_l$; only the symmetric parts couple, by the symmetry of $\vertex{4}$.
Writing the second vertex through the pairing, $\vertex{4}(\cdot, Q) = \pair[2]{\cdot}{\sprod{2}{2}(Q)}$, we get
\begin{equation}
  \label{eq:ch2:amp:collapse}
  \begin{aligned}
    A_4^{(ij | kl)}
      &
      = \frac{\gs^2}{6} \, \vertex{4}(\cdot, P) \circ \frac{12}{\gs^2} \, \sprod{2}{2}^{-1} \circ \frac{\gs^2}{6} \, \vertex{4}(\cdot, Q)
    \\
      &
      = \frac{\gs^2}{3} \, \vertex{4} \Big( P, \sprod{2}{2}^{-1} \circ \sprod{2}{2}(Q) \Big)
    \\
      &
      = \frac{\gs^2}{3} \, \vertex{4}(P, Q)
    \\
      &
      = \frac{\gs^2}{3} \, \vertex{4}(\Psi_1, \Psi_2, \Psi_3, \Psi_4).
  \end{aligned}
\end{equation}
The second line composes the three factors, which is a sum over a complete set of intermediate $2$-string states $\lambda_I$, normalized so that $\sum_I \lambda_I \pair[2]{\lambda_I}{\cdot}$ is the identity.
This is where the second assumption acts: the partial sums
\begin{equation}
  \label{eq:ch2:amp:partial-sums}
  s_N(P, Q)
    = \sum_{I \le N} \vertex{4} \big( P, \sprod{2}{2}^{-1} \lambda_I \big) \, \pair[2]{\lambda_I}{\sprod{2}{2}(Q)}
\end{equation}
must converge to $\vertex{4}(P, Q)$, and this convergence, which allows the sum to be exchanged with $\sprod{2}{2}^{-1}$, is the only analytic input of the collapse.
A level truncation makes the sum over levels finite but does not truncate the integral over the relative momentum of the pair.
That integral leaves no trace in the result because what is inserted is a resolution of the identity on the whole space of pairs, momentum included, which is saturated against the neighbouring vertex.
The third line is where the first assumption acts: $Q$ enters modulo the kernel of $\sprod{2}{2}$, the kernel components are not sourced and decouple, and $\vertex{4}(P, \cdot)$ does not see them.
Nothing stronger is needed: a solution of $\sprod{2}{2}(X) = \sprod{2}{2}(Q)$ exists, $X = Q$, it is unique modulo the kernel by definition of the quotient, and what the first assumption supplies is the uniformity in the level of the quotient on which $\sprod{2}{2}^{-1}$ acts, not a bound on the inverse.
The last line uses the total symmetry of $\vertex{4}$, which makes the result the same for the three matchings; symmetry under the exchange of the two pairs alone would leave them unequal.
Summing the three matchings,
\begin{equation}
  \label{eq:ch2:amp:sum}
  A_4^{\Sigma_2\text{-exchange}}
    = 3 \times \frac{\gs^2}{3} \, \vertex{4}(\Psi_1, \Psi_2, \Psi_3, \Psi_4)
    = \gs^2 \, \vertex{4}(\Psi_1, \Psi_2, \Psi_3, \Psi_4),
\end{equation}
which is the contact term of $S_{\le 4}$ with coefficient one, the chain of factors being $3 \times \tfrac16 \times 12 \times \tfrac16 = 1$.
Each of the three matchings reproduces the whole quartic vertex region with weight one third: the factor three is a symmetry count, not a tiling of the moduli space like the disjoint regions of the channels $s$, $t$, $u$.

\subparagraph{Off-shell and on-shell.}
For arbitrary external states, off-shell, and the same vertex family, we therefore have at order $\gs^2$
\begin{equation}
  \label{eq:ch2:amp:off-shell}
  A_4^{\text{cubic}}(\Psi_1, \ldots, \Psi_4)
    = A_4^{\text{SFT, off-shell}}(\Psi_1, \ldots, \Psi_4),
\end{equation}
under the first two assumptions.
Off-shell amplitudes depend on the vertex family and on the choice of local coordinates, and~\eqref{eq:ch2:amp:off-shell} says that the cubic rewriting adds no dependence of its own: the two theories have the same off-shell Green's functions, not only the same S-matrix.
Since $\Sigma_2$ enters the quartic-order action quadratically, the same holds for all the Green's functions of $\Psi$ (the Gaussian identity of generating functionals, at tree level in the block above).

To go on-shell, four remarks are needed.
First, the collapse~\eqref{eq:ch2:amp:collapse} is algebraic and multilinear in the external states and uses no genericity, so it holds wherever both sides are defined, the on-shell locus included.
Second, the poles of the channels are common to both sides, since the channels $s$, $t$, $u$ are the same diagrams in both theories, so no limit has to be exchanged with the collapse.
Third, the first two assumptions depend on the external momenta, through $\vertex{4}$ and through the states $\sprod{2}{2}(P)$, and they are to be read as including the on-shell momenta.
Fourth, the amputation in~\eqref{eq:ch2:amp:normalization} degenerates on-shell, where $Q_B \Psi_i = 0$ kills the numerator while the external propagator diverges; this degeneration cancels because the statement is relative, both theories amputating with the same source and the same external propagators, and an absolute on-shell statement would need an LSZ argument that we do not give.
The support for the first assumption at the on-shell momenta is genericity: $\sprod{2}{2}$ has no kernel at generic off-shell momentum for a generic vertex family, and no mechanism ties a degeneration to the tachyon shell.
This support is weak: genericity over families is a remedy rather than evidence, and the on-shell locus lies outside the scope of the genericity statement, which is made off-shell.
The failure mode is concrete, though: the level-truncated determinant of $\vertex{4}$ has zeros with no known reason to avoid a given locus, and a zero on the tachyon shell that persists with the level would break the on-shell statement.

For four on-shell tachyons, in $D = 26$, the tree amplitude of closed string field theory is the integral over $\mathcal M_{0, 4}$ given the third assumption, which is standard for the minimal-area and hyperbolic vertices~\cite{Zwiebach:1993:ClosedStringField,Erbin:2021:StringFieldTheory} and, for the $\group{SL}(2, \C)$ vertices, supported by numerical checks at a few on-shell points and on the boundary but not established; that integral is the Virasoro--Shapiro amplitude~\cite{Polchinski:2005:StringTheory-1}:
\begin{equation}
  \label{eq:ch2:amp:vs}
  A_4^{\text{SFT}} \Big|_{\text{on-shell}}
    = A_{\text{VS}}(s, t, u)
    \propto \frac{\Gamma(-1 - \alpha' s / 4) \, \Gamma(-1 - \alpha' t / 4) \, \Gamma(-1 - \alpha' u / 4)}{\Gamma(2 + \alpha' s / 4) \, \Gamma(2 + \alpha' t / 4) \, \Gamma(2 + \alpha' u / 4)},
\end{equation}
with $s = - (k_1 + k_2)^2$ and so on in the signature $(-, +, \ldots, +)$.
For four tachyons, $m^2 = - 4 / \alpha'$, so that $s + t + u = - 16 / \alpha'$: the three arguments of the numerator sum to one and each argument of the denominator is one minus the corresponding argument of the numerator, the canonical form $\prod_i \Gamma(a_i) / \prod_i \Gamma(1 - a_i)$ with $\sum_i a_i = 1$.
Combining with~\eqref{eq:ch2:amp:off-shell}, the cubic action reproduces the on-shell four-tachyon Virasoro--Shapiro amplitude under the three assumptions.
What is shown is the reduction to the standard amplitude of closed string field theory, which is known to equal the Virasoro--Shapiro amplitude on-shell; the Virasoro--Shapiro amplitude is not derived again, and the covering property of the vertex family is assumed, not tested.

At level~$0$, on the block of two tachyons, the first two assumptions are replaced by properties of the kinetic operator of $\Sigma_2$ established in \cref{ch:momentum-space}, for the vertices and under the prescription and the assumptions stated there.
The $\Sigma_2$ exchange then equals the restriction of an entire function of the external momenta at every real external pair, off-shell included, so that its on-shell limit is its on-shell value; the truncations of the eigenvalue sum, with threshold off the spectrum at the on-shell point, converge uniformly on compact subsets of the Euclidean on-shell region $\alpha' s < 16$, so that truncating and going on-shell commute.
The propagator alone has no value at the external data and does not enter.
Above level~$0$ the analogue is open: it needs the first two assumptions on a real neighbourhood of the on-shell locus, and the positivity used at level~$0$ does not carry over, the $2$-string pairing being indefinite there.

\subparagraph{No other diagram at four points.}
The count above uses the quartic-order action, which has no self-coupling of $\Sigma_2$.
The couplings of the higher orders, $\vertex{5}(\Sigma_2, \Sigma_2, \Psi)$, the sextic couplings and so on, do not contribute at four points either, but the reason is the number of legs, not the order.
A vertex with $a$ legs of $\Sigma_2$ and $b$ legs of $\Psi$ comes from a vertex with $n = 2 a + b$ punctures and carries $\gs^{2 a + b - 2}$.
A connected diagram with $V$ vertices, $I_\Sigma$ propagators of $\Sigma_2$ (each of order $\gs^{-2}$), $I_\Psi$ propagators of $\Psi$, $E$ external $\Psi$ and no external $\Sigma_2$ has $\sum_v a_v = 2 I_\Sigma$ and $\sum_v b_v = E + 2 I_\Psi$, so that its order is
\begin{equation}
  \sum_v (2 a_v + b_v - 2) - 2 I_\Sigma
    = E + 2 (I_\Sigma + I_\Psi - V)
    = E + 2 L - 2.
\end{equation}
Every four-point tree diagram is thus of order $\gs^2$, whatever its vertices.
Eliminating $I_\Psi$ in the same way gives, at tree level,
\begin{equation}
  E
    = \sum_v (a_v + b_v - 2) + 2,
\end{equation}
and every interaction vertex has $a_v + b_v \ge 3$, the only object with $a + b = 2$ being the kinetic term of $\Sigma_2$.
A vertex of the tower other than the source, that is with $a_v \ge 2$, or with $a_v = 1$ and $b_v \ge 3$, forces $E \ge 5$.
If $a_v \ge 2$, its legs of $\Sigma_2$ run at tree level to $a_v$ distinct other vertices, so $V \ge 3$ and $E \ge V + 2 \ge 5$.
If $a_v = 1$ and $b_v \ge 3$, it contributes at least $2$ to the sum, and the vertex at the other end of its $\Sigma_2$ line at least $1$, so $E \ge 5$.
The cheapest diagram with a tower vertex has five external legs and sits at order $\gs^3$, and the four diagrams above exhaust the four-point amplitude.
Any statement at order $\gs^3$ needs in addition the kernel condition of the quintic order, \eqref{p:eq:hs5:kernel-condition}.

\subparagraph{What a finite model can and cannot test.}
In a finite-dimensional model, a level truncation in which the relative momentum of the pair is suppressed, the collapse~\eqref{eq:ch2:amp:collapse} is an identity of matrix algebra for any symmetric $\vertex{4}$ and its pseudo-inverse: a numerical comparison of the truncated exchange with the truncated contact term cannot fail and tests nothing, and it passes as well on randomly corrupted vertices.
What such a model tests is the rest: the prefactors of~\eqref{eq:ch2:amp:assembled}, checked by expanding the whole generating functional of the joint Gaussian of $(\Psi, \Sigma_2)$ to order $\gs^2$ with no diagram as input; the ratio two to one between the source and the kinetic term of $\Sigma_2$, whose change to one to one gives exactly a quarter of the contact term; the use of the same $\vertex{4}$ in the kinetic term, the source and the contact term; the pairing; and the enumeration of the channels, since dropping one channel breaks the match.
In a level truncation of the actual theory the sum over intermediate states still contains the integral over the relative momentum, so the collapse at fixed level is not a matrix identity, and whether a comparison at fixed level carries information is open; what it could test is the convergence of that integral.

The physical content of this computation is that the auxiliary field changes nothing in the four-point amplitude, not even off-shell: the $\Sigma_2$ line contracts to a point because the propagator of $\Sigma_2$ is the inverse of the very form that appears in its source.
Two places only can make the statement fail, both in the representation of $\sprod{2}{2}^{-1}$ on the space of pairs: a kernel that does not stabilize with the level or appears at the external momenta, and a sum over intermediate pairs, levels and relative momentum, that does not converge.
On-shell the covering of the moduli space is a third, geometric, condition, which belongs to the vertex family and not to the cubic rewriting.

\medskip

The quintic order is the first at which the saddle of $\Sigma_2$ differs from $\Psi^2$: the quintic coupling dresses the kinetic operator of $\Sigma_2$, the saddle carries its inverse, and the elimination produces a six-point feedback term.
The paper's block derives the coupling, the equations of motion, the completed square, the five-point diagrams, the gauge transformation, and the kernel condition with the reality of the solution.
It states the following results:
\begin{itemize}
  \item the quintic coupling $\tfrac{\gs^3}{5!} \vertex{5}(\Sigma_2, \Sigma_2, \Psi)$ is unique, reproduces the quintic vertex on-shell, and the gauge law extends: \status{Conditional(PH1, PH5) -- Proven}\record{011-b1-quintic-milestone (R17)};
  \item the corrected saddle and the gauge law are real, under the hypotheses of the block below: \status{mixed}\record{039-reality-ogs3 (R46)}.
\end{itemize}
The completion of the square is derived in more detail in the companion sections after the octic order.

\subsubsection{Quintic order, \texorpdfstring{$O(\gs^3)$}{O(gs**3)}}
\label{p:sec:hs:quintic}

\paragraph{Quintic coupling and equations of motion}

\subparagraph{The quintic coupling.}

The next term of the action is the quintic term:
\begin{equation}
  \label{p:eq:hs5:S5}
  S_5
    = \frac{\gs^3}{5!} \, \vertex{5}(\Psi^5),
\end{equation}
where the vertex $\vertex{5}$ has five slots.
The question is how to describe this vertex in terms of auxiliary fields.
We cannot use the same trick as for the quartic order, because it requires an exact square.

Given the form of the saddle-point solution \eqref{p:eq:hs4:saddle} for $\Sigma_2$, we can replace $\Psi^4$ by $\Sigma_2^2$ to turn the quintic term into a cubic coupling.
Two auxiliary fields $\Sigma_2$ fill four of the five slots and the string field fills the last one:
\begin{equation}
  \label{p:eq:hs5:S5tilde}
  \caction[5]
    := \frac{\gs^3}{5!} \, \vertex{5}(\Sigma_2^2, \Psi),
  \qquad
  \caction[\le 5]
    := \caction[\le 4] + \caction[5],
\end{equation}
where $\vertex{5}(\Sigma_2^2, \Psi) := \vertex{5}(\Sigma_2, \Sigma_2, \Psi)$.
At this stage this should be treated as an ansatz: we will now check that eliminating $\Sigma_2$ gives back $S_{\le 5}$.
Note that the partition $5 = 2 + 2 + 1$ of the five strands over three fields drawn from $\{\Sigma_2, \Psi\}$ is the only one, so no other cubic term can be built from the quintic vertex with these two fields.
In principle we could use a $3$-string field to partition $5 = 3 + 1 + 1$, but it is not needed before the sextic order: we want to keep the field content minimal and avoid introducing fields from higher orders into lower orders.
Such a term also appeared when decomposing the quintic term of open SFT with stubs~\cite{Erbin:2024:OpenStringStub}.

Ghost number is conserved as it must: the two $\Sigma_2$ and the $\Psi$ carry $4 + 4 + 2 = 10$, the same total as $\Psi^5$.
Since the vertex is read with two of its slots as outputs, we define the product $\sprod{3}{2} \colon \Hn{3} \to \Hn{2}$ by
\begin{equation}
  \label{p:eq:hs5:p32}
  \vertex{5}(A_1, A_2, A_3, A_4, A_5)
    =: \pair[2]{A_1 \otimes A_2}{\sprod{3}{2}(A_3, A_4, A_5)},
\end{equation}
as $\sprod{2}{2}$ was defined from the quartic vertex, with the $2$-string pairing on the left.
It has ghost number zero, by the general formula~\eqref{p:eq:conv:ghost-numbers}.

Since the cubic term $\caction[5]$ contains $\Sigma_2$, it modifies the equation of motion of $\Sigma_2$, whose solution is no longer $\Psi^2$.
The terms of $\caction[\le 5]$ which contain $\Sigma_2$ are:
\begin{equation}
  \label{p:eq:hs5:F}
  \caction[4] + \caction[5]
    = - \frac{\gs^2}{4!} \, \vertex{4}(\Sigma_2^2)
      + \frac{2 \gs^2}{4!} \, \vertex{4}(\Sigma_2, \Psi^2)
      + \frac{\gs^3}{5!} \, \vertex{5}(\Sigma_2^2, \Psi),
\end{equation}
a quadratic form in $\Sigma_2$ with two terms, and a source term linear in $\Sigma_2$.
The kinetic term and the cubic vertex of $\Psi$ are the same in $\caction[\le 5]$ and in $S_{\le 5}$ and drop out of every comparison.


\subparagraph{Equations of motion.}

Varying the cubic action~\eqref{p:eq:hs5:S5tilde} and writing each term through the pairing, we get:
\begin{equation}
  \label{p:eq:hs5:delta-S}
  \begin{aligned}
    \delta \caction[\le 5]
      &
      = \pair{\delta \Psi}{Q_B \Psi}
        + \frac{\gs}{2} \, \pair{\delta \Psi}{\ell_2(\Psi^2)}
        + \frac{\gs^2}{3!} \, \pair{\delta \Psi}{\ell_3(\Psi, \Sigma_2)}
        + \frac{\gs^3}{5!} \, \pair{\delta \Psi}{\ell_4(\Sigma_2^2)}
        \\ &
        \quad
        - \frac{2 \gs^2}{4!} \, \pair[2]{\delta \Sigma_2}{\sprod{2}{2}(\Sigma_2)}
        + \frac{2 \gs^2}{4!} \, \pair[2]{\delta \Sigma_2}{\sprod{2}{2}(\Psi^2)}
        + \frac{2 \gs^3}{5!} \, \pair[2]{\delta \Sigma_2}{\sprod{3}{2}(\Psi, \Sigma_2)}.
  \end{aligned}
\end{equation}
The two new terms are the variation of~\eqref{p:eq:hs5:S5tilde} with respect to $\Psi$ and to $\Sigma_2$.
The equations of motion are then:
\begin{subequations}
  \label{p:eq:hs5:eom}
  \begin{align}
    \label{p:eq:hs5:eom-psi}
    \mathcal E^{(3)}_\Psi
      &
      := Q_B \Psi
        + \frac{\gs}{2} \, \ell_2(\Psi^2)
        + \frac{\gs^2}{3!} \, \ell_3(\Psi, \Sigma_2)
        + \frac{\gs^3}{5!} \, \ell_4(\Sigma_2^2)
      = 0,
    \\
    \label{p:eq:hs5:eom-sigma}
    \mathcal E^{(3)}_{\Sigma_2}
      &
      := \sprod{2}{2}(\Psi^2)
        - \sprod{2}{2}(\Sigma_2)
        + \frac{\gs}{5} \, \sprod{3}{2}(\Psi, \Sigma_2)
      = 0.
  \end{align}
\end{subequations}
The equation for $\Sigma_2$ is still algebraic, but it no longer states that $\Sigma_2$ and $\Psi^2$ have the same image under $\sprod{2}{2}$: the quintic vertex adds a term linear in $\Sigma_2$ and in $\Psi$.
We group this term with the quadratic form, since both are linear in $\Sigma_2$, and define the operator:
\begin{equation}
  \label{p:eq:hs5:dressed}
  \kin{2}{\Psi}{3}
    := \sprod{2}{2} - \frac{\gs}{5} \, \sprod{3}{2}(\Psi, \cdot),
  \qquad
  \kin{2}{\Psi}{3} \colon \Hn{2} \to \Hn{2},
\end{equation}
which we call the \emph{dressed kinetic operator} of $\Sigma_2$.
The quintic vertex, with one slot filled by $\Psi$, is not a kinetic term: it dresses the kinetic operator $\sprod{2}{2}$, as interaction insertions dress a propagator.
The superscript records that it is the operator of the action truncated at order $\gs^3$, since it receives further corrections at the next orders.
The dressed operator is symmetric with respect to the $2$-string pairing, $\pair[2]{A}{\kin{2}{\Psi}{3}(B)} = \pair[2]{\kin{2}{\Psi}{3}(A)}{B}$, since both vertices are symmetric, and it is affine in $\Psi$.
In terms of it, the quadratic part of~\eqref{p:eq:hs5:F} is a single form:
\begin{equation}
  \label{p:eq:hs5:F-dressed}
  - \frac{\gs^2}{4!} \, \vertex{4}(\Sigma_2^2)
    + \frac{\gs^3}{5!} \, \vertex{5}(\Sigma_2^2, \Psi)
    = - \frac{\gs^2}{4!} \, \pair[2]{\Sigma_2}{\kin{2}{\Psi}{3}(\Sigma_2)}.
\end{equation}
The equation of motion of $\Sigma_2$ then reads:
\begin{equation}
  \label{p:eq:hs5:eom-sigma-dressed}
  \mathcal E^{(3)}_{\Sigma_2}
    = \sprod{2}{2}(\Psi^2) - \kin{2}{\Psi}{3}(\Sigma_2)
    = 0.
\end{equation}
We can make an analogy with non-linear sigma models: the quadratic form of $\Sigma_2$ has a $\Psi$-dependent metric, and the field $\Psi$ enters the kinetic term of $\Sigma_2$ as a background.


\paragraph{Saddle-point solution and induced interactions}

\subparagraph{Saddle-point solution.}

On the quotient, where $\sprod{2}{2}$ is injective, the dressed operator is $\sprod{2}{2}$ composed with a perturbation of the identity:
\begin{equation}
  \label{p:eq:hs5:dressed-factor}
  \kin{2}{\Psi}{3}
    = \sprod{2}{2} \circ \Big( 1 - \frac{\gs}{5} \, \sprod{2}{2}^{-1} \circ \sprod{3}{2}(\Psi, \cdot) \Big),
\end{equation}
provided that $\sprod{3}{2}(\Psi, \cdot)$ takes its values in the image of $\sprod{2}{2}$ and descends to the quotient.
This is an assumption on the quintic vertex, stated precisely in~\eqref{p:eq:hs5:kernel-condition} below, where the two standing assumptions of the construction from this order on are collected.
For small $\gs$, the inverse of the dressed operator then exists as a Neumann series:
\begin{equation}
  \label{p:eq:hs5:neumann}
  \big( \kin{2}{\Psi}{3} \big)^{-1}
    = \sum_{n \ge 0} \Big( \frac{\gs}{5} \, \sprod{2}{2}^{-1} \circ \sprod{3}{2}(\Psi, \cdot) \Big)^n \circ \sprod{2}{2}^{-1},
\end{equation}
which converges for $\gs$ small at fixed $\Psi$ when the composite $\sprod{2}{2}^{-1} \circ \sprod{3}{2}(\Psi, \cdot)$ is bounded on the quotient.
Order by order in $\gs$, no invertibility condition appears beyond the one on $\sprod{2}{2}$ assumed in \cref{p:sec:hs:quartic}.
The statements at finite $\gs$ below assume instead that $\kin{2}{\Psi}{3}$ itself is invertible on the quotient.
The equation of motion~\eqref{p:eq:hs5:eom-sigma-dressed} then has the unique solution:
\begin{equation}
  \label{p:eq:hs5:saddle}
  \Sigma_2^{(3)*}
    := \big( \kin{2}{\Psi}{3} \big)^{-1} \circ \sprod{2}{2}(\Psi^2)
    = \Psi^2 + \frac{\gs}{5} \, \big( \kin{2}{\Psi}{3} \big)^{-1} \circ \sprod{3}{2}(\Psi^3),
\end{equation}
where the second form follows from $\sprod{2}{2}(\Psi^2) = \kin{2}{\Psi}{3}(\Psi^2) + \tfrac{\gs}{5} \, \sprod{3}{2}(\Psi, \Psi^2)$, the definition~\eqref{p:eq:hs5:dressed} applied to $\Psi^2$.
The solution is again fixed entirely by $\Psi$ through an algebraic equation, so that $\Sigma_2$ remains an auxiliary field.
It is no longer $\Psi^2$ but $\Psi^2$ plus corrections at all orders in $\gs$.
The solution now contains more terms, and we could ask whether the correction of order $\gs$ spoils the matching, but this is not the case by construction.
By stationarity of the quartic-order action at $\Psi^2$, the action is quadratic in any fluctuation around it, and the quintic term, of order $\gs^3$, changes linearly in the correction of order $\gs$, so the correction to the eliminated action starts at order $\gs^4$, as we will see below.

Applying $\sprod{2}{2}^{-1}$ to~\eqref{p:eq:hs5:eom-sigma} instead gives the solution as a recursion:
\begin{equation}
  \label{p:eq:hs5:recursion}
  \begin{aligned}
    \Sigma_2^{(3)*}
      &
      = \Psi^2 + \frac{\gs}{5} \, \sprod{2}{2}^{-1} \circ \sprod{3}{2}(\Psi, \Sigma_2^{(3)*})
      \\ &
      = \quad
        \begin{tikzpicture}[diagram]
          \draw (0, 0.1) -- (2, 0.1) pic {ext} node[above right] {$\Psi$};
          \draw (0, -0.1) -- (2, -0.1) pic {ext} node[below right] {$\Psi$};
        \end{tikzpicture}
        \quad
        +
        \quad
        \begin{tikzpicture}[diagram]
          \draw[sigma line] (0, 0.1) -- (4, 0.1) pic {ext};
          \draw[sigma line] (0, -0.1) -- (4, -0.1) pic {ext};
          \node[right=3pt] at (4, 0) {$\Sigma_2^{(3)}$};
          \draw (2, 0.1) -- (2, 1.4) pic {ext} node[above=2pt] {$\Psi$};
          \node[below=3pt] at (0.9, -0.1) {$\sprod{2}{2}^{-1}$};
          \node[below=3pt] at (2, -0.1) {$\sprod{3}{2}$};
        \end{tikzpicture}
  \end{aligned}
\end{equation}
Iterating it, we get:
\begin{equation}
  \label{p:eq:hs5:saddle-series}
  \begin{aligned}
    \Sigma_2^{(3)*}
      &
      = \Psi^2
        + \frac{\gs}{5} \, \sprod{2}{2}^{-1} \circ \sprod{3}{2}(\Psi^3)
        + \frac{\gs^2}{25} \, \sprod{2}{2}^{-1} \circ \sprod{3}{2}\Big( \Psi, \sprod{2}{2}^{-1} \circ \sprod{3}{2}(\Psi^3) \Big)
        + O(\gs^3)
      \\ &
      = \quad
        \begin{tikzpicture}[diagram]
          \draw (0, 0.1) -- (2, 0.1) pic {ext} node[above right] {$\Psi$};
          \draw (0, -0.1) -- (2, -0.1) pic {ext} node[below right] {$\Psi$};
        \end{tikzpicture}
        \quad
        +
        \quad
        \begin{tikzpicture}[diagram]
          \draw[sigma line] (0, 0.1) -- (2, 0.1) pic {ext};
          \draw[sigma line] (0, -0.1) -- (2, -0.1) pic {ext};
          \draw (2, 0.1) -- (4, 0.1) pic {ext} node[above right] {$\Psi$};
          \draw (2, -0.1) -- (4, -0.1) pic {ext} node[below right] {$\Psi$};
          \draw (2, 0.1) -- (2, 1.4) pic {ext} node[above=2pt] {$\Psi$};
          \node[below=3pt] at (0.9, -0.1) {$\sprod{2}{2}^{-1}$};
          \node[below=3pt] at (2, -0.1) {$\sprod{3}{2}$};
        \end{tikzpicture}
        \quad
        +
        \quad
        \cdots
  \end{aligned}
\end{equation}
which is the Neumann series~\eqref{p:eq:hs5:neumann} applied to $\sprod{3}{2}(\Psi^3)$ (we did not include the numerical factors in the graphical representations).
Each iteration inserts one more $\sprod{3}{2}(\Psi, \cdot)$ between two $\sprod{2}{2}^{-1}$, with one more factor $\gs / 5$.
Read as Feynman diagrams, the term of order $\gs^n$ is a chain of $n$ quintic vertices joined by $\Sigma_2$ lines, each carrying one external $\Psi$.
The chain ends on the source $\Psi^2$ at one end and on the open $\Sigma_2$ slot at the other.
Up to the factor $12 / \gs^2$, the inverse of the dressed operator is the $\Sigma_2$ propagator in the background $\Psi$, and the chain is the geometric series of the insertions $\sprod{3}{2}(\Psi, \cdot)$ on the $\Sigma_2$ line.

The field $\Sigma_2$ enters the cubic action through its fluctuation from the solution:
\begin{equation}
  \label{p:eq:hs5:deviation}
  \sigma_2^{(3)}
    := \Sigma_2 - \Sigma_2^{(3)*},
\end{equation}
which replaces the fluctuation $\sigma_2^{(2)} = \Sigma_2 - \Psi^2$ of the quartic order.
The terms~\eqref{p:eq:hs5:F} can be written as a completed square in the fluctuation.
By~\eqref{p:eq:hs5:F-dressed} and the definition~\eqref{p:eq:hs5:saddle} of the solution, $\sprod{2}{2}(\Psi^2) = \kin{2}{\Psi}{3}(\Sigma_2^{(3)*})$, we get:
\begin{equation}
\label{p:eq:hs5:square}
\begin{aligned}
  \caction[4] + \caction[5]
    &
    = - \frac{\gs^2}{4!} \Big[ \pair[2]{\Sigma_2}{\kin{2}{\Psi}{3}(\Sigma_2)} - 2 \, \pair[2]{\Sigma_2}{\kin{2}{\Psi}{3}(\Sigma_2^{(3)*})} \Big]
    \\ &
    = - \frac{\gs^2}{4!} \Big[ \pair[2]{\sigma_2^{(3)}}{\kin{2}{\Psi}{3}(\sigma_2^{(3)})} - \pair[2]{\Sigma_2^{(3)*}}{\kin{2}{\Psi}{3}(\Sigma_2^{(3)*})} \Big],
\end{aligned}
\end{equation}
using the symmetry of the dressed operator.
The first term is the quadratic form of the fluctuation and vanishes on-shell.
The second term does not depend on $\Sigma_2$ and is the value of the action on the solution.


\subparagraph{Induced interactions.}

Expanding the second form of~\eqref{p:eq:hs5:saddle} in the constant term of~\eqref{p:eq:hs5:square}, and using $\kin{2}{\Psi}{3} \circ (\kin{2}{\Psi}{3})^{-1} = 1$ together with $\pair[2]{A}{\sprod{3}{2}(\Psi^3)} = \vertex{5}(A, \Psi^2, \Psi)$, we get:
\begin{equation}
  \label{p:eq:hs5:constant}
  \pair[2]{\Sigma_2^{(3)*}}{\kin{2}{\Psi}{3}(\Sigma_2^{(3)*})}
    = \vertex{4}(\Psi^4)
      + \frac{\gs}{5} \, \vertex{5}(\Psi^5)
      + \frac{\gs^2}{25} \, \vertex{5}\Big( \Psi^3, \big( \kin{2}{\Psi}{3} \big)^{-1} \circ \sprod{3}{2}(\Psi^3) \Big).
\end{equation}
The first two terms come from $\pair[2]{\Psi^2}{\kin{2}{\Psi}{3}(\Psi^2)} = \vertex{4}(\Psi^4) - \tfrac{\gs}{5} \vertex{5}(\Psi^5)$ and from twice the cross term $\tfrac{\gs}{5} \vertex{5}(\Psi^5)$.
Multiplied by $\gs^2 / 4!$, the first two terms are the quartic and quintic terms of $S_{\le 5}$.
This is how the coefficient $1 / 5!$ of~\eqref{p:eq:hs5:S5tilde} enters: the quintic term of the target is the cross term minus the quintic part of the dressed operator, and the two combine to $+ \gs / 5$ only with this coefficient.
Putting~\eqref{p:eq:hs5:square} and~\eqref{p:eq:hs5:constant} together, and adding back the terms of $\Psi$ alone, we get:
\begin{equation}
  \label{p:eq:hs5:identity}
  \caction[\le 5]
    = S_{\le 5}
      - \frac{\gs^2}{4!} \, \pair[2]{\sigma_2^{(3)}}{\kin{2}{\Psi}{3}(\sigma_2^{(3)})}
      + \frac{\gs^4}{600} \, \vertex{5}\Big( \Psi^2, \Psi, \big( \kin{2}{\Psi}{3} \big)^{-1} \circ \sprod{3}{2}(\Psi^3) \Big),
\end{equation}
This is an identity in the fields $(\Psi, \Sigma_2)$, exact in $\gs$, and the counterpart at this order of the completed square~\eqref{p:eq:hs4:action-tilde}.
The three terms are the target, a quadratic form in the fluctuation which vanishes on-shell, and a term independent of $\Sigma_2$.
The last term starts at order $\gs^4$, where it is of degree six in $\Psi$.
It is a six-point interaction induced by the exchange of $\Sigma_2$, which the Gaussian elimination produces in addition to the vertex for which the auxiliary field was introduced.

On-shell, $\sigma_2^{(3)} = 0$ and the elimination of $\Sigma_2$ gives:
\begin{equation}
  \label{p:eq:hs5:eliminate}
  \caction[\le 5] \big|_{\Sigma_2 = \Sigma_2^{(3)*}}
    = S_{\le 5}
      + \frac{\gs^4}{600} \, \vertex{5 \cdot 5}(\Psi^6)
      + O(\gs^5),
\end{equation}
where we expanded the dressed inverse with~\eqref{p:eq:hs5:neumann} and introduced the composite vertex
\begin{equation}
  \label{p:eq:hs5:exchange-vertex}
  \begin{aligned}
    \vertex{5 \cdot 5}(A_1, \ldots, A_6)
      &
      := \vertex{5}\Big( A_1, A_2, A_3, \sprod{2}{2}^{-1} \circ \sprod{3}{2}(A_4, A_5, A_6) \Big)
      \\ &
      = \pair[2]{\sprod{3}{2}(A_1, A_2, A_3)}{\sprod{2}{2}^{-1} \circ \sprod{3}{2}(A_4, A_5, A_6)},
  \end{aligned}
\end{equation}
a composite six-point vertex obtained from two quintic vertices joined by a $\Sigma_2$ line.
It is a six-linear form, graded-symmetric within each triple of arguments and under the exchange of the two triples, since $\sprod{2}{2}^{-1}$ is symmetric with respect to the $2$-string pairing as $\sprod{2}{2}$ is.
We read it also as a trilinear form on $2$-string states: $\vertex{5 \cdot 5}(B_1, B_2, B_3)$ with $B_i \in \Hn{2}$ places $B_1$ in the first two slots, $B_3$ in the last two, and the two strands of $B_2$ in the third and sixth slots.
The result is averaged over the three choices of the split argument, so that the form is symmetric in $(B_1, B_2, B_3)$.
Off the diagonal, this symmetrized form is the one meant whenever $\vertex{5 \cdot 5}$ is evaluated on $2$-string states, in particular in the reduced sextic vertex of \cref{p:sec:hs:sextic}.
On the diagonal this reading gives $\vertex{5 \cdot 5}(\Psi^2, \Psi^2, \Psi^2) = \vertex{5 \cdot 5}(\Psi^6)$.


Once the quintic vertex is included, the auxiliary field thus induces interactions at all orders.
A $\Sigma_2$ line collapses to a contact term when one of its ends is a source vertex, as in~\eqref{p:eq:hs4:telescoping}.
The line of $\vertex{5 \cdot 5}$ joins two quintic vertices and in general does not collapse.
The sextic order must therefore reproduce $\vertex{6}$ and cancel this induced term at the same time, which is the subject of \cref{p:sec:hs:sextic}.
\Cref{p:sec:worldsheet} gives the interpretation of this exchange in terms of moduli space.

This also means that this approach of using $\Sigma_2$ directly in the quintic term works only because SFT contains an infinite number of higher-order terms.
Indeed, any term induced at higher orders can be re-packaged with the higher-order elementary vertices.
In normal QFT, we would not be able to use the same technique, without also introducing an infinite tower of auxiliary fields.
Another way to deal with odd powers was suggested in~\cite{Rivasseau:2010:LoopVertexExpansion,Lionni:2019:IntermediateFieldRepresentation}.


\subparagraph{The five-point diagrams.}

The identity~\eqref{p:eq:hs5:identity} guarantees that every tree amplitude of the cubic action agrees with that of $S_{\le 5}$ through order $\gs^3$, under the assumptions on the exchange stated in \cref{p:sec:hs:quartic} and the standing assumptions of this subsection.
We show how the new diagram of the five-point amplitude reduces to the quintic vertex, since this is the first order at which the auxiliary field gives a new diagram.
Tree amplitudes with $n$ external legs scale as $\gs^{n - 2}$, so order $\gs^3$ is the five-point tree level.
The Feynman rules are those of \cref{p:sec:hs:quartic}, in the same convention: the propagator $\tfrac{12}{\gs^2} \sprod{2}{2}^{-1}$ and the source vertex $\tfrac{\gs^2}{3!} \vertex{4}(\Sigma_2, \Psi^2)$ of~\eqref{p:eq:hs4:feynman-rules}.
The quintic coupling~\eqref{p:eq:hs5:S5tilde} adds a vertex with two $\Sigma_2$ legs and one $\Psi$ leg:
\begin{equation}
  \label{p:eq:hs5:feynman-rule}
  \frac{2 \gs^3}{5!} \, \vertex{5}(\Sigma_2, \Sigma_2, \Psi).
\end{equation}
A $\Sigma_2$ line ends either on a source vertex, linear in $\Sigma_2$, or on the quintic vertex, quadratic in it.
\begin{figure}[htp]
  \centering
  \begin{tikzpicture}[baseline={(0, -0.1)}]
    \begin{feynman}
      \vertex[dot] (v) at (0, 0) {};
      \vertex (a) at (-1.2, 0.8) {$\Sigma_2$};
      \vertex (b) at (-1.2, -0.8) {$\Sigma_2$};
      \vertex (c) at (1.2, 0) {$\Psi$};
      \diagram* {
      (a) -- [double, double distance=1.5pt] (v) -- [double, double distance=1.5pt] (b),
      (v) -- [plain] (c),
      };
    \end{feynman}
  \end{tikzpicture}
  \qquad \qquad
  \begin{tikzpicture}[baseline={(0, -0.1)}]
    \begin{feynman}
      \vertex[dot] (s1) at (0, 0.9) {};
      \vertex[dot] (s2) at (0, -0.9) {};
      \vertex[dot] (q) at (1.6, 0) {};
      \vertex (p1) at (-1, 1.5) {$\Psi$};
      \vertex (p2) at (-1, 0.4) {$\Psi$};
      \vertex (p3) at (-1, -0.4) {$\Psi$};
      \vertex (p4) at (-1, -1.5) {$\Psi$};
      \vertex (p5) at (2.7, 0) {$\Psi$};
      \diagram* {
      (p1) -- [plain] (s1),
      (p2) -- [plain] (s1),
      (p3) -- [plain] (s2),
      (p4) -- [plain] (s2),
      (s1) -- [double, double distance=1.5pt, edge label={$\Sigma_2$}] (q),
      (s2) -- [double, double distance=1.5pt, edge label'={$\Sigma_2$}] (q),
      (q) -- [plain] (p5),
      };
    \end{feynman}
  \end{tikzpicture}
  \quad $\longrightarrow$ \quad
  \begin{tikzpicture}[baseline={(0, -0.1)}]
    \begin{feynman}
      \vertex[dot] (c) at (0, 0) {};
      \vertex (l1) at (90:1.1) {$\Psi$};
      \vertex (l2) at (162:1.1) {$\Psi$};
      \vertex (l3) at (234:1.1) {$\Psi$};
      \vertex (l4) at (306:1.1) {$\Psi$};
      \vertex (l5) at (18:1.1) {$\Psi$};
      \diagram* {
      (c) -- [plain] (l1),
      (c) -- [plain] (l2),
      (c) -- [plain] (l3),
      (c) -- [plain] (l4),
      (c) -- [plain] (l5),
      };
    \end{feynman}
  \end{tikzpicture}
  \caption{Left, the vertex~\eqref{p:eq:hs5:feynman-rule} of the quintic coupling, with two $\Sigma_2$ legs (double lines) and one $\Psi$ leg (single line); right, the five-point exchange diagram, two source vertices joined by $\Sigma_2$ lines to the quintic vertex, and its collapse to the contact vertex $\vertex{5}$.}
  \label{p:fig:hs5:five-point}
\end{figure}

The new diagram joins two source vertices by $\Sigma_2$ lines to the quintic vertex, whose $\Psi$ leg carries the fifth external state (\cref{p:fig:hs5:five-point}).
Both lines collapse at once, since the propagated source is the pair state: the exchange diagram evaluates $\vertex{5}(\Psi^2, \Psi^2, \Psi)$.
The collapse holds up to an element of $\ker \sprod{2}{2}$, on which the quintic vertex vanishes by~\eqref{p:eq:hs5:kernel-condition}.
A quintic vertex joined to a cubic vertex by its $\Psi$ leg, or to a second quintic vertex by a $\Sigma_2$ line, has six external legs and is of order $\gs^4$.
The second one is the diagram behind the induced term of~\eqref{p:eq:hs5:eliminate}.


\paragraph{Gauge transformation}

The transformation law~\eqref{p:eq:hs4:gauge-sigma} of $\Sigma_2$ at quartic order does not leave the action invariant at order $\gs^3$.

Since the solution~\eqref{p:eq:hs5:saddle} contains new terms which depend on $\Psi$, it transforms under~\eqref{p:eq:hs1:gauge-psi}, and the transformation of $\Sigma_2$ must be modified to compensate.
We generalize the pattern of the quartic order, where the transformation law~\eqref{p:eq:hs4:gauge-sigma} left the fluctuation invariant, and require the quadratic form $\pair[2]{\sigma_2^{(3)}}{\kin{2}{\Psi}{3}(\sigma_2^{(3)})}$ to be invariant.
We impose the invariance of this form because it is the only term of the identity~\eqref{p:eq:hs5:identity} which depends on $\Sigma_2$.
If it is invariant, the variation of the cubic action equals that of $S_{\le 5}$ up to the variation of the last term, which is of order $\gs^4$, the order of the truncation.
Its variation has two pieces, from the fluctuation and from the $\Psi$ inside the dressed operator.
By~\eqref{p:eq:hs5:dressed} and the symmetry of the dressed operator, we get:
\begin{equation}
  \label{p:eq:hs5:form-variation}
  \delta_\Lambda \pair[2]{\sigma_2^{(3)}}{\kin{2}{\Psi}{3}(\sigma_2^{(3)})}
    = 2 \, \pair[2]{\delta_\Lambda \sigma_2^{(3)}}{\kin{2}{\Psi}{3}(\sigma_2^{(3)})}
      - \frac{\gs}{5} \, \vertex{5}(\sigma_2^{(3)}, \sigma_2^{(3)}, \delta_\Lambda \Psi).
\end{equation}
This variation vanishes if the fluctuation transforms as:
\begin{equation}
  \label{p:eq:hs5:gauge-deviation}
  \delta_\Lambda \sigma_2^{(3)}
    = \frac{\gs}{10} \, \big( \kin{2}{\Psi}{3} \big)^{-1} \circ \sprod{3}{2}(\delta_\Lambda \Psi, \sigma_2^{(3)}).
\end{equation}
The fluctuation transforms homogeneously, by a $\Psi$-dependent linear rotation which vanishes on-shell.
This is the transformation law~\eqref{p:eq:gauge:law-quintic} of \cref{p:sec:gauge}, where its algebra is worked out.

Translating back to $\Sigma_2 = \Sigma_2^{(3)*} + \sigma_2^{(3)}$, the auxiliary field transforms as its own solution, plus the rotation of the fluctuation:
\begin{equation}
  \label{p:eq:hs5:gauge-sigma}
  \delta_\Lambda \Sigma_2
    = \delta_\Lambda \Sigma_2^{(3)*}
      + \frac{\gs}{10} \, \big( \kin{2}{\Psi}{3} \big)^{-1} \circ \sprod{3}{2}(\delta_\Lambda \Psi, \sigma_2^{(3)}),
\end{equation}
where $\delta_\Lambda \Sigma_2^{(3)*}$ is the variation of~\eqref{p:eq:hs5:saddle} induced by $\delta_\Lambda \Psi$.
The transformation law of the quartic order is the order-$\gs^0$ truncation of $\delta_\Lambda \Sigma_2 = \delta_\Lambda \Sigma_2^{(3)*}$.
Through order $\gs$, the transformation law reads:
\begin{equation}
  \label{p:eq:hs5:gauge-sigma-g1}
  \begin{aligned}
    \delta_\Lambda \Sigma_2
      = 2 \, \delta_\Lambda \Psi \otimes \Psi
        &
        + \frac{3 \gs}{5} \, \sprod{2}{2}^{-1} \circ \sprod{3}{2}(\delta_\Lambda \Psi, \Psi^2)
        \\ &
        + \frac{\gs}{10} \, \sprod{2}{2}^{-1} \circ \sprod{3}{2}(\delta_\Lambda \Psi, \Sigma_2 - \Psi^2)
        + O(\gs^2),
    \end{aligned}
\end{equation}
the middle term being the variation of $\tfrac{\gs}{5} \sprod{2}{2}^{-1} \circ \sprod{3}{2}(\Psi^3)$, cubic in $\Psi$.
The last term vanishes on-shell but cannot be dropped off-shell: it compensates the variation of the dressed operator.

With~\eqref{p:eq:hs5:gauge-sigma} the quadratic form is exactly invariant.
Varying the identity~\eqref{p:eq:hs5:identity} gives:
\begin{equation}
  \label{p:eq:hs5:delta-S-equal}
  \delta_\Lambda \caction[\le 5] - \delta_\Lambda S_{\le 5}
    = \frac{\gs^4}{600} \, \delta_\Lambda \vertex{5}\Big( \Psi^2, \Psi, \big( \kin{2}{\Psi}{3} \big)^{-1} \circ \sprod{3}{2}(\Psi^3) \Big)
    = O(\gs^4),
\end{equation}
identically, off-shell, for any variation $\delta_\Lambda \Psi$ and exactly in $\gs$.
Since $\delta_\Lambda S_{\le 5} = O(\gs^4)$, the cubic action is invariant up to terms of order $\gs^4$, the order of the truncation it should match.
We discuss the nature of the rotation term, the commutator of two gauge transformations and the Batalin--Vilkovisky structure of the cubic action in \cref{p:sec:gauge}.


\paragraph{Kernel condition and reality}

\subparagraph{Kernel condition.}

The auxiliary field lives on the quotient~\eqref{p:eq:hs4:quotient}, and the quintic coupling must be a function on it.
The two quartic terms of~\eqref{p:eq:hs5:F} are invariant under $\Sigma_2 \to \Sigma_2 + N_2$ for $N_2 \in \ker \sprod{2}{2}$, by~\eqref{p:eq:hs4:kernel}, while the quintic term shifts by
\begin{equation}
  \vertex{5}((\Sigma_2 + N_2)^2, \Psi) - \vertex{5}(\Sigma_2^2, \Psi)
    \to 2 \, \vertex{5}(N_2, \Sigma_2, \Psi) + \vertex{5}(N_2^2, \Psi).
\end{equation}
For this to vanish for all $\Sigma_2$, the term linear in $\Sigma_2$ must vanish, which then kills the quadratic one.
The cubic action therefore descends to the quotient if and only if the quintic vertex vanishes on the kernel:
\begin{equation}
  \label{p:eq:hs5:kernel-condition}
  \forall N_2 \in \ker \sprod{2}{2}, \;
  \forall A \in \Hn{2}, \;
  \forall B \in \Hs:
  \quad
  \vertex{5}(N_2, A, B)
    = 0,
\end{equation}
that is, $\sprod{3}{2}(A, B)$ pairs to zero with every kernel element.
In infinite dimension, the states which pair to zero with the kernel of $\sprod{2}{2}$ form only the closure of its image: the condition thus places the values of $\sprod{3}{2}$ in the closure of the image of $\sprod{2}{2}$, and the first assumption below, that they lie in the image itself, is what makes $\sprod{2}{2}^{-1}$ defined on them.
It also says that $\sprod{3}{2}(\Psi, \cdot)$ does not depend on the representative of its argument.
The condition thus makes $\sprod{3}{2}(\Psi, \cdot)$ descend to the quotient, and with it the dressed operator~\eqref{p:eq:hs5:dressed-factor} and its inverse~\eqref{p:eq:hs5:neumann}.
Since $\vertex{5}$ is symmetric, the condition is equivalent by polarization to its restriction to the diagonal, $\vertex{5}(N_2, \Psi^2, \Psi) = 0$ for all $\Psi$, which is the only evaluation of the quintic vertex in the matching at order $\gs^3$.

We have no proof that the quintic vertex vanishes on a null vector of $\sprod{2}{2}$, and nothing imposes it.
  At the same time, no null vector of $\sprod{2}{2}$ is known above level $0$ (at level $0$ there is none at any real momentum for the sections with $\beta > \sqrt{2}$, \cref{p:sec:hs:quartic}), and none is forced by the mechanisms we have checked at the levels and momenta where the couplings are known: the dilaton theorem with the ghost-number counting behind the zeros of the cubic vertex, BRST-exact strands, the universality of the zero-momentum couplings, and the worldsheet twist.
  If $\sprod{2}{2}$ has an accidental zero at some real momentum, with a single null vector $N_2$, the condition at that momentum is that $N_2$ pairs to zero with the quintic source $\sprod{3}{2}(\Psi^2, \Psi)$, one number for each $\Psi$, a property of $\vertex{5}$ at that momentum which nothing enforces and which we expect to fail in general; the matching is then restored by the multiplier described next.

When it fails, the quotient is not available and one must choose a complement of the kernel.
The eliminated action then differs from $S_{\le 5}$ at order $\gs^3$ by a five-point term which depends on the complement and cannot vanish for all of them.
The kernel directions can then be handled by a multiplier rather than by inversion, which restores the matching at this order and keeps the action cubic; this is taken up in \cref{p:sec:all-orders}.

Hence, unless said otherwise, we make two assumptions from this order on: 1) the states on which $\sprod{2}{2}^{-1}$ acts, here the values of $\sprod{3}{2}$, lie in the image of $\sprod{2}{2}$ and not only in its closure, which is all the kernel condition gives, so that the inverse exists on them level by level (\cref{p:sec:hs:quartic}); 2) the quintic vertex vanishes on the kernel of the quartic vertex,~\eqref{p:eq:hs5:kernel-condition} (``kernel condition'').
The sextic vertex is subject to an analogous condition at the next order.


\subparagraph{Reality.}

The quintic order is the first where the solution and the gauge law carry an inverse, so that reality is trickier to prove.

We use the reality of the vertices~\eqref{p:eq:hs4:vertex-real} for $n = 4$ and $n = 5$, and the condition~\eqref{p:eq:hs5:kernel-condition}.
  The kernel of $\sprod{2}{2}$ is stable under the star conjugation.
  Indeed, if $\vertex{4}(A, N_2) = 0$ for all $A$, then $\vertex{4}(A, N_2^\star) = \vertex{4}(A^\star, N_2)^* = 0$ for all $A$, since $A^\star$ ranges over all $2$-string states as $A$ does.
  Hence $\star$ descends to the quotient, and the condition~\eqref{p:eq:hs4:reality}, $\Sigma_2^\star = \Sigma_2$, makes sense there.

The two maps from which the solution is built, $\sprod{2}{2}^{-1} \circ \sprod{3}{2}$ and $(\kin{2}{\Psi}{3})^{-1} \circ \sprod{2}{2}$, commute with $\star$ for real $\Psi$.
This follows by conjugating both sides of the two relations which characterize them:
\begin{equation}
\vertex{4}\big( A, \sprod{2}{2}^{-1} \circ \sprod{3}{2}(C) \big)
  = \vertex{5}(A, C),
\qquad
\pair[2]{A}{\kin{2}{\Psi}{3}(B')}
  = \vertex{4}(A, B),
\end{equation}
with $B' := (\kin{2}{\Psi}{3})^{-1} \circ \sprod{2}{2}(B)$.
The conjugation is applied only to states whose strands have the ghost number of the string field, in $\Hs$, $\Hn{2}$ and $\Hn{3}$, where it acts strand by strand.
It is never applied to the outputs of $\sprod{2}{2}$ and $\sprod{3}{2}$, whose strands have ghost number three.
All coefficients are real rationals times $\gs$.

With $\Psi^\star = - \Psi$, the state $\Psi \otimes \Psi$ is star-even.
For $\gs$ real, the solution~\eqref{p:eq:hs5:saddle} therefore satisfies $\Sigma_2^\star = \Sigma_2$, to all orders of the formal expansion in $\gs$ and, within the quintic truncation, wherever the dressed operator is invertible on the quotient.
For a real variation $\delta_\Lambda \Psi$, each term of the transformation law~\eqref{p:eq:hs5:gauge-sigma} is real for the same reason, and the transformation law maps the real locus of $(\Psi, \Sigma_2)$ to itself.
A gauge parameter satisfying the reality condition $\Lambda^\star = - \Lambda$ of \cref{p:app:conventions} produces such a real variation, the vertices being real.


At the sextic order the feedback of the quintic order must be cancelled together with the production of the sextic vertex, which leads to the reduced vertex and to the $3$-string field $\Sigma_3$.
The paper's block decomposes the reduced vertex, eliminates the auxiliary fields and states the consistency conditions on the kernels; the hybrid form of the kernel condition, which in the paper follows here, is in \cref{ch:multiplier-hybrid}.
The block states that, within $(\Psi, \Sigma_2)$ and for generic vertices, a sextic coupling of three $\Sigma_2$ alone leaves the feedback uncancelled: \status{Conditional(PH1, PH5, PH14 [$O(g_s^4)$]) -- Proven}\record{013-b3-sextic-milestone (R19)}.

\subsubsection{Sextic order, \texorpdfstring{$O(\gs^4)$}{O(gs**4)}}
\label{p:sec:hs:sextic}

The sextic order is the first where the auxiliary field of the previous order induces a composite vertex to be decomposed.
The vertex to add is therefore $\vertex{6}$ minus the exchange term of \cref{p:sec:hs:quintic}.
In its simplest implementation, the $3$-string field $\Sigma_3$ introduced at this order has a kinetic operator which is not a vertex of closed string field theory.


\paragraph{Decomposition of the reduced vertex}
\label{p:sec:hs:sextic:vertex}

The target at this order is:
\begin{equation}
  S_{\le 6}
    = S_{\le 5} + \frac{\gs^4}{6!} \vertex{6}(\Psi^6).
\end{equation}
The cubic action $\caction[\le 5]$ of the quintic order induces a new term beyond $S_{\le 5}$: by~\eqref{p:eq:hs5:eliminate}, eliminating $\Sigma_2$ from it produces the six-point term $\tfrac{\gs^4}{600} \vertex{5 \cdot 5}(\Psi^6)$, the exchange of $\Sigma_2$ between two quintic vertices.
The sextic order must thus reproduce on-shell the difference:
\begin{equation}
  \label{p:eq:hs6:remainder}
  S_{\le 6} - \caction[\le 5] \big|_{\Sigma_2 = \Sigma_2^{(3)*}}
    = \frac{\gs^4}{6!} \Big[ \vertex{6}(\Psi^6) - \frac{6}{5} \, \vertex{5 \cdot 5}(\Psi^6) \Big]
      + O(\gs^5).
\end{equation}
Two things have to happen at once: the sextic vertex must be produced, and the induced exchange term must be cancelled.
The second term is not a vertex of closed string field theory but a diagram built from $\vertex{4}$ and $\vertex{5}$.

We define the \emph{reduced vertex} $\dvertex{6}(A_1, \ldots, A_6)$ with general arguments as the fundamental vertex minus the composite vertex ~\eqref{p:eq:hs5:exchange-vertex}:
\begin{equation}
  \label{p:eq:hs6:dressed-vertex}
  \dvertex{6}(A_1, \ldots, A_6)
    := \vertex{6}(A_1, \ldots, A_6) - \frac{6}{5} \, \vertex{5 \cdot 5}(A_1, \ldots, A_6).
\end{equation}
This is represented graphically in \cref{p:fig:hs6:dressed-vertex}.
It is a six-linear form, graded-symmetric within each triple of arguments and under the exchange of the two triples, which is the symmetry of its exchange part.
It is not symmetric under the permutations mixing the two triples.
Read as a form in three $2$-string arguments, $\dvertex{6}(B_1, B_2, B_3)$ with $B_i \in \Hn{2}$, it is symmetrized in them, as $\vertex{5 \cdot 5}$ is in \cref{p:sec:hs:quintic}.
On the diagonal, where every argument is built from $\Psi$ alone, all readings agree: $\dvertex{6}(\Psi^3, \Psi^3) = \dvertex{6}(\Psi^2, \Psi^2, \Psi^2) = \dvertex{6}(\Psi^6)$.
By~\eqref{p:eq:hs6:remainder}, the term to add on-shell is thus the sextic term with $\dvertex{6}(\Psi^6)$ in place of $\vertex{6}(\Psi^6)$.
\begin{figure}[htp]
  \centering
  \begin{tikzpicture}
    \begin{feynman}
      \vertex (w) at (0, 0) [dot, label=above:{$\mathcal{W}_6$}];
      \vertex (w1) at (-0.9, 0.8) {};
      \vertex (w2) at (-1.1, 0) {};
      \vertex (w3) at (-0.9, -0.8) {};
      \vertex (w4) at (0.9, 0.8) {};
      \vertex (w5) at (1.1, 0) {};
      \vertex (w6) at (0.9, -0.8) {};
      \vertex (eq) at (1.9, 0) {$=$};
      \vertex (v) at (3.3, 0) [dot, label=above:{$\mathcal{V}_6$}];
      \vertex (v1) at (2.4, 0.8) {};
      \vertex (v2) at (2.2, 0) {};
      \vertex (v3) at (2.4, -0.8) {};
      \vertex (v4) at (4.2, 0.8) {};
      \vertex (v5) at (4.4, 0) {};
      \vertex (v6) at (4.2, -0.8) {};
      \vertex (minus) at (5.3, 0) {$- \; \dfrac{6}{5}$};
      \vertex (a) at (7.0, 0) [dot];
      \vertex (b) at (9.0, 0) [dot];
      \vertex (a1) at (6.2, 0.8) {};
      \vertex (a2) at (6.0, 0) {};
      \vertex (a3) at (6.2, -0.8) {};
      \vertex (b1) at (9.8, 0.8) {};
      \vertex (b2) at (10.0, 0) {};
      \vertex (b3) at (9.8, -0.8) {};
      \diagram* {
      (w) -- (w1), (w) -- (w2), (w) -- (w3), (w) -- (w4), (w) -- (w5), (w) -- (w6),
      (v) -- (v1), (v) -- (v2), (v) -- (v3), (v) -- (v4), (v) -- (v5), (v) -- (v6),
      (a) -- (a1), (a) -- (a2), (a) -- (a3), (b) -- (b1), (b) -- (b2), (b) -- (b3),
      (a) -- [plain, double, edge label={$\Sigma_2$}] (b),
      };
    \end{feynman}
  \end{tikzpicture}
  \caption{The reduced sextic vertex~\eqref{p:eq:hs6:dressed-vertex}: the sextic vertex minus the composite vertex, two quintic vertices joined by a $\Sigma_2$ line.}
  \label{p:fig:hs6:dressed-vertex}
\end{figure}


We have several options to decompose this vertex, since the matching constrains the sextic term only through its value on the diagonal, which is a sextic polynomial in $\Psi$.
Any trilinear form on $\Hn{2}$ with the value $\dvertex{6}(\Psi^6)$ on $\Sigma_2 = \Psi^2$ reproduces the target at this order.
Two forms differing by one that vanishes on the diagonal give the same eliminated action: they are equivalent at this truncation.
The construction we keep is the one that stays exact at all orders.
We see that we cannot immediately apply the trick of the quartic order with a $3$-string field $\Sigma_3 \in \Hn{3}$, because we would get a quartic term $\mc V_6(\Sigma_3, \Psi^3)$.
Another solution could be to write $\mc V_6(\Sigma_2^3)$, but this would not generalize at higher orders since $\mc V_8(\Sigma_2^4)$ is quartic.\footnotemark{}
\footnotetext{
  Such a term appeared in~\cite{Erbin:2024:OpenStringStub}.
  This term together with the terms linear and quadratic in the auxiliary field were sufficient to reproduce the complete action, with a single auxiliary field.
  This is specific to the structure of the interactions induced by the stubs.
}
The way out is to take an intermediary path and replace $\Psi^2$ by $\Sigma_2$ in the tensor power, since the solution $\Sigma_2 = \Psi^2 + O(\gs)$ guarantees that only $\dvertex{6}(\Psi^6)$ is produced at this order.
Moreover, this form is suggested by the form of the composite vertex, whose arguments are identical, $\vertex{5 \cdot 5}(\Sigma_2^2, \Psi^2)$.


We write the sextic term as $\dvertex{6}(\Psi \otimes \Sigma_2, \Psi \otimes \Sigma_2)$, with the $3$-string state $\Psi \otimes \Sigma_2$ in each triple of slots, and add the square of the fluctuation of $\Sigma_3$ from it:
\begin{equation}
  \label{p:eq:hs6:DeltaS6}
  \Delta S_6
    := - \frac{\gs^4}{6!} \, \dvertex{6}(\Sigma_3 - \Psi \otimes \Sigma_2, \Sigma_3 - \Psi \otimes \Sigma_2),
\end{equation}
where
\begin{equation}
  \Sigma_3 \in \Hn{3},
  \qquad
  \Ngh(\Sigma_3)
    = 6.
\end{equation}
By construction, $\Delta S_6$ vanishes when $\Sigma_3 = \Psi \otimes \Sigma_2$.
Expanding the square, we get the sextic term of the cubic action:
\begin{equation}
  \label{p:eq:hs6:S6tilde}
  \begin{aligned}
    \caction[6]
      &
      := \frac{\gs^4}{6!} \, \dvertex{6}(\Psi \otimes \Sigma_2, \Psi \otimes \Sigma_2) + \Delta S_6
      \\ &
      = - \frac{\gs^4}{6!} \, \dvertex{6}(\Sigma_3^2)
        + \frac{2 \gs^4}{6!} \, \dvertex{6}(\Psi \otimes \Sigma_2, \Sigma_3),
  \end{aligned}
\end{equation}
a quadratic term for $\Sigma_3$ and a cubic coupling of $\Psi$, $\Sigma_2$ and $\Sigma_3$.
As a consequence, the action through this order,
\begin{equation}
  \label{p:eq:hs6:action-tilde}
  \caction[\le 6]
    := \caction[\le 5] + \caction[6],
\end{equation}
is cubic in $(\Psi, \Sigma_2, \Sigma_3)$.
In terms of the elementary vertices of closed string field theory, we get:
\begin{equation}
  \label{p:eq:hs6:S6tilde-explicit}
  \begin{aligned}
    \caction[6]
      =
        &
        - \frac{\gs^4}{6!} \, \vertex{6}(\Sigma_3^2)
        + \frac{\gs^4}{600} \, \vertex{5}\Big( \Sigma_3, \sprod{2}{2}^{-1} \circ \sprod{3}{2}(\Sigma_3) \Big)
      \\ &
      + \frac{2 \gs^4}{6!} \, \vertex{6}(\Psi, \Sigma_2, \Sigma_3)
        - \frac{2 \gs^4}{600} \, \vertex{5}\Big( \Sigma_3, \sprod{2}{2}^{-1} \circ \sprod{3}{2}(\Psi, \Sigma_2) \Big),
  \end{aligned}
\end{equation}
The second and fourth terms, the kinetic term of $\Sigma_3$ and its coupling to $(\Psi, \Sigma_2)$, are both non-local in the sense of effective field theory, since they contain the inverse of a vertex, $\sprod{2}{2}^{-1}$.
This is the price of cancelling an exchange diagram with a contact term, and a typical phenomenon in effective field theory.
For generic vertices, in which the exchange term is not a linear combination of the elementary vertices, the elementary vertices alone cannot cancel the exchange term, which contains $\sprod{2}{2}^{-1}$: the subtraction must itself be an exchange term, produced either by the reduced vertex or by a second $2$-string field.
An alternative is to enlarge the field content with a second $2$-string field.
  We show it at the sextic order in \cref{p:app:wrong-sign} and will not develop this approach further.

Since the sextic vertex is now read with three of its slots as outputs, we define the product $\sprod{3}{3} \colon \Hn{3} \to \Hn{3}$ by:
\begin{equation}
  \label{p:eq:hs6:p33}
  \vertex{6}(A_1, \ldots, A_6)
    =: \pair[3]{A_1 \otimes A_2 \otimes A_3}{\sprod{3}{3}(A_4, A_5, A_6)},
\end{equation}
with the $3$-string pairing on the left, as $\sprod{2}{2}$ and $\sprod{3}{2}$ were defined from $\vertex{4}$ and $\vertex{5}$.
It has ghost number three and it is symmetric with respect to the $3$-string pairing since $\vertex{6}$ is graded-symmetric.


\paragraph{Elimination of the auxiliary fields}
\label{p:sec:hs:sextic:eom}

The $3$-string field enters $\caction[\le 6]$ only through $\dvertex{6}$.
We write this form through the $3$-string pairing:
\begin{equation}
  \label{p:eq:hs6:kinetic}
  \begin{aligned}
    \dvertex{6}(A, B)
      &
      =: \pair[3]{A}{\wprod{3}{3}(B)},
    \\
    \wprod{3}{3}
      &
      := \sprod{3}{3} - \frac{6}{5} \, \sprod{2}{3} \circ \sprod{2}{2}^{-1} \circ \sprod{3}{2}
    \qquad
    \wprod{3}{3} \colon \Hn{3} \to \Hn{3},
  \end{aligned}
\end{equation}
for $A, B \in \Hn{3}$, where $\sprod{2}{3} \colon \Hn{2} \to \Hn{3}$ is the transpose of $\sprod{3}{2}$,~\eqref{p:eq:conv:transpose}.
This is the kinetic operator of $\Sigma_3$.
It is symmetric with respect to the $3$-string pairing, as $\sprod{3}{3}$ and $\sprod{2}{2}^{-1}$ are, and it does not depend on the fields.
It receives no correction at higher orders.

Varying~\eqref{p:eq:hs6:S6tilde} with respect to the three fields, we get:
\begin{equation}
  \label{p:eq:hs6:delta-S}
  \begin{aligned}
    \delta \caction[6]
      &
      = \frac{2 \gs^4}{6!} \, \dvertex{6}(\delta \Psi \otimes \Sigma_2, \Sigma_3)
        + \frac{2 \gs^4}{6!} \, \dvertex{6}(\Psi \otimes \delta \Sigma_2, \Sigma_3)
      \\ &
      \quad
      - \frac{2 \gs^4}{6!} \, \pair[3]{\delta \Sigma_3}{\wprod{3}{3}(\Sigma_3)}
        + \frac{2 \gs^4}{6!} \, \pair[3]{\delta \Sigma_3}{\wprod{3}{3}(\Psi \otimes \Sigma_2)}.
  \end{aligned}
\end{equation}
The equations of motion of $\caction[\le 6]$ are:
\begin{subequations}
  \label{p:eq:hs6:eom}
  \begin{align}
    \label{p:eq:hs6:eom-psi}
    \mathcal E^{(4)}_\Psi
      &
      := \mathcal E^{(3)}_\Psi
        + \frac{2 \gs^4}{6!} \, \dvertex{6}(\, \cdot \, \otimes \Sigma_2, \Sigma_3)
      = 0,
    \\
    \label{p:eq:hs6:eom-sigma2}
    \mathcal E^{(4)}_{\Sigma_2}
      &
      := \mathcal E^{(3)}_{\Sigma_2}
        + \frac{\gs^2}{30} \, \dvertex{6}(\Psi \otimes \cdot \,, \Sigma_3)
      = 0,
    \\
    \label{p:eq:hs6:eom-sigma3}
    \mathcal E^{(4)}_{\Sigma_3}
      &
      := \wprod{3}{3}(\Psi \otimes \Sigma_2) - \wprod{3}{3}(\Sigma_3)
      = 0,
  \end{align}
\end{subequations}
where $\mathcal E^{(3)}_\Psi$ and $\mathcal E^{(3)}_{\Sigma_2}$ are the equations~\eqref{p:eq:hs5:eom} of the quintic order, and where $\dvertex{6}(\, \cdot \, \otimes \Sigma_2, \Sigma_3)$ denotes the state paired with the open slot.

The equation for $\Sigma_3$ is algebraic and linear, with the source $\wprod{3}{3}(\Psi \otimes \Sigma_2)$ and the field-independent kinetic operator $\wprod{3}{3}$.
A kernel direction $N_3 \in \ker \wprod{3}{3}$ appears neither in the quadratic term nor in the coupling of~\eqref{p:eq:hs6:S6tilde}, by the symmetry of $\wprod{3}{3}$.
It has zero action and couples to nothing at this order, as the kernel of $\sprod{2}{2}$ at the quartic order.
We therefore take the $3$-string field on the quotient:
\begin{equation}
  \label{p:eq:hs6:quotient}
  \Sigma_3
    \in \Hn{3} / \ker \wprod{3}{3},
\end{equation}
where $\wprod{3}{3}$ is injective by construction.
The equation of motion~\eqref{p:eq:hs6:eom-sigma3} then has the unique solution:
\begin{equation}
  \label{p:eq:hs6:saddle-sigma3}
  \Sigma_3^{(4)*}
    = \Psi \otimes \Sigma_2,
\end{equation}
for any $\Sigma_2$, and $\Psi \otimes \Sigma_2^{(4)*}$ on the joint solution.
The $3$-string field satisfies the reality condition $\Sigma_3^\star = - \Sigma_3$, the sign of three star-odd strands, and the kernel of $\wprod{3}{3}$ is stable under $\star$, so that the condition is meaningful on the quotient~\eqref{p:eq:hs6:quotient}.
The solution is real for real $\Psi$ and $\Sigma_2$ (\cref{p:app:reality}).
The $3$-string field is thus an auxiliary field: it is fixed by $(\Psi, \Sigma_2)$ through an algebraic equation, and its solution lies on the diagonal, $\Sigma_3 = \Psi^3 + O(\gs)$ once $\Sigma_2$ is on-shell.
Since $\Delta S_6$ is the square of $\Sigma_3 - \Psi \otimes \Sigma_2$, it vanishes on the solution.
Eliminating $\Sigma_3$ thus gives back the sextic term we started from, exactly in $\gs$ and for any $\Sigma_2$, as in~\eqref{p:eq:hs4:eliminate}:
\begin{equation}
  \label{p:eq:hs6:eliminate-sigma3}
  \caction[\le 6] \big|_{\Sigma_3 = \Sigma_3^{(4)*}}
    = \caction[\le 5] + \frac{\gs^4}{6!} \, \dvertex{6}(\Psi \otimes \Sigma_2, \Psi \otimes \Sigma_2).
\end{equation}

The elimination of $\Sigma_2$ then proceeds as at the quintic order, with one new ingredient.
The sextic term is quadratic in $\Sigma_2$ and of order $\gs^4$, so it adds a correction of order $\gs^2$ to the dressed operator of $\Sigma_2$ and none to its source, which remains $\sprod{2}{2}(\Psi^2)$.
The solution therefore becomes $\Sigma_2^{(4)*} = \Sigma_2^{(3)*} + O(\gs^2)$.
By stationarity of $\caction[\le 5]$ at $\Sigma_2^{(3)*}$, as in \cref{p:sec:hs:quintic}, a shift of order $\gs^2$ changes $\caction[\le 5]$ at order $\gs^2 \times \gs^4$, and the sextic term, of order $\gs^4$, at the same order.
At order $\gs^4$ only the value of the sextic term on $\Sigma_2 = \Psi^2$ enters, $\dvertex{6}(\Psi^6)$.
With~\eqref{p:eq:hs5:eliminate} and~\eqref{p:eq:hs6:remainder}, we get:
\begin{equation}
  \label{p:eq:hs6:eliminate}
  \begin{aligned}
    \caction[\le 6] \big|_{\Sigma_3 = \Sigma_3^{(4)*}, \, \Sigma_2 = \Sigma_2^{(4)*}}
      &
      = S_{\le 5}
        + \frac{\gs^4}{600} \, \vertex{5 \cdot 5}(\Psi^6)
        + \frac{\gs^4}{6!} \Big[ \vertex{6}(\Psi^6) - \frac{6}{5} \, \vertex{5 \cdot 5}(\Psi^6) \Big]
        + O(\gs^5)
      \\ &
      = S_{\le 6} + O(\gs^5).
  \end{aligned}
\end{equation}
The exchange term produced by $\Sigma_2$ at the quintic order and the exchange part of the reduced vertex cancel exactly, and the sextic vertex is reproduced.
The elimination of $\Sigma_3$ is exact, and the elimination of $\Sigma_2$ matches the target through order $\gs^4$.
Its remainder of order $\gs^5$ is the next induced interaction, built from the reduced vertex and from chains of quintic vertices on the $\Sigma_2$ line, which the septic order must absorb in its turn (\cref{p:sec:hs:septic}).
Off-shell, $\caction[\le 6]$ contains the square of the fluctuation $\Sigma_3 - \Psi \otimes \Sigma_2$ with the field-independent operator $\wprod{3}{3}$.
Once $\Sigma_3$ is eliminated, what is left is quadratic in $\Sigma_2$, with the dressed operator $\kin{2}{\Psi}{3}$ corrected at order $\gs^2$ by the sextic term.
The gauge transformations built on these two fluctuations are the subject of \cref{p:sec:gauge}.


\paragraph{Consistency conditions on the kernels}
\label{p:sec:hs:sextic:kernel}

The quadratic term of $\Sigma_3$ in~\eqref{p:eq:hs6:S6tilde-explicit} is not given by the sextic vertex alone.
It contains the exchange part $\sprod{2}{3} \circ \sprod{2}{2}^{-1} \circ \sprod{3}{2}$ of~\eqref{p:eq:hs6:kinetic}, two quintic vertices joined by a $\Sigma_2$ line, read with three strands in and three strands out.
The kinetic operator of $\Sigma_3$ is thus non-local in the sense of effective field theory since, unlike that of $\Sigma_2$, it contains the inverse of a vertex.
A $\Sigma_2$ line collapses to a contact term if one of its ends is the source vertex $\sprod{2}{2}(\Psi^2)$, since $\sprod{2}{2}^{-1} \circ \sprod{2}{2}$ is the identity on the quotient, and, for generic vertices, in which the exchange term is not a linear combination of the elementary vertices, only then.
  The line of $\vertex{5 \cdot 5}$ joins two quintic vertices, so it is the first that does not collapse, and this is why $\vertex{4}$ and $\vertex{5}$ needed no reduction while $\vertex{6}$ does.

The interpretation is the same one as of the quintic order.
The field $\Sigma_2$ describes more than the quartic interaction once the quintic vertex is included: it induces interactions at all orders.
The part of the sextic interaction left for $\Sigma_3$ to describe is the sextic vertex minus what $\Sigma_2$ already produces.
We can make an analogy with linear algebra: $\wprod{3}{3}$ is the quadratic form left on $\Hn{3}$ once $\Sigma_2$ is eliminated from the sextic form, that is, its Schur complement.\footnote{
    For a block matrix with blocks $A$, $B$, $B^T$ and $D$, with $D$ invertible, the Schur complement of $D$ is $A - B D^{-1} B^T$: the quadratic form left on the first block once the second is eliminated by its equation of motion.
    Here $A = \sprod{3}{3}$, $B = \sprod{2}{3}$ and $D = \tfrac{5}{6} \, \sprod{2}{2}$.
  }
  The completed square~\eqref{p:eq:hs6:DeltaS6} is built on this form.

Two conditions make $\Sigma_3$ an auxiliary field, in addition to the invertibility of $\sprod{2}{2}$ on its quotient, which the exchange part needs.
First, $\wprod{3}{3}$ must be invertible on the quotient~\eqref{p:eq:hs6:quotient}, so that $\Sigma_3$ has a propagator and is integrated out by a Gaussian integral.
Second, its kernel must couple to nothing.
Injectivity, which holds on the quotient by construction, already makes the solution~\eqref{p:eq:hs6:saddle-sigma3} unique.
Unlike the dressed operator of $\Sigma_2$, $\wprod{3}{3}$ carries no power of $\gs$: its two terms are of the same order.
Its invertibility is therefore a genuine condition on the vertex family, on the sextic vertex and on the exchange diagram together, which no expansion in $\gs$ reduces to the condition on $\sprod{2}{2}$.
It holds in the finite-dimensional models of \cref{p:app:finite-models}, with generic vertices.
For the string vertices, it is a condition on the spectrum of $\sprod{3}{3}$ on $\Hn{3}$ and on the exchange part, which can be checked level by level in level truncation.
The kernels of $\sprod{3}{3}$ and $\wprod{3}{3}$ on $\Hn{3}$ are meant, as for $\sprod{2}{2}$, in the square-integrable sense: their elements are square-integrable functions of the relative momenta at fixed total momentum.
We assume it as we did for $\sprod{2}{2}$ in \cref{p:sec:hs:quartic,p:sec:hs:quintic}: kernels in the square-integrable sense at fixed total momentum, sources of $(\wprod{3}{3})^{-1}$ and of its dressed forms lying in the image of $\wprod{3}{3}$, and kernels that stabilize as the level grows, in the sense that the quotient and the inverse have a limit at generic total momentum:
\begin{equation}
\label{p:eq:hs6:invertibility-3}
\big( \wprod{3}{3} \big)^{-1} \circ \wprod{3}{3}
  = \id
\quad \text{on } \Hn{3} / \ker \wprod{3}{3}.
\end{equation}

A kernel direction of $\wprod{3}{3}$ is a flat direction of $\Sigma_3$, with zero action.
At this order it decouples automatically, since $\Sigma_3$ enters only through $\wprod{3}{3}$.
At the next order $\Sigma_3$ enters the septic vertex.
The quotient then stays consistent only if the kernel directions are null vectors of every vertex in which $\Sigma_3$ appears, as the quintic vertex had to vanish on $\ker \sprod{2}{2}$ by~\eqref{p:eq:hs5:kernel-condition}.
The construction needs this assumption on the kernel of $\wprod{3}{3}$ from the septic order on, and we state it there (\cref{p:sec:hs:septic}).

We have taken the auxiliary fields on quotients, $\Sigma_2$ on $\Hn{2} / \ker \sprod{2}{2}$ and $\Sigma_3$ on $\Hn{3} / \ker \wprod{3}{3}$, and the action must be a function on them.
  For $\Sigma_3$ nothing is required at this order, since the field enters only through $\wprod{3}{3}$, whose kernel is quotiented.
  For $\Sigma_2$, the quintic coupling descends to the quotient if and only if the quintic vertex vanishes on the kernel,~\eqref{p:eq:hs5:kernel-condition}.
  The sextic term adds one coupling, $\dvertex{6}(\Psi \otimes \Sigma_2, \Sigma_3)$.
  Under $\Sigma_2 \to \Sigma_2 + N_2$ with $N_2 \in \ker \sprod{2}{2}$, its exchange part shifts by a term containing $\sprod{3}{2}(\Psi, N_2)$, which vanishes by~\eqref{p:eq:hs5:kernel-condition}, and its contact part by $\vertex{6}(\Psi, N_2, \Sigma_3)$.
  The cubic action therefore descends to the quotient if and only if the sextic vertex vanishes on the kernel as well:
  \begin{equation}
  \label{p:eq:hs6:kernel-condition}
  \forall N_2 \in \ker \sprod{2}{2}, \;
  \forall A \in \Hn{3}, \;
  \forall B \in \Hs:
  \quad
  \vertex{6}(N_2, A, B)
    = 0.
\end{equation}
  The quotient is thus consistent through this order when the quartic form has no kernel or, if it has one, when the quintic and sextic vertices vanish on it.
  When the vertices do not vanish on the kernel of $\sprod{2}{2}$, the quotient is not available: the quintic coupling depends on the representative of $\Sigma_2$, a complement of the kernel has to be chosen, and the eliminated action then depends on the complement, already at order $\gs^3$.
  The constraint can then be enforced by a Lagrange multiplier instead: on the kernel component only, the hybrid below, or on the whole field, the multiplier form of \cref{p:sec:all-orders}.
  The condition holds trivially at generic momentum if the quartic form has no kernel there.
  We do not try to settle here whether it holds for the actual vertex family; in principle, this can be checked level by level in level truncation.
  The quotient needs no inner product and no choice: it is canonical, and the kernel directions, which have zero action and couple to nothing, are not fields.



The second block of the sextic order gives the gauge transformations of $\Sigma_3$ and $\Sigma_2$ in the paper's presentation, the one with the reduced vertex, and the localization of the action by a second $2$-string field, which \cref{ch:partner-field} develops.
The gauge law in the elementary presentation, and the comparison of the two, are in \cref{sec:order:sextic:both-gauge}.



\paragraph{Gauge transformation}
\label{p:sec:hs:sextic:gauge}

Under the gauge transformation~\eqref{p:eq:hs1:gauge-psi} the truncated action is invariant up to the order of the truncation, $\delta_\Lambda S_{\le 6} = O(\gs^5)$, which should be reproduced by the cubic action.
The two auxiliary fields enter $\caction[\le 6]$ through two completed squares, and we read off their transformation laws by demanding that each square be invariant, as at the quintic order.

The $3$-string field enters only through $\Delta S_6$, which is the square of the fluctuation:
\begin{equation}
  \label{p:eq:hs6:deviation-3}
  \sigma_3^{(4)}
    := \Sigma_3 - \Sigma_3^{(4)*}
    = \Sigma_3 - \Psi \otimes \Sigma_2,
  \qquad
  \Delta S_6
    = - \frac{\gs^4}{6!} \, \pair[3]{\sigma_3^{(4)}}{\wprod{3}{3}(\sigma_3^{(4)})},
\end{equation}
with the field-independent kinetic operator~\eqref{p:eq:hs6:kinetic}.
Its variation is:
\begin{equation}
  \label{p:eq:hs6:square-variation-3}
  \delta_\Lambda \Delta S_6
    = - \frac{2 \gs^4}{6!} \, \pair[3]{\delta_\Lambda \sigma_3^{(4)}}{\wprod{3}{3}(\sigma_3^{(4)})},
\end{equation}
with no term from the operator.
It vanishes for all $\sigma_3^{(4)}$ when $\delta_\Lambda \sigma_3^{(4)}$ lies in the kernel of $\wprod{3}{3}$, that is, when it vanishes on the quotient~\eqref{p:eq:hs6:quotient}.
We take this transformation law.
As for $\Sigma_2$ below, it is unique up to the trivial symmetries of \cref{p:sec:gauge}.

The solution~\eqref{p:eq:hs6:saddle-sigma3} and the propagator of $\Sigma_3$ rest on the invertibility of $\wprod{3}{3}$ on the quotient, assumed in~\eqref{p:eq:hs6:invertibility-3}.
We thus take the fluctuation of the $3$-string field inert, and $\Sigma_3$ transforms as its solution:
\begin{equation}
  \label{p:eq:hs6:gauge-sigma3}
  \delta_\Lambda \sigma_3^{(4)}
    = 0,
  \qquad
  \delta_\Lambda \Sigma_3
    = \delta_\Lambda \Psi \otimes \Sigma_2 + \Psi \otimes \delta_\Lambda \Sigma_2.
\end{equation}
There is no rotation term: the rotation at quintic order compensated the variation of the dressed operator, which depends on $\Psi$, and $\wprod{3}{3}$ does not depend on $\Psi$.
This is the quartic mechanism one level up, with $\Psi \otimes \Sigma_2$ in place of $\Psi^2$.

Once $\Sigma_3$ is eliminated, the action~\eqref{p:eq:hs6:eliminate-sigma3} is quadratic in $\Sigma_2$, and its quadratic part defines the dressed kinetic operator of $\Sigma_2$ at this order:
\begin{equation}
  \label{p:eq:hs6:dressed-2}
  \pair[2]{A}{\kin{2}{\Psi}{4}(B)}
    := \pair[2]{A}{\kin{2}{\Psi}{3}(B)}
      - \frac{\gs^2}{30} \, \dvertex{6}(\Psi \otimes A, \Psi \otimes B),
\end{equation}
for $A, B \in \Hn{2}$, so that the terms of $\caction[\le 6] |_{\Sigma_3 = \Sigma_3^{(4)*}}$ quadratic in $\Sigma_2$ read
\begin{equation}
  - \frac{\gs^2}{4!} \, \pair[2]{\Sigma_2}{\kin{2}{\Psi}{4}(\Sigma_2)},
\end{equation}
as in~\eqref{p:eq:hs5:F-dressed}.
The correction is the $\Sigma_3$ exchange between two of the cubic couplings $\dvertex{6}(\Psi \otimes \Sigma_2, \Sigma_3)$, collapsed to a contact term by the elimination of $\Sigma_3$.
The operator is symmetric with respect to the $2$-string pairing and depends on $\Psi$ through both terms.
The solution and the fluctuation are defined as at the quintic order:
\begin{equation}
  \label{p:eq:hs6:saddle-sigma2}
  \Sigma_2^{(4)*}
    := \big( \kin{2}{\Psi}{4} \big)^{-1} \circ \sprod{2}{2}(\Psi^2),
  \qquad
  \sigma_2^{(4)}
    := \Sigma_2 - \Sigma_2^{(4)*},
\end{equation}
the inverse being a Neumann series in $\gs$ around $\sprod{2}{2}^{-1}$ with one more kind of insertion, the order-$\gs^2$ term of~\eqref{p:eq:hs6:dressed-2}.
Completing the square as in~\eqref{p:eq:hs5:square}, we find that the cubic action is quadratic in $\sigma_2^{(4)}$ with the operator $\kin{2}{\Psi}{4}$.
We demand that this quadratic form be invariant.
Its variation has a term from the fluctuation and one from the $\Psi$ inside the operator, and it vanishes for:
\begin{equation}
  \label{p:eq:hs6:gauge-sigma2}
  \delta_\Lambda \sigma_2^{(4)}
    = - \frac{1}{2} \, \big( \kin{2}{\Psi}{4} \big)^{-1} \circ \big( \delta_\Lambda \kin{2}{\Psi}{4} \big) (\sigma_2^{(4)}),
  \qquad
  \delta_\Lambda \Sigma_2
    = \delta_\Lambda \Sigma_2^{(4)*} + \delta_\Lambda \sigma_2^{(4)},
\end{equation}
where $\delta_\Lambda \kin{2}{\Psi}{4}$ is the variation of the operator~\eqref{p:eq:hs6:dressed-2} through its argument $\Psi$, by the symmetry of the operator.
At order $\gs$ the variation of the operator is $\sprod{3}{2}(\delta_\Lambda \Psi, \cdot)$, and~\eqref{p:eq:hs6:gauge-sigma2} is the transformation law~\eqref{p:eq:hs5:gauge-deviation} of the quintic order.
The order-$\gs^2$ term is the variation of the sextic correction.
This is the transformation law of \cref{p:sec:gauge}, written there as~\eqref{p:eq:gauge:law-quintic} with the dressed operator of the order at hand.
It has the structure of the quintic one: the auxiliary field transforms as its solution, plus a rotation of the fluctuation which is linear in it, carries no parameter of its own and vanishes on-shell.
The transformation law is unique up to the trivial symmetries of \cref{p:sec:gauge}.

With the transformation laws~\eqref{p:eq:hs6:gauge-sigma3} and~\eqref{p:eq:hs6:gauge-sigma2}, both squares are exactly invariant.
The off-shell identity behind~\eqref{p:eq:hs6:eliminate}, the analogue of~\eqref{p:eq:hs5:identity}, reads:
\begin{equation}
\caction[\le 6]
  = S_{\le 6}
    + (\text{exchange terms of order } \gs^5)
    + (\text{squares}).
\end{equation}
It gives:
\begin{equation}
  \label{p:eq:hs6:delta-S-equal}
  \delta_\Lambda \caction[\le 6] - \delta_\Lambda S_{\le 6}
    = O(\gs^5),
\end{equation}
identically, off-shell, for any $\delta_\Lambda \Psi$: the cubic action is invariant to the order of the truncation it reproduces, as at the two previous orders.
The algebra of these transformations is that of closed string field theory through this order (\cref{p:sec:gauge}).


\paragraph{A second \texorpdfstring{$2$}{2}-string field}
\label{p:sec:hs:sextic:partner}

\subparagraph{Localizing the action.}

The reduced vertex is one way of subtracting the exchange term.
  Another way keeps every vertex elementary, at the price of a second $2$-string field.
  We give a field $\Upsilon_2 \in \Hn{2} / \ker \sprod{2}{2}$, with the ghost number and the parity of $\Sigma_2$, a quadratic term and a single elementary coupling, with coefficients $a$ and $b$:
  \begin{equation}
  \label{p:eq:hs6:partner}
  S_\Upsilon
    := \frac{\gs^2}{4!} \, a \, \vertex{4}(\Upsilon_2^2)
      + \frac{\gs^3}{5!} \, b \, \vertex{5}(\Upsilon_2, \Sigma_2, \Psi),
  \qquad
  \frac{b^2}{a}
    = 4.
\end{equation}
  Since $\Upsilon_2$ has no source at order $\gs^2$, its classical value starts at order $\gs$.
  Its own exchange, two quintic vertices joined by an $\Upsilon_2$ line, is $- b^2 / (4 a)$ times the exchange term of $\Sigma_2$.
  The two cancel for $b^2 / a = 4$, and we take $(a, b) = (1, 2)$, the $1 : 2$ ratio of the quadratic and source terms of $\Sigma_2$ in~\eqref{p:eq:hs4:S4tilde}, so that the quadratic form of $\Upsilon_2$ has the sign opposite to that of $\Sigma_2$.
  The action
  \begin{equation}
  \caction[\le 5] + S_\Upsilon + \frac{\gs^4}{6!} \, \vertex{6}(\Sigma_2^3)
\end{equation}
  then reproduces $S_{\le 6}$ through order $\gs^4$ with the fields $(\Psi, \Sigma_2, \Upsilon_2)$, elementary vertices only, and no $\Sigma_3$.
  This holds at tree level, for generic off-shell $\Psi$, under the same assumptions as the construction with $\Sigma_3$, and it was checked in the finite-dimensional models of \cref{p:app:finite-models} with generic vertices.
  For a real field the opposite sign of the quadratic form is required, and the sign is harmless at tree level.
  We leave open whether the sign has a physical consequence beyond tree level, $\Upsilon_2$ being an auxiliary field without propagating states, as well as the gauge transformation of $\Upsilon_2$ and the completed-square form of its sector.
  We give the construction in \cref{p:app:wrong-sign}.

A same-sign field $\Xi_2$ can instead absorb the inverse of $\sprod{2}{2}$ in the kinetic term of $\Sigma_3$, as the Liouville field makes the induced gravitational action local.
It cannot do the same in the coupling of $\Sigma_3$ to $(\Psi, \Sigma_2)$ at the same time: one field cannot serve two terms without producing a cross term.
A fully local action would need several of them, as in the intermediate-field decompositions of~\cite{Rivasseau:2010:LoopVertexExpansion,Lionni:2018:NoteIntermediateField}.
This too is in \cref{p:app:wrong-sign}.


At the septic order the $3$-string field becomes indispensable: the auxiliary fields are eliminated one after the other by exact Gaussian integrals, and the cubic coupling of $\vertex{7}$ is fixed by cancelling the interactions this elimination induces.
The closure at this order, by sequential exact Gaussian elimination, has status \status{Conditional(PH1, PH2 (b), PH5, PH14 [$O(g_s^4)$, $O(g_s^5)$]) + R70 -- Proven}\record{018-b6-ogs5-septic-closure (R25)}.

\subsubsection{Septic order, \texorpdfstring{$O(\gs^5)$}{O(gs**5)}}
\label{p:sec:hs:septic}

The septic vertex $\vertex{7}$ is the first one that cannot be written with $\Sigma_2$ alone.
This is the order at which the $3$-string field introduced in \cref{p:sec:hs:sextic} becomes indispensable, and at which the pattern repeated at all orders in \cref{p:sec:all-orders} becomes visible.
We first display the structure of the action obtained so far, which allows the auxiliary fields to be eliminated one after the other by Gaussian integrals, exact at tree level.
We then compute the interactions this elimination induces at order $\gs^5$, and fix the cubic coupling of $\vertex{7}$ that cancels them.

\paragraph{Elimination of the auxiliary fields}
\label{p:sec:hs:septic:structure}

\subparagraph{Triangular form of the action.}

The terms of $\caction[\le 6]$ which contain an auxiliary field are, by~\eqref{p:eq:hs5:F} and~\eqref{p:eq:hs6:S6tilde}:
\begin{equation}
  \label{p:eq:hs7:action-6}
  \begin{aligned}
    \caction[\le 6] - S_{\le 3}
      &
      = - \frac{\gs^2}{4!} \, \vertex{4}(\Sigma_2^2)
        + \frac{2 \gs^2}{4!} \, \vertex{4}(\Sigma_2, \Psi^2)
        + \frac{\gs^3}{5!} \, \vertex{5}(\Sigma_2^2, \Psi)
      \\ &
      \quad
      - \frac{\gs^4}{6!} \, \pair[3]{\Sigma_3}{\wprod{3}{3}(\Sigma_3)}
        + \frac{2 \gs^4}{6!} \, \pair[3]{\Sigma_3}{\wprod{3}{3}(\Psi \otimes \Sigma_2)},
  \end{aligned}
\end{equation}
where $S_{\le 3}$ is the kinetic term and the cubic vertex of $\Psi$, common to the cubic action and to the target, and where the sextic term is written through the kinetic operator $\wprod{3}{3}$ of~\eqref{p:eq:hs6:kinetic}.
The rest of the construction relies on two features of this expression.
First, the action is quadratic in each auxiliary field separately: in $\Sigma_3$, with the field-independent operator $\wprod{3}{3}$, and in $\Sigma_2$, with the dressed operator $\kin{2}{\Psi}{3}$ of~\eqref{p:eq:hs5:dressed}.
Second, the source of $\Sigma_3$ is linear in $\Sigma_2$ and $\Psi$.
The source of $\Sigma_2$ involves $\Psi$ only, apart from the coupling to $\Sigma_3$, which is the source of $\Sigma_3$ read the other way.
Once the higher-order auxiliary field $\Sigma_3$ is eliminated, the action is again quadratic in $\Sigma_2$, with a source that depends on $\Psi$ only.
We can therefore eliminate the auxiliary fields one at a time, in decreasing order, $\Sigma_3$ first, and each elimination is a Gaussian integral, exact at tree level.
The value of the action on the solution then comes out in closed form rather than order by order.
We call this the triangular form of the action.


\subparagraph{Eliminating the $3$-string field.}

We eliminated the $3$-string field in~\eqref{p:eq:hs6:eliminate-sigma3}, exactly in $\gs$ and for any $\Sigma_2$.
The solution is $\Sigma_3^{(4)*} = \Psi \otimes \Sigma_2$, and the value of the $\Sigma_3$ terms on it is:
\begin{equation}
  \label{p:eq:hs7:sigma3-value}
    \pair[3]{\wprod{3}{3}(\Psi \otimes \Sigma_2)}{\big( \wprod{3}{3} \big)^{-1} \circ \wprod{3}{3}(\Psi \otimes \Sigma_2)}
    = \dvertex{6}(\Psi \otimes \Sigma_2, \Psi \otimes \Sigma_2),
\end{equation}
which is exactly quadratic in $\Sigma_2$, since the source is linear in $\Sigma_2$.
The inverse of $\wprod{3}{3}$ is taken on the quotient~\eqref{p:eq:hs6:quotient}, under the assumption~\eqref{p:eq:hs6:invertibility-3}.
The action left for $\Sigma_2$ is therefore quadratic too, since the quintic vertex has only two $\Sigma_2$ slots, and the second elimination is again a Gaussian integral.
We expand around the quartic-order solution, $\Sigma_2 = \Psi^2 + \sigma_2^{(2)}$, which does not depend on $\gs$; the solution of this order then follows by completing the square in the fluctuation $\sigma_2^{(2)}$.
The terms of $\caction[\le 6]$ with $\Sigma_3$ eliminated then sort into three groups:
\begin{subequations}
\label{p:eq:hs7:sigma2-groups}
\begin{align}
  \label{p:eq:hs7:sigma2-contact}
  \caction[\le 6] \big|_{\Sigma_3 = \Sigma_3^{(4)*}} - S_{\le 3}
    &
    = \frac{\gs^2}{4!} \, \vertex{4}(\Psi^4)
      + \frac{\gs^3}{5!} \, \vertex{5}(\Psi^5)
      + \frac{\gs^4}{6!} \, \dvertex{6}(\Psi^6)
  \\
  \label{p:eq:hs7:sigma2-quadratic}
  &
  \quad
  - \frac{\gs^2}{4!} \, \pair[2]{\sigma_2^{(2)}}{\kin{2}{\Psi}{4}(\sigma_2^{(2)})}
  \\
  \label{p:eq:hs7:sigma2-linear}
  &
  \quad
  + \frac{\gs^3}{60} \, \Big[
    \vertex{5}(\sigma_2^{(2)}, \Psi^2, \Psi)
    + \frac{\gs}{6} \, \dvertex{6}(\Psi^3, \Psi \otimes \sigma_2^{(2)})
  \Big],
\end{align}
\end{subequations}
where the first line collects the contact terms, the second is the quadratic form of the fluctuation and the third is linear in it.
The contact terms are the target through order $\gs^4$ except for the sextic deficit: $\dvertex{6}(\Psi^6)$ differs from $\vertex{6}(\Psi^6)$ by $- \tfrac{6}{5} \vertex{5 \cdot 5}(\Psi^6)$, which the exchange term of the quintic order compensates, as in~\eqref{p:eq:hs6:eliminate}.
The quadratic form is that of the dressed kinetic operator $\kin{2}{\Psi}{4}$ of~\eqref{p:eq:hs6:dressed-2}.
It is the quartic form dressed by the quintic vertex, as in~\eqref{p:eq:hs5:dressed}, and by the exchange of $\Sigma_3$ between two sextic couplings: by~\eqref{p:eq:hs7:sigma3-value}, this exchange is the last term of $\kin{2}{\Psi}{4}$.
The linear term defines the source of the fluctuation, the $2$-string state $\beta_\Psi \in \Hn{2}$ such that
\begin{equation}
  \label{p:eq:hs7:source-2}
  \pair[2]{\beta_\Psi}{A}
    := \vertex{5}(A, \Psi^2, \Psi) + \frac{\gs}{6} \, \dvertex{6}(\Psi^3, \Psi \otimes A),
  \qquad
  \beta_\Psi
    = \sprod{3}{2}(\Psi^3) + O(\gs),
\end{equation}
so that~\eqref{p:eq:hs7:sigma2-linear} reads $\tfrac{\gs^3}{60} \pair[2]{\beta_\Psi}{\sigma_2^{(2)}}$.
Completing the square, we get the value of $\caction[\le 6]$ on the joint solution as a single expression, exact in $\gs$:
\begin{equation}
  \label{p:eq:hs7:closed-form}
  \caction[\le 6] \big|_{\Sigma_3 = \Sigma_3^{(4)*}, \, \Sigma_2 = \Sigma_2^{(4)*}} - S_{\le 6}
    = \frac{\gs^4}{600} \Big[ \pair[2]{\beta_\Psi}{\big( \kin{2}{\Psi}{4} \big)^{-1}(\beta_\Psi)} - \vertex{5 \cdot 5}(\Psi^6) \Big],
\end{equation}
where the last term is the sextic deficit of~\eqref{p:eq:hs7:sigma2-contact}.
The right-hand side is a single function of $\gs$, rational at any finite level truncation: it describes the the whole sequence of interactions induced by the two auxiliary fields.
It is as explicit as the inverse of the dressed operator, which is a Neumann series for small $\gs$ as in \cref{p:sec:hs:quintic}.
It holds at finite $\gs$ when $\wprod{3}{3}$ and $\kin{2}{\Psi}{4}$ are invertible on their quotients.
We checked~\eqref{p:eq:hs7:closed-form} against a direct solution of the joint equations of motion in the finite-dimensional models of \cref{p:app:finite-models}, with generic vertices.

\subparagraph{Induced interactions at order $\gs^5$.}

We expand~\eqref{p:eq:hs7:closed-form} to first order in $\gs$ beyond the sextic order, with the inverse of the dressed operator:
\begin{equation}
\big( \kin{2}{\Psi}{4} \big)^{-1}
  = \sprod{2}{2}^{-1}
    + \frac{\gs}{5} \, \sprod{2}{2}^{-1} \circ \sprod{3}{2}(\Psi, \cdot) \circ \sprod{2}{2}^{-1}
    + O(\gs^2).
\end{equation}
The $\Sigma_3$ exchange inside $\kin{2}{\Psi}{4}$ enters only at order $\gs^6$.
We get:
\begin{equation}
  \label{p:eq:hs7:induced}
  \caction[\le 6] \big|_{\Sigma_3 = \Sigma_3^{(4)*}, \, \Sigma_2 = \Sigma_2^{(4)*}}
    = S_{\le 6}
      + \frac{\gs^5}{600} \Big[ \frac{1}{3} \, \mathcal S_2(\Psi^7) - \frac{1}{5} \, \mathcal S_1(\Psi^7) \Big]
      + O(\gs^6),
\end{equation}
with the two seven-point structures
\begin{subequations}
  \label{p:eq:hs7:structures}
  \begin{align}
    \label{p:eq:hs7:S1}
    \mathcal S_1(\Psi^7)
      &
      := \vertex{5}\Big( \sprod{2}{2}^{-1} \circ \sprod{3}{2}(\Psi^3), \, \sprod{2}{2}^{-1} \circ \sprod{3}{2}(\Psi^3), \, \Psi \Big),
    \\
    \label{p:eq:hs7:S2}
    \mathcal S_2(\Psi^7)
      &
      := \vertex{6}\Big( \Psi^4, \, \sprod{2}{2}^{-1} \circ \sprod{3}{2}(\Psi^3) \Big).
  \end{align}
\end{subequations}
Both are exchange diagrams of $\Sigma_2$ between interaction vertices, and neither collapses, by the criterion of \cref{p:sec:hs:sextic}.
The first is a chain of three quintic vertices joined by two $\Sigma_2$ lines, the two end vertices carrying three external $\Psi$ each and the middle one a single $\Psi$.
It is the next term of the Neumann series~\eqref{p:eq:hs5:saddle-series} inserted in the exchange term of~\eqref{p:eq:hs5:identity}, which gives $+ \tfrac{\gs^5}{3000} \mathcal S_1$.
The second is a sextic vertex with four external $\Psi$ joined by a $\Sigma_2$ line to a quintic vertex with three.
It comes from the reduced sextic vertex evaluated on the quintic-order solution, that is, from the $\Sigma_3$ exchange~\eqref{p:eq:hs7:sigma3-value} with one $\Sigma_2$ slot of $\dvertex{6}(\Psi \otimes \Sigma_2, \Psi \otimes \Sigma_2)$ receiving the correction $\tfrac{\gs}{5} \sprod{2}{2}^{-1} \circ \sprod{3}{2}(\Psi^3)$ of~\eqref{p:eq:hs5:saddle-series}.
This cross term is $\tfrac{2 \gs^5}{5 \cdot 6!} \dvertex{6}(\Psi^3, \Psi \otimes B)$ with $B := \sprod{2}{2}^{-1} \circ \sprod{3}{2}(\Psi^3)$.
By $\dvertex{6}(\Psi^3, \Psi \otimes B) = \vertex{6}(\Psi^4, B) - \tfrac{6}{5} \vertex{5 \cdot 5}(\Psi^3, \Psi \otimes B)$, whose exchange part is $\mathcal S_1$ again, it equals $\tfrac{\gs^5}{1800} [\mathcal S_2 - \tfrac{6}{5} \mathcal S_1]$.
The two contributions add up to~\eqref{p:eq:hs7:induced}.
No other term enters at order $\gs^5$: the correction of order $\gs^2$ to the solution for $\Sigma_2$ from the sextic term contributes only at order $\gs^6$, by stationarity of $\caction[\le 5]$ at its own solution, as at the sextic order.
In both structures, each triple of external legs is attached to one and the same vertex.
  We use this property to build the reduced vertex below.

\paragraph{Septic coupling}
\label{p:sec:hs:septic:dressed}

\subparagraph{Reduced septic vertex.}


The septic order must reproduce $\vertex{7}$ and cancel the induced term of~\eqref{p:eq:hs7:induced} at the same time.
A term added at order $\gs^5$ contributes to the eliminated action at that order through its value on the diagonal $\Sigma_k = \Psi^k$, where it must be of degree seven in $\Psi$, and it must be cubic in the fields.
The only partitions of seven strands over three fields drawn from $\{\Psi, \Sigma_2, \Sigma_3\}$ are $7 = 3 + 3 + 1$ and $7 = 3 + 2 + 2$.
We take the first partition, in which the $3$-string field plays the role $\Sigma_2$ played in the quintic coupling:
\begin{equation}
  \label{p:eq:hs7:S7tilde}
  \caction[7]
    := \frac{\gs^5}{7!} \, \dvertex{7}(\Sigma_3, \Sigma_3, \Psi),
  \qquad
  \caction[\le 7]
    := \caction[\le 6] + \caction[7],
\end{equation}
with a trilinear form $\dvertex{7}$ on $\Hn{3} \times \Hn{3} \times \Hs$, symmetric in its two $3$-string arguments, which the matching determines.
  We discuss the second partition below.
We define the reduced septic vertex, for $A, B \in \Hn{3}$ and $C \in \Hs$:
\begin{equation}
  \label{p:eq:hs7:dressed-vertex}
  \begin{aligned}
    \dvertex{7}(A, B, C)
      &
      := \vertex{7}(A, B, C)
        + \frac{42}{25} \, \vertex{5}\Big( \sprod{2}{2}^{-1} \circ \sprod{3}{2}(A), \, \sprod{2}{2}^{-1} \circ \sprod{3}{2}(B), \, C \Big)
      \\ &
      \quad
      - \frac{14}{5} \cdot \frac{1}{2} \Big[
        \vertex{6}\Big( A, C, \sprod{2}{2}^{-1} \circ \sprod{3}{2}(B) \Big)
        + \vertex{6}\Big( B, C, \sprod{2}{2}^{-1} \circ \sprod{3}{2}(A) \Big)
      \Big],
  \end{aligned}
\end{equation}
where $\sprod{3}{2}(A)$ fills the three input slots of $\sprod{3}{2}$ with the $3$-string state $A$, and where $\vertex{6}(A, C, D)$ fills three slots with $A$, one with $C$ and two with $D \in \Hn{2}$.
The second term is the chain of three quintic vertices of $\mathcal S_1$, the two end vertices read with their three single-string slots as one $\Sigma_3$ slot.
The third term is the sextic--quintic exchange of $\mathcal S_2$, with one $\Sigma_3$ in the sextic vertex and the other in the quintic one, symmetrized.
On the diagonal
\begin{equation}
  \dvertex{7}(\Psi^3, \Psi^3, \Psi)
    = \vertex{7}(\Psi^7)
      + \frac{42}{25} \, \mathcal S_1(\Psi^7)
      - \frac{14}{5} \, \mathcal S_2(\Psi^7),
\end{equation}
by~\eqref{p:eq:hs7:structures}.
\begin{figure}[htp]
  \centering
  \begin{tikzpicture}
    \begin{feynman}
      \vertex (a1) at (0, 0) {$\Sigma_3$};
      \vertex (a) at (1.2, 0) [dot];
      \vertex (b) at (2.8, 0) [dot];
      \vertex (c) at (4.4, 0) [dot];
      \vertex (c1) at (5.6, 0) {$\Sigma_3$};
      \vertex (p) at (2.8, 1.2) {$\Psi$};
      \diagram* {
      (a1) -- [plain, double distance=2.5pt] (a) -- [plain, double, edge label'={$\Sigma_2$}] (b) -- [plain, double, edge label'={$\Sigma_2$}] (c) -- [plain, double distance=2.5pt] (c1),
      (b) -- [plain] (p),
      };
      \vertex (d1) at (7.5, 0) {$\Sigma_3$};
      \vertex (d) at (8.7, 0) [dot];
      \vertex (e) at (10.3, 0) [dot];
      \vertex (e1) at (11.5, 0) {$\Sigma_3$};
      \vertex (q) at (8.7, 1.2) {$\Psi$};
      \diagram* {
      (d1) -- [plain, double distance=2.5pt] (d) -- [plain, double, edge label'={$\Sigma_2$}] (e) -- [plain, double distance=2.5pt] (e1),
      (d) -- [plain] (q),
      };
    \end{feynman}
  \end{tikzpicture}
  \caption{The two exchange terms of the reduced septic vertex~\eqref{p:eq:hs7:dressed-vertex}: a thick double line is a $\Sigma_3$ leg (three strands), a double line a $\Sigma_2$ line carrying $\sprod{2}{2}^{-1}$. Left, three quintic vertices on two $\Sigma_2$ lines; right, a sextic vertex and a quintic vertex on one $\Sigma_2$ line. In both, the three strands of each $\Sigma_3$ leg enter one and the same vertex.}
  \label{p:fig:hs7:dressed-vertex}
\end{figure}

The elimination proceeds as before, $\Sigma_3$ first.
The septic coupling is quadratic in $\Sigma_3$ and does not change its source, so the $3$-string sector stays Gaussian, with the dressed kinetic operator:
\begin{equation}
  \label{p:eq:hs7:dressed-3}
  \pair[3]{A}{\kin{3}{\Psi}{5}(B)}
    := \pair[3]{A}{\wprod{3}{3}(B)} - \frac{\gs}{7} \, \dvertex{7}(A, B, \Psi),
  \qquad
  A, B \in \Hn{3}.
\end{equation}
The equation of motion of $\Sigma_3$ and its solution are:
\begin{equation}
  \label{p:eq:hs7:saddle-3}
  \begin{aligned}
    \mathcal E^{(5)}_{\Sigma_3}
      &
      := \wprod{3}{3}(\Psi \otimes \Sigma_2) - \kin{3}{\Psi}{5}(\Sigma_3)
      = 0,
    \\
    \Sigma_3^{(5)*}
      &
      := \big( \kin{3}{\Psi}{5} \big)^{-1} \circ \wprod{3}{3}(\Psi \otimes \Sigma_2),
  \end{aligned}
\end{equation}
the inverse existing as a Neumann series in $\tfrac{\gs}{7} (\wprod{3}{3})^{-1} \circ \dvertex{7}(\, \cdot \,, \cdot \,, \Psi)$ for small $\gs$.
The solution is now $\Psi \otimes \Sigma_2$ plus a correction of order $\gs$, as $\Sigma_2^{(3)*}$ at the quintic order.
The value of the $\Sigma_3$ terms on the solution is the completed square:
\begin{equation}
  \label{p:eq:hs7:sigma3-value-5}
  \begin{aligned}
    &
    \frac{\gs^4}{6!} \, \pair[3]{\wprod{3}{3}(\Psi \otimes \Sigma_2)}{\big( \kin{3}{\Psi}{5} \big)^{-1} \circ \wprod{3}{3}(\Psi \otimes \Sigma_2)}
    \\ &
    \qquad
    = \frac{\gs^4}{6!} \, \dvertex{6}(\Psi \otimes \Sigma_2, \Psi \otimes \Sigma_2)
      + \frac{\gs^5}{7!} \, \dvertex{7}(\Psi \otimes \Sigma_2, \Psi \otimes \Sigma_2, \Psi)
      + O(\gs^6),
  \end{aligned}
\end{equation}
where the first term of the Neumann series gives back~\eqref{p:eq:hs7:sigma3-value} and the second is the septic coupling.
The order-$\gs^5$ contribution of the septic coupling on-shell is thus its naive contact value, as for the quintic coupling in~\eqref{p:eq:hs5:constant}.
The action left for $\Sigma_2$ is still exactly quadratic, since $\Sigma_3^{(5)*}$ is linear in $\Sigma_2$.
Its solution $\Sigma_2^{(5)*}$ differs from $\Sigma_2^{(4)*}$ at order $\gs^3$.
By stationarity this does not affect order $\gs^5$, and at that order the septic term enters only through its value at $\Sigma_2 = \Psi^2$.

The target at this order is:
  \begin{equation}
  S_{\le 7}
    = S_{\le 6} + \frac{\gs^5}{7!} \, \vertex{7}(\Psi^7).
\end{equation}
  When we add the contact value of the septic coupling to~\eqref{p:eq:hs7:induced}, the matching with this target requires the single condition:
\begin{equation}
  \label{p:eq:hs7:coefficient}
  \frac{\gs^5}{600} \Big[ \frac{1}{3} \, \mathcal S_2 - \frac{1}{5} \, \mathcal S_1 \Big]
    + \frac{\gs^5}{7!} \, \dvertex{7}(\Psi^7)
    = \frac{\gs^5}{7!} \, \vertex{7}(\Psi^7),
\end{equation}
that is,
\begin{equation}
  \label{p:eq:hs7:diagonal}
  \dvertex{7}(\Psi^7)
    = \vertex{7}(\Psi^7) + \frac{42}{25} \, \mathcal S_1(\Psi^7) - \frac{14}{5} \, \mathcal S_2(\Psi^7),
\end{equation}
which is the diagonal value of~\eqref{p:eq:hs7:dressed-vertex}.
The matching determines the three coefficients $(1, \tfrac{42}{25}, - \tfrac{14}{5})$, since $\vertex{7}(\Psi^7)$, $\mathcal S_1$ and $\mathcal S_2$ are independent functions of $\Psi$ in a finite-dimensional model with generic vertices.
We assume this independence for the string vertices, as at the sextic order.
  Changing one coefficient then leaves a seven-point term proportional to the corresponding invariant.
With these coefficients, we get:
\begin{equation}
  \label{p:eq:hs7:eliminate}
  \caction[\le 7] \big|_{\Sigma_3 = \Sigma_3^{(5)*}, \, \Sigma_2 = \Sigma_2^{(5)*}}
    = S_{\le 7} + O(\gs^6),
\end{equation}
and the septic order closes: the exchange diagrams induced at order $\gs^5$ are canceled by the exchange parts of $\dvertex{7}$, one structure at a time, and the septic vertex is reproduced.
We checked~\eqref{p:eq:hs7:eliminate} in the finite-dimensional models, where the residual scales as $\gs^6$ and the matching through order $\gs^4$ is untouched.
The remainder of order $\gs^6$ is the next induced interaction, which the octic order absorbs in its turn (\cref{p:sec:hs:octic}).

\subparagraph{Dressed kinetic operators.}

The septic coupling dresses the kinetic operator of $\Sigma_3$ with $\Psi$, through~\eqref{p:eq:hs7:dressed-3}, as the quintic coupling did for $\Sigma_2$.
We can read its two exchange parts in terms of the exchanges of the lower orders, and this reading shows the pattern of the higher orders.
The chain term of~\eqref{p:eq:hs7:dressed-vertex} enters $\kin{3}{\Psi}{5}$ with the coefficient $- \tfrac{6}{5} \cdot \tfrac{\gs}{5}$.
This is the order-$\gs$ term of the exchange part of $\wprod{3}{3}$ with the bare $\Sigma_2$ propagator replaced by the dressed one,
\begin{equation}
  \label{p:eq:hs7:running}
  \begin{aligned}
    \sprod{2}{3} \circ \big( \kin{2}{\Psi}{3} \big)^{-1} \circ \sprod{3}{2}
      &
      = \sprod{2}{3} \circ \sprod{2}{2}^{-1} \circ \sprod{3}{2}
        + \frac{\gs}{5} \, \sprod{2}{3} \circ \sprod{2}{2}^{-1} \circ \sprod{3}{2}(\Psi, \cdot) \circ \sprod{2}{2}^{-1} \circ \sprod{3}{2}
        \\ & \qquad
        + O(\gs^2),
  \end{aligned}
\end{equation}
by the Neumann series~\eqref{p:eq:hs5:neumann}.
At this order, the chain term of the dressing is thus the first Taylor coefficient of the operator obtained from $\wprod{3}{3}$ by replacing the bare $\Sigma_2$ propagator with the dressed one, and $\wprod{3}{3}$ is the value of that operator at $\gs = 0$.
The higher coefficients are not tested here.
The sextic--quintic term, whose coefficient $+ \tfrac{2 \gs}{5}$ follows from the matching, is the $\Sigma_2$ exchange between the contact part $\vertex{6}$ of the sextic source vertex of $\Sigma_3$ and the quintic vertex read as $\sprod{3}{2}(\Sigma_3)$.
Conversely, the dressed operator of $\Sigma_2$ at this order is read off the action once $\Sigma_3$ is eliminated with~\eqref{p:eq:hs7:sigma3-value-5}:
\begin{equation}
  \label{p:eq:hs7:dressed-2}
  \begin{aligned}
    \pair[2]{A}{\kin{2}{\Psi}{5}(B)}
      &
      := \vertex{4}(A, B)
        - \frac{\gs}{5} \, \vertex{5}(A, B, \Psi)
      \\ &
      \quad
      - \frac{\gs^2}{30} \, \pair[3]{\wprod{3}{3}(\Psi \otimes A)}{\big( \kin{3}{\Psi}{5} \big)^{-1} \circ \wprod{3}{3}(\Psi \otimes B)},
  \end{aligned}
\end{equation}
This is the operator $\kin{2}{\Psi}{4}$ of~\eqref{p:eq:hs6:dressed-2} with the $\Sigma_3$ propagator dressed in its turn.
Each order thus dresses the kinetic operators of the lower-order auxiliary fields by the exchanges of the higher-order ones, and the propagators inside these exchanges are themselves the dressed ones.
The diagonal value of every reduced vertex is fixed in the same way: it is the vertex of the target minus what the previous orders induce:
\begin{equation}
  \label{p:eq:hs7:deficit}
  \dvertex{n}(\Psi^n)
    = \vertex{n}(\Psi^n) - n! \, \big[ \gs^{n - 2} \big] \, \caction[\le n - 1] \big|_{\text{solution}},
\end{equation}
where the cubic action of the previous order is evaluated on the solution of its auxiliary fields and $[\gs^{n - 2}]$ extracts the coefficient of $\gs^{n - 2}$; the reduced sextic vertex~\eqref{p:eq:hs6:dressed-vertex} and the condition~\eqref{p:eq:hs7:coefficient} are its cases $n = 6$ and $n = 7$.

\paragraph{Choices in the construction}
\label{p:sec:hs:septic:alternatives}

\subparagraph{Choice of the partition.}

The matching~\eqref{p:eq:hs7:coefficient} constrains the septic term only through its value on the diagonal, a septic polynomial in $\Psi$.
Any trilinear form with this diagonal value closes the order, and two forms differing by one that vanishes on the diagonal give the same eliminated action at this order.
The exchange parts of~\eqref{p:eq:hs7:dressed-vertex} are unique only up to this freedom, as those of the reduced sextic vertex were, and the freedom now includes the partition of the strands.
The second partition, $7 = 3 + 2 + 2$, gives a coupling $\vertex{7}(\Sigma_3, \Sigma_2, \Sigma_2)$ which, reduced by the same two structures with the same diagonal value, closes the order as well by the argument just given.
So does any combination of the two partitions whose coefficients add up to one.
To realize $\mathcal S_1$ and $\mathcal S_2$ in the second partition, the two strands of a $\Sigma_2$ slot must be split between two constituent vertices.
We prefer the first partition for three reasons.
In the first partition all the strands of each slot sit inside one constituent vertex, as in \cref{p:fig:hs7:dressed-vertex}, so that each exchange part of the reduced vertex is one of the induced diagrams~\eqref{p:eq:hs7:structures} with its triples of legs read as $\Sigma_3$ slots, off the diagonal as well as on it.
The first partition also preserves the triangular form.
A coupling $\vertex{7}(\Sigma_3, \Sigma_2, \Sigma_2)$ makes the source of $\Sigma_3$ quadratic in $\Sigma_2$, the action left after eliminating $\Sigma_3$ becomes quartic in $\Sigma_2$, and the closed form~\eqref{p:eq:hs7:closed-form} is lost at the next order.
With the first partition, finally, $\Sigma_3$ comes with an even and an odd coupling, $\dvertex{6}$ and $\dvertex{7}$, as $\Sigma_2$ came with $\vertex{4}$ and $\vertex{5}$.
We therefore keep the first partition,~\eqref{p:eq:hs7:S7tilde} with the reduced vertex~\eqref{p:eq:hs7:dressed-vertex}.

\subparagraph{The necessity of the $3$-string field.}

We could ask whether the $3$-string field could be introduced at this order, where it is first indispensable, rather than at the sextic order.
Suppose that the sextic order had been closed with $\Sigma_2$ alone, with a sextic coupling in $\Sigma_2^3$, see \cref{p:sec:hs:sextic}, and that $\Sigma_3$ were introduced here with a kinetic term and a single coupling, $\vertex{7}(\Sigma_3, \Sigma_2, \Sigma_2)$ or $\vertex{7}(\Sigma_3, \Sigma_3, \Psi)$.
The only kinetic term for $\Sigma_3$ built from the vertices is of sextic type, $\vertex{6}(\Sigma_3^2)$ or its reduced form, and is of order $\gs^4$.
Without a source at that order, the solution for $\Sigma_3$ is of order $\gs$ and the value of the coupling on it is of order $\gs^6$ at least.
The contact term $\vertex{7}(\Psi^7)$ is then never produced at order $\gs^5$, for any kernel and any coefficient.
Within cubic actions in $\Psi$, $\Sigma_2$ and $\Sigma_3$ with kinetic terms of the same sign, a kinetic term and a source for $\Sigma_3$ at the sextic order, giving it a solution of order one, are therefore required before the septic order can close.
The completed square~\eqref{p:eq:hs6:DeltaS6} is the form we use.
Within this class, a sextic order closed with $\Sigma_2$ alone has to be rebuilt with $\Sigma_3$ as soon as $\vertex{7}$ is reached: that form reproduces the same eliminated action through order $\gs^4$ and does not extend beyond it, the distinction made in \cref{p:sec:hs:sextic}.

This is why we introduce each auxiliary field through a completed square at the even order $\gs^{2 k - 2}$ and let the odd vertex dress its kinetic operator at the next order.
The pattern is established here for $k = 3$ and repeated in \cref{p:sec:all-orders}.
The septic order thus introduces no new field: it couples $\Sigma_3$ beyond the completed square of the sextic order, and gives it its first dressed kinetic operator.

\subparagraph{Kernel condition.}

The septic coupling must be a function on the quotient~\eqref{p:eq:hs6:quotient} on which $\Sigma_3$ lives, as the quintic coupling had to be a function on $\Hn{2} / \ker \sprod{2}{2}$.
The sextic source $\wprod{3}{3}(\Psi \otimes \Sigma_2)$ pairs to zero with every kernel element by the symmetry of $\wprod{3}{3}$, so it descends automatically, whereas the septic coupling does not.
From this order on we make a further standing assumption, besides those of \cref{p:sec:hs:quintic,p:sec:hs:sextic}.
The kernel of the kinetic operator with which each auxiliary field is introduced, $\wprod{3}{3}$ for $\Sigma_3$ here and its analogue for $\Sigma_k$ in \cref{p:sec:all-orders}, consists of joint null vectors of every vertex in which that field appears.
It need only hold for one choice of the terms vanishing on the diagonal, up to which the reduced vertices are defined.
At this order the condition reads:
\begin{equation}
  \label{p:eq:hs7:kernel-condition}
  \forall N_3 \in \ker \wprod{3}{3}, \;
  \forall A \in \Hn{3}, \;
  \forall B \in \Hs:
  \quad
  \dvertex{7}(N_3, A, B)
    = 0,
\end{equation}
and it makes $\dvertex{7}(\, \cdot \,, \cdot \,, \Psi)$, hence $\kin{3}{\Psi}{5}$ and its inverse, descend to the quotient.
It holds trivially when $\wprod{3}{3}$ has no kernel, and it is a condition on the choice of the reduced vertex rather than on the vertices alone, since $\dvertex{7}$ may be changed by terms vanishing on the diagonal.
When it fails for some kernel direction $N_3$, that direction has no kinetic term at order $\gs^4$ and acquires a quadratic term from the background at order $\gs^5$, $\tfrac{\gs^5}{7!} \dvertex{7}(N_3, N_3, \Psi)$.
Integrating it out divides by this field-dependent mass, and the eliminated action acquires a term rational in $\Psi$, $- \tfrac{\gs^5}{7!} \dvertex{7}(N_3, \Psi^3, \Psi)^2 / \dvertex{7}(N_3, N_3, \Psi)$ when the denominator does not vanish, which no polynomial vertex can match.
We established this in a finite-dimensional model.
Under the assumption, the auxiliary fields reproduce a polynomial action order by order.
When it fails along a direction where the denominator does not vanish, the eliminated action is not polynomial.
We do not try to settle here whether the string vertices satisfy this condition.
\Cref{p:sec:all-orders} and \cref{p:sec:gauge} refer to it in this form.

\paragraph{Gauge transformation}
\label{p:sec:hs:septic:gauge}

\subparagraph{Rotation of the fluctuations.}

With the septic coupling, the action decomposes exactly, off-shell, into the target, the induced interactions and two squares:
\begin{equation}
  \label{p:eq:hs7:identity}
  \begin{aligned}
    \caction[\le 7]
      &
      = S_{\le 7}
        + \big( \text{exchange terms of order } \gs^6 \big)
      \\ &
      \quad
      - \frac{\gs^4}{6!} \, \pair[3]{\sigma_3^{(5)}}{\kin{3}{\Psi}{5}(\sigma_3^{(5)})}
      - \frac{\gs^2}{4!} \, \pair[2]{\sigma_2^{(5)}}{\kin{2}{\Psi}{5}(\sigma_2^{(5)})},
  \end{aligned}
\end{equation}
where the exchange terms of order $\gs^6$ are the induced interaction of~\eqref{p:eq:hs7:eliminate}, a function of $\Psi$ alone.
The fluctuations are:
\begin{equation}
  \label{p:eq:hs7:deviations}
  \sigma_3^{(5)}
    := \Sigma_3 - \Sigma_3^{(5)*},
  \qquad
  \sigma_2^{(5)}
    := \Sigma_2 - \Sigma_2^{(5)*},
\end{equation}
the first from the solution~\eqref{p:eq:hs7:saddle-3} at the given $\Sigma_2$, the second from the solution of the $\Sigma_2$ equation once $\Sigma_3$ is eliminated.
Both squares come with dressed kinetic operators.
As at the quintic order, we fix the gauge law by demanding that each square be exactly invariant.
Since the action is Gaussian in both fluctuations, the transformation laws are explicit rotations:
\begin{equation}
  \label{p:eq:hs7:gauge-deviations}
  \begin{aligned}
    \delta_\Lambda \sigma_3^{(5)}
      &
      = - \frac{1}{2} \, \big( \kin{3}{\Psi}{5} \big)^{-1} \circ \big( \delta_\Lambda \kin{3}{\Psi}{5} \big) (\sigma_3^{(5)})
      = \frac{\gs}{14} \, \big( \kin{3}{\Psi}{5} \big)^{-1} \circ \dvertex{7}(\, \cdot \,, \sigma_3^{(5)}, \delta_\Lambda \Psi),
    \\
    \delta_\Lambda \sigma_2^{(5)}
      &
      = - \frac{1}{2} \, \big( \kin{2}{\Psi}{5} \big)^{-1} \circ \big( \delta_\Lambda \kin{2}{\Psi}{5} \big) (\sigma_2^{(5)}),
  \end{aligned}
\end{equation}
where $\dvertex{7}(\, \cdot \,, \sigma_3^{(5)}, \delta_\Lambda \Psi)$ denotes the $3$-string state paired with the open slot, as in~\eqref{p:eq:hs6:eom}.
Each auxiliary field transforms as its solution plus the rotation of its fluctuation:
\begin{equation}
  \label{p:eq:hs7:gauge-fields}
  \delta_\Lambda \Sigma_3
    = \delta_\Lambda \Sigma_3^{(5)*} + \delta_\Lambda \sigma_3^{(5)},
  \qquad
  \delta_\Lambda \Sigma_2
    = \delta_\Lambda \Sigma_2^{(5)*} + \delta_\Lambda \sigma_2^{(5)}.
\end{equation}
The variation of $\kin{2}{\Psi}{5}$ collects those of the quintic term, of the sextic source $\wprod{3}{3}(\Psi \otimes \cdot)$ and of the inverse of $\kin{3}{\Psi}{5}$ in~\eqref{p:eq:hs7:dressed-2}.
With these transformation laws, varying~\eqref{p:eq:hs7:identity} gives:
\begin{equation}
  \label{p:eq:hs7:delta-S-equal}
  \delta_\Lambda \caction[\le 7] - \delta_\Lambda \big[ S_{\le 7} + \big( \text{exchange terms of order } \gs^6 \big) \big]
    = 0
\end{equation}
identically, off-shell, for any variation $\delta_\Lambda \Psi$ and exactly in $\gs$.
Since $\delta_\Lambda S_{\le 7} = O(\gs^6)$, the cubic action is invariant up to order $\gs^6$, the order of the truncation it reproduces.
We checked the identity, and the residuals left by dropping either rotation, in the finite-dimensional model.
At the sextic order, the fluctuation of $\Sigma_3$ needed no rotation, since its kinetic operator was independent of $\Psi$.
From this order on, both fluctuations rotate.
The rotations are homogeneous, $\delta_\Lambda X \propto X$, and vanish on-shell.
Their algebra and the Batalin--Vilkovisky structure of the cubic action with both auxiliary fields are worked out in \cref{p:sec:gauge} through the sextic order.
A flat representative for the square of $\sigma_3^{(5)}$, expected as for $\Sigma_2$, has not been constructed.

\subparagraph{Reality.}

The fields are subject to the conditions $\Sigma_3^\star = - \Sigma_3$ and $\Sigma_2^\star = \Sigma_2$ of \cref{p:app:conventions}, and the septic coupling is compatible with them for the reason given at the quintic and sextic orders.
The reduced vertex $\dvertex{7}$ is built from $\vertex{5}$, $\vertex{6}$, $\vertex{7}$, $\sprod{3}{2}$ and $\sprod{2}{2}^{-1}$ with real coefficients, so the septic term is real on real fields by the reality of the vertices~\eqref{p:eq:hs4:vertex-real}, and $\kin{3}{\Psi}{5}$ is a real form when $\Psi$ is real.
The map $(\kin{3}{\Psi}{5})^{-1} \circ \wprod{3}{3}$ then commutes with $\star$, by conjugating the relation $\pair[3]{A}{\kin{3}{\Psi}{5}(B')} = \dvertex{6}(A, B)$, with $B' := (\kin{3}{\Psi}{5})^{-1} \circ \wprod{3}{3}(B)$, that characterizes it, as at the quintic order.
Hence the solution $\Sigma_3^{(5)*}$ is real when $\Psi$ and $\Sigma_2$ are, and the rotations~\eqref{p:eq:hs7:gauge-deviations} map the real locus to itself.
The argument at every order is given in \cref{p:app:reality}.


The octic order repeats the sextic one one step up the tower, with the $4$-string field.
This block is the paper's octic order in full; the paper keeps only its general part, in its all-orders section.
It states its own scope: the rational term produced by a kernel direction that violates the kernel condition is checked in a finite-dimensional model at $k = 3$ and argued for general $k$.

\subsubsection{Octic order, \texorpdfstring{$O(\gs^6)$}{O(gs**6)}}
\label{p:sec:hs:octic}

The octic order repeats the sextic one with every field moved one order up the tower.
We introduce each auxiliary field by completing a square at the even order $\gs^{2 k - 2}$, as explained in \cref{p:sec:hs:septic}, and we reach $k = 4$ here.
We introduce the $4$-string field, write the dressed operators of the three auxiliary fields at this order, and show that they depend on the string field only through $\gs \Psi$.
We write the equations of motion of all the fields, and then state the condition on the kernels under which the auxiliary fields can still be taken on quotients, and what goes wrong when it fails.
The matching follows the rule~\eqref{p:eq:hs7:deficit} without change, and we state it at all orders in the next section.

\subparagraph{The $4$-string field.}

The remainder of order $\gs^6$ of~\eqref{p:eq:hs7:eliminate} collects the exchange diagrams of $\Sigma_2$ and $\Sigma_3$ between the quintic, sextic and septic vertices, among them the $\Sigma_3$ exchange inside $\kin{2}{\Psi}{4}$, which did not contribute at order $\gs^5$.
Its diagonal value fixes the reduced octic vertex by~\eqref{p:eq:hs7:deficit} with $n = 8$:
\begin{equation}
  \label{p:eq:hs8:diagonal}
  \dvertex{8}(\Psi^8)
    = \vertex{8}(\Psi^8) - 8! \, \big[ \gs^6 \big] \, \caction[\le 7] \big|_{\Sigma_3 = \Sigma_3^{(5)*}, \, \Sigma_2 = \Sigma_2^{(5)*}}.
\end{equation}
The right-hand side is an octic polynomial in $\Psi$.
Any symmetric form with this diagonal value closes the order.
We have not written out the form built from the diagrams themselves, with the strands of each slot inside one constituent vertex as in~\eqref{p:eq:hs7:dressed-vertex}.
  In the finite-dimensional model, we use the polarization of the diagonal.
Only this diagonal value enters the closure of the order, and \cref{p:sec:all-orders} shows that the same holds at every order.
A $4$-string field fills half of the eight slots.
We read $\dvertex{8}$ as a symmetric bilinear form on $\Hn{4}$, write the octic term with $\dvertex{8}(\Psi \otimes \Sigma_3, \Psi \otimes \Sigma_3)$, and complete its square with a $4$-string field:
\begin{equation}
  \label{p:eq:hs8:DeltaS8}
  \begin{gathered}
    \Delta S_8
      := - \frac{\gs^6}{8!} \, \dvertex{8}(\Sigma_4 - \Psi \otimes \Sigma_3, \Sigma_4 - \Psi \otimes \Sigma_3),
    \\
    \Sigma_4 \in \Hn{4},
    \qquad
    \Ngh(\Sigma_4)
      = 8.
  \end{gathered}
\end{equation}
We use the state $\Psi \otimes \Sigma_3$ instead of $\Psi^4$ for the reason given at the sextic order: the cross term of the square must be cubic in the fields, and $\Sigma_3 = \Psi^3 + O(\gs)$ on-shell guarantees that only $\dvertex{8}(\Psi^8)$ is produced at this order.
The state $\Sigma_2 \otimes \Sigma_2$ would also give a cubic cross term, but the source of $\Sigma_4$ would then be quadratic in $\Sigma_2$ rather than linear in the field just below it, and the action would leave the triangular form of \cref{p:sec:hs:septic}.
The octic term of the cubic action is then:
\begin{equation}
  \label{p:eq:hs8:S8tilde}
  \begin{aligned}
    \caction[8]
      &
      := \frac{\gs^6}{8!} \, \dvertex{8}(\Psi \otimes \Sigma_3, \Psi \otimes \Sigma_3) + \Delta S_8
    \\ &
      = - \frac{\gs^6}{8!} \, \dvertex{8}(\Sigma_4^2)
        + \frac{2 \gs^6}{8!} \, \dvertex{8}(\Psi \otimes \Sigma_3, \Sigma_4),
    \qquad
    \caction[\le 8]
      := \caction[\le 7] + \caction[8].
  \end{aligned}
\end{equation}
It contains a quadratic term for $\Sigma_4$ and a cubic coupling of $\Sigma_4$, $\Sigma_3$ and $\Psi$, whose slots split the eight strands as $8 = 4 + 3 + 1$.
As at the sextic order, we write the form through the $4$-string pairing,
\begin{equation}
  \label{p:eq:hs8:kinetic}
  \dvertex{8}(A, B)
    =: \pair[4]{A}{\kin{4}{}{6}(B)},
  \qquad
  A, B \in \Hn{4},
\end{equation}
which defines the kinetic operator $\kin{4}{}{6} \colon \Hn{4} \to \Hn{4}$ of the $4$-string field, symmetric with respect to the $4$-string pairing.
This operator does not depend on the fields, since the odd vertex $\dvertex{9}$ which dresses it is of order $\gs^7$ and no higher-order field has been introduced whose exchange could dress it.

We now write the equations of motion of all four fields at this order, to show their pattern and the couplings between the fields.
  As at the previous orders, the equation of each auxiliary field is normalized with coefficient one on its source, that is, up to the overall normalization of the term in the action:
  \begin{equation}
  \label{p:eq:hs8:eom-normalization}
  \delta \caction[\le 8]
    = \pair{\delta \Psi}{\mathcal E^{(6)}_\Psi}
      + \sum_{k = 2}^{4} \frac{2 \gs^{2 k - 2}}{(2 k)!} \, \pair[k]{\delta \Sigma_k}{\mathcal E^{(6)}_{\Sigma_k}},
\end{equation}
  which is the normalization of~\eqref{p:eq:hs4:eom}, \eqref{p:eq:hs5:eom}, \eqref{p:eq:hs6:eom} and~\eqref{p:eq:hs7:saddle-3}.
  The equations of motion of $\caction[\le 8]$ are $\mathcal E^{(6)}_X = 0$ for $X = \Psi, \Sigma_2, \Sigma_3, \Sigma_4$, with
  \begin{subequations}
  \label{p:eq:hs8:eom}
  \begin{align}
    \label{p:eq:hs8:eom-psi}
    \mathcal E^{(6)}_\Psi
      &
      := Q_B \Psi
        + \frac{\gs}{2} \, \ell_2(\Psi^2)
        + \frac{\gs^2}{3!} \, \ell_3(\Psi, \Sigma_2)
        + \frac{\gs^3}{5!} \, \ell_4(\Sigma_2^2)
        + \frac{2 \gs^4}{6!} \, \dvertex{6}(\, \cdot \, \otimes \Sigma_2, \Sigma_3)
    \notag
    \\ &
      \quad
        + \frac{\gs^5}{7!} \, \dvertex{7}(\Sigma_3, \Sigma_3, \, \cdot \,)
        + \frac{2 \gs^6}{8!} \, \dvertex{8}(\, \cdot \, \otimes \Sigma_3, \Sigma_4),
    \\
    \label{p:eq:hs8:eom-sigma2}
    \mathcal E^{(6)}_{\Sigma_2}
      &
      := \sprod{2}{2}(\Psi^2)
        - \sprod{2}{2}(\Sigma_2)
        + \frac{\gs}{5} \, \sprod{3}{2}(\Psi, \Sigma_2)
        + \frac{\gs^2}{30} \, \dvertex{6}(\Psi \otimes \, \cdot \,, \Sigma_3),
    \\
    \label{p:eq:hs8:eom-sigma3}
    \mathcal E^{(6)}_{\Sigma_3}
      &
      := \wprod{3}{3}(\Psi \otimes \Sigma_2)
        - \wprod{3}{3}(\Sigma_3)
        + \frac{\gs}{7} \, \dvertex{7}(\, \cdot \,, \Sigma_3, \Psi)
        + \frac{\gs^2}{56} \, \dvertex{8}(\Psi \otimes \, \cdot \,, \Sigma_4),
    \\
    \label{p:eq:hs8:saddle-4}
    \mathcal E^{(6)}_{\Sigma_4}
      &
      := \kin{4}{}{6}(\Psi \otimes \Sigma_3)
        - \kin{4}{}{6}(\Sigma_4),
  \end{align}
\end{subequations}
  where the dot denotes the state paired with the open slot, as in~\eqref{p:eq:hs6:eom}.
  Each term is the variation of one term of the action.
  In the equation of $\Psi$ these are, in order, the kinetic term and the cubic vertex, the quartic source, the quintic coupling, the sextic source, the septic coupling and the octic source.
  In the equation of each auxiliary field they are its source, its kinetic term, its odd coupling to $\Psi$, and the source of the next field, varied with respect to $\Sigma_k$ rather than to that field.
  The equations of the auxiliary fields are linear in the tower and tridiagonal, which is the triangular form of \cref{p:sec:hs:septic} read on the equations: the equation of $\Sigma_k$ involves the lower-order field $\Sigma_{k - 1}$ through its source, $\Sigma_k$ itself through its kinetic term and its odd coupling, and the higher-order field $\Sigma_{k + 1}$ only through one coupling, which comes with $\gs^2$.
  The equation of $\Psi$ is the only non-linear one, and the only one that involves every field.

We solve the equations from the top.
  The last one has the solution
  \begin{equation}
  \label{p:eq:hs8:solution-4}
  \Sigma_4^{(6)*}
    = \Psi \otimes \Sigma_3
\end{equation}
  on the quotient $\Hn{4} / \ker \kin{4}{}{6}$, where the $4$-string field lives for the reason given at~\eqref{p:eq:hs6:quotient}: a kernel direction has zero action and couples to nothing at this order.
  Once substituted in the equation of $\Sigma_3$, the last term becomes the $\gs^2 / 56$ term of the operator~\eqref{p:eq:hs8:dressed-3} below, and the equation of $\Sigma_2$ then receives the operator~\eqref{p:eq:hs8:dressed-2}.
  The elimination from the top is thus the Gaussian elimination of a tridiagonal system, with one inversion per field.
  The solution depends on $\Sigma_3$, and $\Sigma_4 = \Psi^4 + O(\gs)$ once the lower-order fields are on-shell.
  We write the equations of motion at every order, with sums, in~\eqref{p:eq:all:eom} of \cref{p:sec:all-orders}.
Since $\Delta S_8$ vanishes on this solution, eliminating $\Sigma_4$ gives back the octic term exactly in $\gs$ and for any $\Sigma_3$:
\begin{equation}
\caction[\le 8] \big|_{\Sigma_4 = \Sigma_4^{(6)*}}
  = \caction[\le 7]
    + \frac{\gs^6}{8!} \, \dvertex{8}(\Psi \otimes \Sigma_3, \Psi \otimes \Sigma_3).
\end{equation}
The induced action is still quadratic in terms $\Sigma_3$.
At order $\gs^6$ only the diagonal value $\dvertex{8}(\Psi^8)$ of the octic term enters, by stationarity of $\caction[\le 7]$ at its own solution, and~\eqref{p:eq:hs8:diagonal} gives:
\begin{equation}
  \label{p:eq:hs8:eliminate}
  \caction[\le 8] \big|_{\text{solution}}
    = S_{\le 8} + O(\gs^7).
\end{equation}
The octic order closes, and its remainder of order $\gs^7$ is absorbed by the odd vertex $\dvertex{9}$ through the coupling $\dvertex{9}(\Sigma_4, \Sigma_4, \Psi)$, which dresses $\kin{4}{}{6}$ as $\dvertex{7}$ dressed $\wprod{3}{3}$ in~\eqref{p:eq:hs7:dressed-3}.
The pair $(\dvertex{8}, \dvertex{9})$ is for $\Sigma_4$ what $(\dvertex{6}, \dvertex{7})$ was for $\Sigma_3$, and we checked the closure of both orders in the finite-dimensional model, with $\dvertex{8}$ and $\dvertex{9}$ the polarizations of their diagonal values.

\subparagraph{The dressed kinetic operators.}

Once $\Sigma_4$ is eliminated, the quadratic term in $\Sigma_3$ of the octic term adds to~\eqref{p:eq:hs7:dressed-3}.
For $A, B \in \Hn{3}$, the dressed operator of $\Sigma_3$ at this order is:
\begin{equation}
  \label{p:eq:hs8:dressed-3}
  \pair[3]{A}{\kin{3}{\Psi}{6}(B)}
    := \pair[3]{A}{\kin{3}{\Psi}{5}(B)}
      - \frac{\gs^2}{56} \, \pair[4]{\kin{4}{}{6}(\Psi \otimes A)}{\big( \kin{4}{}{6} \big)^{-1} \circ \kin{4}{}{6}(\Psi \otimes B)}.
\end{equation}
The last term is the exchange of $\Sigma_4$ between two octic couplings.
At this order the $\Sigma_4$ line carries the bare operator, and $\dvertex{9}$ dresses it at the next order.
Since $(\kin{4}{}{6})^{-1} \circ \kin{4}{}{6}$ is the identity on the quotient, the term is at this order the contact term $- \tfrac{\gs^2}{56} \dvertex{8}(\Psi \otimes A, \Psi \otimes B)$.
We keep the exchange form because at the next order the line carries the operator dressed by $\dvertex{9}$, and the inverse no longer cancels.
The dressed operator of $\Sigma_2$ is~\eqref{p:eq:hs7:dressed-2} with $\kin{3}{\Psi}{5}$ replaced by $\kin{3}{\Psi}{6}$ inside the $\Sigma_3$ exchange, for $A, B \in \Hn{2}$:
\begin{equation}
  \label{p:eq:hs8:dressed-2}
  \begin{aligned}
    \pair[2]{A}{\kin{2}{\Psi}{6}(B)}
      &
      := \vertex{4}(A, B)
        - \frac{\gs}{5} \, \vertex{5}(A, B, \Psi)
      \\ &
      \quad
      - \frac{\gs^2}{30} \, \pair[3]{\wprod{3}{3}(\Psi \otimes A)}{\big( \kin{3}{\Psi}{6} \big)^{-1} \circ \wprod{3}{3}(\Psi \otimes B)}.
  \end{aligned}
\end{equation}
The three operators are nested: each one contains the inverse of the next one up the tower, and the chain stops at $\kin{4}{}{6}$, which is bare.


Every occurrence of $\Psi$ in these operators comes with one power of $\gs$.
The odd vertex enters as $\gs \, \dvertex{2 k + 1}(\, \cdot \,, \cdot \,, \Psi)$, and the source of the next field, linear in $\Psi$, enters the exchange term twice, with $\gs^2$.
The dressed kinetic operators are therefore functions of $\gs \Psi$ alone.
Their expansion in that variable starts with the bare operator, and the exchange term reduces at this order to the reduced even vertex as in~\eqref{p:eq:hs6:dressed-2},
\begin{equation}
  \label{p:eq:hs8:scaling}
  \begin{aligned}
    \pair[k]{A}{\kin{k}{\Psi}{n}(B)}
      &
      = \pair[k]{A}{\kin{k}{}{2 k - 2}(B)}
        - \frac{\gs}{2 k + 1} \, \dvertex{2 k + 1}(A, B, \Psi)
      \\ &
      \quad
      - \frac{\gs^2}{(2 k + 1) (2 k + 2)} \, \dvertex{2 k + 2}(\Psi \otimes A, \Psi \otimes B)
        + O(\gs^3 \Psi^3),
  \end{aligned}
\end{equation}
for $k = 2, 3$ at this order and $n \ge 2 k$, with $\dvertex{5} := \vertex{5}$ and $\kin{2}{}{2} := \sprod{2}{2}$.

The dependence through $\gs \Psi$ alone holds at every order.
The reduction to the bare operator at $\gs \Psi = 0$ holds as soon as the bare operators of the higher-order fields are invertible on their quotients at generic momentum and the condition~\eqref{p:eq:hs8:kernel-condition} below holds: the inverse of $\kin{k + 1}{\Psi}{n}$ is then analytic at $\gs \Psi = 0$.
The operator introduced with the square at order $\gs^{2 k - 2}$ is thus the value at $\gs = 0$ of the dressed operator of $\Sigma_k$ at every later order, and the background enters through $\gs \Psi$ only.
This also turns the invertibility assumption on the dressed operators into a condition on the bare ones.
At a fixed level truncation, with the momenta fixed or discretized, and under the same condition, a dressed operator whose bare value is invertible on the quotient stays so for $\gs \Psi$ in a neighborhood of zero.

\subparagraph{The condition on the nested kernels.}

The auxiliary fields live on quotients, $\Sigma_k$ on $\Hn{k} / \ker \kin{k}{}{2 k - 2}$, and the action must be a function on them.
At the quintic and septic orders this gave a condition on the odd vertex, \eqref{p:eq:hs5:kernel-condition} and~\eqref{p:eq:hs7:kernel-condition}.
At the sextic order it gave~\eqref{p:eq:hs6:kernel-condition}, a condition on the source of the next field.
Besides its kinetic term, a kernel direction $N_k \in \ker \kin{k}{}{2 k - 2}$ appears in three terms of the action: the odd coupling, the source of $\Sigma_{k + 1}$, and its own source $\pair[k]{\Sigma_k}{\kin{k}{}{2 k - 2}(\Psi \otimes \Sigma_{k - 1})}$.
The last one pairs to zero with $N_k$ by the symmetry of $\kin{k}{}{2 k - 2}$, and the action descends to the quotients if the other two vanish on every kernel direction:
\begin{equation}
  \label{p:eq:hs8:kernel-condition}
  \begin{aligned}
    &
    \forall N_k \in \ker \kin{k}{}{2 k - 2}, \;
    \forall A \in \Hn{k}, \;
    \forall B \in \Hs:
    \\ &
    \qquad
    \dvertex{2 k + 1}(N_k, A, B)
      = 0,
    \qquad
    \kin{k + 1}{}{2 k}(B \otimes N_k)
      = 0,
  \end{aligned}
\end{equation}
for every $k \ge 2$, with the conventions of~\eqref{p:eq:hs8:scaling} at $k = 2$.
The first clause is the condition of the odd orders, and the second that of the even ones.
At $k = 2$, the second clause is~\eqref{p:eq:hs6:kernel-condition}, since the exchange part of $\dvertex{6}$ vanishes on the kernel by the first clause.
The second clause says that a null direction of one level, tensored with any string field, is a null direction of the next, to which the condition applies in its turn: null directions propagate upward through the whole tower.
Equivalently, every null vector of the kinetic operator of $\Sigma_k$ must be a joint null vector of every vertex in which $\Sigma_k$ appears.
Here, as for $\sprod{2}{2}$ and for the kinetic operator of $\Sigma_3$ in \cref{p:sec:hs:sextic}, a kernel is meant in the square-integrable sense: at fixed total momentum, its elements are square-integrable functions of the relative momenta.
Through the second clause, a null vector of $\sprod{2}{2}$ at a single total momentum gives every higher operator null vectors at every total momentum, but these are plane waves, sharp in one relative momentum, and are not counted in their kernels, so that the condition does not by itself exclude zeros of $\sprod{2}{2}$, even when the higher operators have no kernel.

This is the standing assumption of \cref{p:sec:hs:septic}, extended to the whole tower.
At $k = 2$ it is a condition on the quintic and sextic vertices themselves.
For $k \ge 3$ it is in part a condition on the choice of the reduced vertices, since each is defined up to terms vanishing on the diagonal.
  A choice satisfying it exists, however, only if the diagonal values fixed by the matching vanish on the kernel directions.
The condition holds trivially when the kinetic operators have no kernel.

The kernels in~\eqref{p:eq:hs8:kernel-condition} are those of the bare operators, while the operators that are inverted depend on the background.
The quotient prescription is meaningful only if $\ker \kin{k}{}{2 k - 2}$ stays in the radical of $\kin{k}{\Psi}{n}$ for every $\Psi$ and every small $\gs$, so that every dressed operator descends to the same quotient.
Read order by order in~\eqref{p:eq:hs8:scaling}, this is condition~\eqref{p:eq:hs8:kernel-condition}: the first clause makes the term of order $\gs$ vanish on the kernel, and the second the exchange term.
Conversely, the second clause is required as soon as $\kin{k + 1}{}{2 k}$ is definite on its quotient.
When the condition fails, a kernel direction of the bare operator is lifted, at order $\gs$ when the first clause fails and at order $\gs^2$ when the second does.
The inverse of $\kin{k}{\Psi}{n}$ then has a pole at $\gs = 0$, and the quotient is not defined for $\gs \neq 0$.
A second requirement, independent of the first, is that the quotients and the inverses have a limit as the mass level of the truncation grows, at generic momentum.
We assume it for every $\kin{k}{}{2 k - 2}$, in the weak form we assumed for the quartic form: the kernels at successive levels need not be nested, and only the quotients and the inverses are required to converge.

In the construction used here, a failure of the condition at one order of the tower does not show in the dressed operator of the field below it at $\gs = 0$, for poles of order at most two, as we now check.
Suppose that a single null vector $N$ of $\kin{k + 1}{}{2 k}$, the operators of the fields above $\Sigma_{k + 1}$ being regular, violates~\eqref{p:eq:hs8:kernel-condition} and is not quotiented.
Restricted to the line of $N$ after eliminating the complement, $\kin{k + 1}{\Psi}{n}$ is $\gs \alpha_1 + \gs^2 \alpha_2 + O(\gs^3)$, with $\alpha_1 = - \tfrac{1}{2 k + 3} \dvertex{2 k + 3}(N, N, \Psi)$.
When $\alpha_1 = 0$, $\alpha_2$ is the order-$\gs^2$ entry along $N$ corrected by the order-$\gs$ mixing of $N$ with the complement, a quadratic form in $\Psi$.
The inverse of the restricted operator has a pole at $\gs = 0$, simple if $\alpha_1 \neq 0$, double if $\alpha_1$ vanishes identically and $\alpha_2$ does not.

To see the effect of the pole, we consider a variant of the construction in which the kernel direction couples to the source of $\Sigma_{k + 1}$, through a form $\ell$ on $\Hn{k}$ linear in $\Psi$.
The exchange term in the operator of $\Sigma_k$ then picks up the pole.
For a double pole the value at $\gs = 0$ of the dressed operator of $\Sigma_k$ is then shifted from the bare one by $- \ell(A) \, \ell(B) / ((2 k + 1) (2 k + 2) \, \alpha_2)$.
For a simple pole the same term, with $\alpha_1$ in place of $\alpha_2$, enters at order $\gs$.
The shift is homogeneous of degree zero in $\Psi$, so that the operator has a direction-dependent limit at $\gs \Psi = 0$, which the reduction to the bare operator misses.
In the construction used here the source of $\Sigma_{k + 1}$ is $\kin{k + 1}{}{2 k}(\Psi \otimes \Sigma_k)$, and $\ell$ vanishes by the symmetry of $\kin{k + 1}{}{2 k}$: the value at $\gs = 0$ is protected for poles of order at most two.
A failed order of the tower is then invisible in the dressed operators at $\gs = 0$.
Hence the invertibility of the quartic form on its quotient never certifies the fields above $\Sigma_2$, and the failure shows up in the eliminated action.

\subparagraph{Failure of the kernel condition.}

Let a kernel direction $N$ of $\kin{k}{}{2 k - 2}$ violate the first clause of~\eqref{p:eq:hs8:kernel-condition} and satisfy the second, the kinetic operators of the lower-order fields being invertible, with no kernel.
That direction has no kinetic term at order $\gs^{2 k - 2}$ and acquires a quadratic term from the background at the next order, $\gs \, \dvertex{2 k + 1}(N, N, \Psi)$ up to normalization.
Integrating it out divides by this field-dependent mass, and the eliminated action acquires a term rational in $\Psi$,
\begin{equation}
  \label{p:eq:hs8:rational}
  \Delta S \big|_{\gs^{2 k - 1}}
    = - \frac{\gs^{2 k - 1}}{(2 k + 1)!} \, \frac{\dvertex{2 k + 1}(N, \Psi^k, \Psi)^2}{\dvertex{2 k + 1}(N, N, \Psi)},
\end{equation}
when the denominator does not vanish identically.
When it does but $\dvertex{2 k + 1}(N, \, \cdot \,, \Psi)$ does not, the pole is double and the term appears one order earlier, at $\gs^{2 k - 2}$, with a quadratic form in $\Psi$ in the denominator.
The case $k = 3$ is the term of \cref{p:sec:hs:septic}, which we checked in the finite-dimensional model.
The general $k$ is argued by the same computation, the shift of the dressed operator of $\Sigma_{k - 1}$ entering the value through its diagonal, and we have not checked it beyond $k = 3$.
The ratio has the degree $2 k + 1$ of its order but is not a polynomial for generic vertices, and no vertex, elementary or reduced, can absorb it: the matching~\eqref{p:eq:hs7:deficit} adjusts a polynomial diagonal.
The condition~\eqref{p:eq:hs8:kernel-condition} is thus what makes the auxiliary tower reproduce a polynomial action.

When the condition holds and the bare operators are invertible on their quotients at generic momentum, the elimination goes through at every truncation.
The value of the action is then a continued fraction in the nested operators, which \cref{p:sec:all-orders} states for the whole tower.
When the condition fails for the actual vertices, the auxiliary field concerned cannot be taken on a quotient, and we use instead the multiplier form~\eqref{p:eq:all:multiplier-action} of the action.
Its fields are constrained to be $\Psi^k$ by multipliers, including their kernel components, and nothing is inverted, so that no kernel condition arises, but the exchange structure is lost.
The hybrid of \cref{p:sec:hs:sextic:kernel}, which constrains the kernel directions only, keeps the exchange structure at the sextic order, where its reduced vertex depends on the complement; we have not carried it beyond that order.
Which of the two forms applies to the string vertices depends on the kernels of the operators $\kin{k}{}{2 k - 2}$.
We do not compute them here, but level-truncated data settle them mass level by mass level.


\subsection{The quintic order in the companion note}
\label{sec:order:companion-quintic}

The next three blocks are sections of a companion note on the quintic order.
The first completes the square with the dressed kernel line by line and shows that the saddle is the resummation of the recursion of the paper's quintic order; the second explains why the correction of order $\gs$ to the saddle does not enter the action at order $\gs^3$, and that any other coefficient of the quintic coupling fails; the third shows that the kernel condition of the quintic order is necessary, not only sufficient.
The note writes the coupling as $g_s$ and is otherwise close to the paper's notation; the differences that occur are listed below.

\begin{quote}
  \small
  \textbf{Translation box: companion note on the quintic order $\to$ paper.}
  \begin{itemize}
    \setlength{\itemsep}{0.3ex}
    \item $g_s$ $\to$ $\gs$; $\mc V_n$ is the paper's $\vertex{n}$; $p_{2,2}$, $p_{3,2}$ are the paper's $\sprod{2}{2}$, $\sprod{3}{2}$.
    \item $\Sym^2 \mc H_0$, the symmetric level-matched $2$-string states, $\to$ $\Hn{2}$ (input strands of ghost number $2$, symmetric components; see \cref{ch:conventions}); $\braket{A, B}$ $\to$ $\pair[2]{A}{B}$.
    \item The dressed kernel $K_\Psi(A, B) = \braket{A, k_\Psi(B)}$, $k_\Psi = p_{2,2} - \tfrac{g_s}{5} \, p_{3,2}(\Psi, \cdot)$ $\to$ the dressed kinetic operator $\kin{2}{\Psi}{3}$ of~\eqref{p:eq:hs5:dressed}.
    \item $\widetilde S_{\le 5}$, $S_{\le 5}$ $\to$ $\caction[\le 5]$, $S_{\le 5}$.
    \item The deviation $X = \Sigma_2 - \Psi^2$, the shifted saddle $\Sigma_2^*$ and the shifted deviation $\tilde X = \Sigma_2 - \Sigma_2^*$ $\to$ $\sigma_2^{(2)}$, $\Sigma_2^{(3)*}$ and $\sigma_2^{(3)}$ of~\eqref{p:eq:hs5:deviation}.
    \item The source form $v(A) = \mc V_5(A, \Psi^2, \Psi) = \braket{A, p_{3,2}(\Psi^3)}$; the feedback $\tfrac{g_s^4}{600} \, v(K_\Psi^{-1} v)$ is, at its leading order, the paper's $\tfrac{\gs^4}{600} \vertex{5 \cdot 5}(\Psi^6)$ of~\eqref{p:eq:hs5:eliminate}.
    \item $\Xi$, $N$ denote kernel elements of $\mc V_4$, the paper's $N_2 \in \ker \sprod{2}{2}$.
  \end{itemize}
\end{quote}

\subsection{Completing the square with the dressed kernel}
\label{q:sec:square}

The whole content of the quintic step is one exact identity: $\widetilde S_{\le5}$ equals $S_{\le5}$ plus a quadratic form in a shifted deviation, which vanishes on shell, plus an explicit $\Sigma_2$-independent term of order $g_s^4$.
This section derives it line by line, then reads off the saddle and shows that it is the resummation of the recursion \eqref{p:eq:hs5:recursion}.
Throughout, $\mc V_4$ is invertible on $\Sym^2\mc H_0/\ker\mc V_4$, and $\mc V_5$ is assumed compatible with the quotient in the sense of the kernel condition \eqref{p:eq:hs5:kernel-condition}; the inverse of the dressed kernel below then exists for small $g_s$ as a Neumann series, and the note assumes it exists at the level of the conformal field theory.

\subsubsection{Collecting the \texorpdfstring{$\Sigma_2$}{Sigma2} terms}
\label{q:sec:square:collect}

The terms of $\caction[\le 4]$ plus \eqref{p:eq:hs5:S5tilde} that contain $\Sigma_2$ are
\begin{equation}
  \label{q:eq:square:F}
  F(\Sigma_2)
    := - \frac{g_s^2}{4!} \, \mc V_4(\Sigma_2, \Sigma_2)
      + \frac{2 g_s^2}{4!} \, \mc V_4(\Sigma_2, \Psi^2)
      + \frac{g_s^3}{5!} \, \mc V_5(\Sigma_2, \Sigma_2, \Psi);
\end{equation}
the kinetic and cubic terms of $\Psi$ are common to $\widetilde S_{\le5}$ and $S_{\le5}$ and drop out of every difference below.
$F$ is a quadratic polynomial in $\Sigma_2$: a quadratic form (the two terms with two $\Sigma_2$) and a linear form (the source).

\subsubsection{Absorbing the coupling into the kernel}
\label{q:sec:square:kernel}

The two quadratic terms combine into a single bilinear form once the coefficients are aligned.
Since $4!/5!=1/5$,
\begin{equation}
  \frac{g_s^3}{5!}
    = \frac{g_s^2}{4!} \cdot \frac{g_s}{5},
\end{equation}
so that
\begin{equation}
  \label{q:eq:square:absorb}
  - \frac{g_s^2}{4!} \, \mc V_4(\Sigma_2, \Sigma_2)
    + \frac{g_s^3}{5!} \, \mc V_5(\Sigma_2, \Sigma_2, \Psi)
    = - \frac{g_s^2}{4!} \, K_\Psi(\Sigma_2, \Sigma_2),
\end{equation}
where
\begin{equation}
  \label{q:eq:square:K}
  \boxed{\;
  \begin{aligned}
    K_\Psi(A, B)
      &
      := \mc V_4(A, B) - \frac{g_s}{5} \, \mc V_5(A, B, \Psi)
      = \braket{A, k_\Psi(B)},
    \\
    k_\Psi
      &
      := p_{2,2} - \frac{g_s}{5} \, p_{3,2}(\Psi, \cdot).
  \end{aligned}
  \;}
\end{equation}
We call $K_\Psi$ the \emph{dressed kernel}: it is the kinetic form of $\Sigma_2$ with the quintic coupling absorbed, so that $\Sigma_2$ propagates in the background of $\Psi$.
It is a small perturbation of $\mc V_4$, symmetric because $\mc V_4$ and $\mc V_5$ are, and linear in $\Psi$ at order $g_s$.
Its inverse on the quotient is
\begin{equation}
  \label{q:eq:square:neumann}
  k_\Psi^{-1}
    = \sum_{n \ge 0} \Big( \frac{g_s}{5} \, p_{2,2}^{-1} \, p_{3,2}(\Psi, \cdot) \Big)^n \, p_{2,2}^{-1},
\end{equation}
a Neumann series which converges for small $g_s$ whenever $p_{2,2}^{-1}$ is bounded: for the order-by-order statements no new invertibility condition arises at this order beyond the one on $p_{2,2}$, while the statements at finite $g_s$ below assume that $K_\Psi$ itself is invertible on the quotient.

\subsubsection{The source through the dressed kernel}
\label{q:sec:square:source}

The source term of \eqref{q:eq:square:F} is still built from the bare $\mc V_4$.
To complete the square with $K_\Psi$ we rewrite it through $K_\Psi$, using \eqref{q:eq:square:K} with $B=\Psi^2$:
\begin{equation}
  \label{q:eq:square:source}
  \mc V_4(\Sigma_2, \Psi^2)
    = K_\Psi(\Sigma_2, \Psi^2) + \frac{g_s}{5} \, \mc V_5(\Sigma_2, \Psi^2, \Psi)
    = K_\Psi(\Sigma_2, \Psi^2) + \frac{g_s}{5} \, v(\Sigma_2),
\end{equation}
where
\begin{equation}
  \label{q:eq:square:v}
  v(A)
    := \mc V_5(A, \Psi^2, \Psi)
    = \braket{A, p_{3,2}(\Psi^3)}
\end{equation}
is a linear form on $\Sym^2\mc H_0$, cubic in $\Psi$ (two powers through $\Psi^2$, one explicit).
Inserting \eqref{q:eq:square:absorb} and \eqref{q:eq:square:source} into \eqref{q:eq:square:F},
\begin{equation}
  \label{q:eq:square:F-dressed}
  F(\Sigma_2)
    = - \frac{g_s^2}{4!} \Big[
        K_\Psi(\Sigma_2, \Sigma_2)
        - 2 K_\Psi(\Sigma_2, \Psi^2)
        - \frac{2 g_s}{5} \, v(\Sigma_2)
      \Big].
\end{equation}

\subsubsection{The shifted saddle and the square}
\label{q:sec:square:saddle}

The two linear terms in the bracket of \eqref{q:eq:square:F-dressed} are both of the form $2K_\Psi(\Sigma_2,\cdot)$, since $v(\Sigma_2)=K_\Psi(\Sigma_2,K_\Psi^{-1}v)$ by definition of $K_\Psi^{-1}v$:
\begin{equation}
  2 K_\Psi(\Sigma_2, \Psi^2) + \frac{2 g_s}{5} \, v(\Sigma_2)
    = 2 K_\Psi(\Sigma_2, \Sigma_2^*),
\end{equation}
where the shifted saddle is
\begin{equation}
  \label{q:eq:square:saddle}
  \boxed{\;
  \Sigma_2^*
    := \Psi^2 + \frac{g_s}{5} \, K_\Psi^{-1} v
    = \Psi^2 + \frac{g_s}{5} \, k_\Psi^{-1} \, p_{3,2}(\Psi^3).
  \;}
\end{equation}
Let's complete the square.
With the shifted deviation
\begin{equation}
  \label{q:eq:square:Xtilde}
  \tilde X
    := \Sigma_2 - \Sigma_2^*
    = X - \frac{g_s}{5} \, K_\Psi^{-1} v,
\end{equation}
the symmetry of $K_\Psi$ gives
\begin{equation}
  \label{q:eq:square:square}
  K_\Psi(\Sigma_2, \Sigma_2) - 2 K_\Psi(\Sigma_2, \Sigma_2^*)
    = K_\Psi(\tilde X, \tilde X) - K_\Psi(\Sigma_2^*, \Sigma_2^*).
\end{equation}
The variation of $F$ is now $\delta F=-\frac{2g_s^2}{4!}K_\Psi(\delta\Sigma_2,\tilde X)$, so the equation of motion of $\Sigma_2$ is exactly
\begin{equation}
  \tilde X = 0,
  \qquad\text{i.e.}\qquad
  \Sigma_2 = \Sigma_2^* .
\end{equation}

\subsubsection{Reassembling the effective vertices}
\label{q:sec:square:reassemble}

The constant term $K_\Psi(\Sigma_2^*,\Sigma_2^*)$ contains the target.
Expanding \eqref{q:eq:square:saddle} and using $K_\Psi(\Psi^2,K_\Psi^{-1}v)=v(\Psi^2)$ and $K_\Psi(K_\Psi^{-1}v,K_\Psi^{-1}v)=v(K_\Psi^{-1}v)$,
\begin{equation}
  \label{q:eq:square:const}
  K_\Psi(\Sigma_2^*, \Sigma_2^*)
    = K_\Psi(\Psi^2, \Psi^2)
      + \frac{2 g_s}{5} \, v(\Psi^2)
      + \frac{g_s^2}{25} \, v(K_\Psi^{-1} v).
\end{equation}
The first two terms reassemble the quartic and quintic vertices with their coefficients.
From \eqref{q:eq:square:K} and \eqref{q:eq:square:v}, $K_\Psi(\Psi^2,\Psi^2)=\mc V_4(\Psi^4)-\frac{g_s}{5}\mc V_5(\Psi^5)$ and $v(\Psi^2)=\mc V_5(\Psi^5)$, hence
\begin{equation}
  K_\Psi(\Psi^2, \Psi^2) + \frac{2 g_s}{5} \, v(\Psi^2)
    = \mc V_4(\Psi^4) + \frac{g_s}{5} \, \mc V_5(\Psi^5),
\end{equation}
and multiplying by the prefactor $\frac{g_s^2}{4!}$,
\begin{equation}
  \frac{g_s^2}{4!} \, \mc V_4(\Psi^4) + \frac{g_s^2}{4!} \cdot \frac{g_s}{5} \, \mc V_5(\Psi^5)
    = \frac{g_s^2}{4!} \, \mc V_4(\Psi^4) + \frac{g_s^3}{5!} \, \mc V_5(\Psi^5),
\end{equation}
which are exactly the quartic and quintic terms of $S_{\le5}$.
This is how the normalization $1/5!$ of \eqref{p:eq:hs5:S5tilde} is forced: the quintic coefficient of the target is produced by the cross term $\frac{2g_s}{5}v(\Psi^2)$ of the square minus the $-\frac{g_s}{5}\mc V_5(\Psi^5)$ inside $K_\Psi(\Psi^2,\Psi^2)$, and the two combine to $+\frac{g_s}{5}$ only if the coupling carries $1/5!$.
The last term of \eqref{q:eq:square:const} carries the prefactor
\begin{equation}
  \frac{g_s^2}{4!} \cdot \frac{g_s^2}{25}
    = \frac{g_s^4}{600},
  \qquad
  600 = 4! \cdot 25,
\end{equation}
the square of the shift $\frac{g_s}{5}$ against the $\frac{g_s^2}{4!}$ of the kinetic term.

\subsubsection{The identity}
\label{q:sec:square:identity}

Putting \eqref{q:eq:square:F-dressed}, \eqref{q:eq:square:square} and \eqref{q:eq:square:const} together, and adding back the common $\Psi$ terms,
\begin{equation}
  \label{q:eq:square:identity}
  \boxed{\;
  \widetilde S_{\le 5}
    = S_{\le 5}
      - \frac{g_s^2}{4!} \, K_\Psi(\tilde X, \tilde X)
      + \frac{g_s^4}{600} \, v(K_\Psi^{-1} v).
  \;}
\end{equation}
This is an identity in the fields $(\Psi,\Sigma_2)$, valid at every order in $g_s$; it has been checked symbolically as such.
The three terms are the target, a quadratic form in the shifted deviation, and a term independent of $\Sigma_2$.
We call the last one the \emph{feedback}: it is the surplus that the Gaussian elimination of $\Sigma_2$ produces beyond the vertex it was introduced to reproduce.
It is of order $g_s^4$ and of degree six in $\Psi$ ($v$ is cubic and appears twice), so it is a six-point interaction; it does not affect the matching at order $g_s^3$, and it is the exact form of the statement that the auxiliary field feeds back into the action only one order later.

On shell, $\tilde X=0$ and
\begin{equation}
  \label{q:eq:square:onshell}
  \widetilde S_{\le 5}\big(\Psi, \Sigma_2^*\big)
    = S_{\le 5}(\Psi) + \frac{g_s^4}{600} \, v(K_\Psi^{-1} v),
\end{equation}
which is \eqref{p:eq:hs5:eliminate} with the remainder made explicit.
All tree amplitudes of $\widetilde S_{\le5}$ therefore agree with those of $S_{\le5}$ through order $g_s^3$; \cref{p:fig:hs5:five-point} shows the mechanism diagram by diagram.

\subsubsection{The saddle and the recursion}
\label{q:sec:square:recursion}

The saddle \eqref{q:eq:square:saddle} is the resummed form of the recursive solution \eqref{p:eq:hs5:recursion}.
The equation of motion \eqref{p:eq:hs5:eom-sigma}, up to an overall sign $0=p_{2,2}(\Sigma_2)-p_{2,2}(\Psi^2)-\frac{g_s}{5}p_{3,2}(\Psi,\Sigma_2)$, is $k_\Psi(\Sigma_2)=p_{2,2}(\Psi^2)$ in the notation \eqref{q:eq:square:K}, so that
\begin{equation}
  \label{q:eq:square:saddle-R}
  \Sigma_2^*
    = k_\Psi^{-1} \, p_{2,2}(\Psi^2).
\end{equation}
To see that this agrees with \eqref{q:eq:square:saddle}, apply \eqref{q:eq:square:K} to $\Psi^2$: $p_{2,2}(\Psi^2)=k_\Psi(\Psi^2)+\frac{g_s}{5}p_{3,2}(\Psi,\Psi^2)$, and $p_{3,2}(\Psi,\Psi^2)=p_{3,2}(\Psi^3)$ by symmetry; applying $k_\Psi^{-1}$ gives $\Sigma_2^*=\Psi^2+\frac{g_s}{5}k_\Psi^{-1}p_{3,2}(\Psi^3)$.
Expanding $k_\Psi^{-1}$ with the Neumann series \eqref{q:eq:square:neumann},
\begin{equation}
  \label{q:eq:square:recursion}
  \Sigma_2^*
    = \Psi^2
      + \frac{g_s}{5} \, p_{2,2}^{-1} \, p_{3,2}(\Psi^3)
      + \frac{g_s^2}{25} \, p_{2,2}^{-1} \, p_{3,2}\Big( \Psi, p_{2,2}^{-1} \, p_{3,2}(\Psi^3) \Big)
      + O(g_s^3),
\end{equation}
which is the recursion \eqref{p:eq:hs5:recursion} iterated once: each iteration inserts one more $p_{3,2}(\Psi,\cdot)$ between two $p_{2,2}^{-1}$, with one more factor $\frac{g_s}{5}$.
The dressed-kernel inverse is thus the closed form of the recursion: what looks like interactions at all orders generated by $\Sigma_2$ is the geometric series of the self-energy insertions $p_{3,2}(\Psi,\cdot)$ on the $\Sigma_2$ line.
The same series, read in the action rather than in the saddle, expands the feedback as
\begin{equation}
  \frac{g_s^4}{600} \, v(K_\Psi^{-1} v)
    = \frac{g_s^4}{600} \Big[
      v(\mc V_4^{-1} v)
      + \frac{g_s}{5} \, v\big(\mc V_4^{-1} W_\Psi \mc V_4^{-1} v\big)
      + O(g_s^2)
    \Big],
\end{equation}
with $W_\Psi:=\mc V_5(\cdot,\cdot,\Psi)$, where $\mc V_4^{-1}W_\Psi\mc V_4^{-1}v$ is shorthand for $p_{2,2}^{-1}p_{3,2}(\Psi,p_{2,2}^{-1}p_{3,2}(\Psi^3))$.



The second block of the note isolates the reason why the correction to the saddle costs no order: stationarity, not the size of the correction.

\subsection{Orders: why the corrected saddle does not spoil the matching}
\label{q:sec:orders}

One may ask why the order-$g_s$ correction to $\Sigma_2^*$ does not enter the action at order $g_s^3$: the kinetic term is quadratic in $\Sigma_2$ and of order $g_s^2$, so a correction of order $g_s$ costs $g_s^4$ there, but the source term is linear in $\Sigma_2$ and also of order $g_s^2$, so the same correction seems to cost only $g_s^3$.
The objection is right about the source term taken alone and wrong about the sum, and the reason is stationarity.
We give the one-line computation, then the general statement.

\subsubsection{The one-line computation}
\label{q:sec:orders:oneline}

Write the saddle as $\Sigma_2^*=\Psi^2+\delta$ with $\delta=\frac{g_s}{5}K_\Psi^{-1}v=O(g_s)$, and insert it into \eqref{q:eq:square:F}.
The kinetic term gives
\begin{equation}
  - \frac{g_s^2}{4!} \, \mc V_4(\Psi^2 + \delta, \Psi^2 + \delta)
    = - \frac{g_s^2}{4!} \, \mc V_4(\Psi^4)
      - \frac{2 g_s^2}{4!} \, \mc V_4(\Psi^2, \delta)
      - \frac{g_s^2}{4!} \, \mc V_4(\delta, \delta),
\end{equation}
and the source term gives
\begin{equation}
  \frac{2 g_s^2}{4!} \, \mc V_4(\Psi^2 + \delta, \Psi^2)
    = \frac{2 g_s^2}{4!} \, \mc V_4(\Psi^4)
      + \frac{2 g_s^2}{4!} \, \mc V_4(\delta, \Psi^2).
\end{equation}
The two terms linear in $\delta$ are equal and opposite, by the symmetry of $\mc V_4$:
\begin{equation}
  \label{q:eq:orders:cancel}
  - \frac{2 g_s^2}{4!} \, \mc V_4(\Psi^2, \delta)
    + \frac{2 g_s^2}{4!} \, \mc V_4(\delta, \Psi^2)
    = 0 .
\end{equation}
This cancellation is the statement that $\Psi^2$ is the saddle of the order-$g_s^2$ action: the first-order variation of $\widetilde S_{\le4}$ in $\Sigma_2$ vanishes at $\Sigma_2=\Psi^2$, whatever the direction $\delta$.
The source term alone does shift the action at order $g_s^3$; the kinetic term shifts it back by the same amount.
What survives is the term quadratic in $\delta$, $-\frac{g_s^2}{4!}\mc V_4(\delta,\delta)=O(g_s^2)\cdot O(g_s)^2=O(g_s^4)$, and the variation of the quintic term, $\frac{2g_s^3}{5!}\mc V_5(\delta,\Psi^2,\Psi)=O(g_s^3)\cdot O(g_s)=O(g_s^4)$.
Both are of order $g_s^4$.

\subsubsection{The general statement}
\label{q:sec:orders:general}

The same argument in one sentence: $F$ is quadratic in $\Sigma_2$, so its value at $\Psi^2+\delta$ is exactly its value at $\Psi^2$, plus its first variation at $\Psi^2$ in the direction $\delta$, plus one half of its second variation; the first variation of the order-$g_s^2$ part vanishes because $\Psi^2$ is its stationary point, the first variation of the quintic part is $O(g_s^3)\times O(g_s)$, and the second variation is the quadratic form $-\frac{g_s^2}{4!}K_\Psi(\delta,\delta)$, of order $g_s^2$ in a shift of order $g_s$.
Hence a shift of the saddle of order $g_s$ changes the on-shell value at order $g_s^4$.
The exact identity \eqref{q:eq:square:identity} says the same without any expansion: on shell the quadratic form vanishes identically and the only remainder is the feedback, which is $O(g_s^4)$ by inspection.
The naive counting, that a correction of order $g_s$ times a coupling of order $g_s^2$ gives $g_s^4$, lands on the right order, but the mechanism is stationarity and not the order of the correction: a correction of order $g_s$ inserted into a term linear in $\Sigma_2$ of order $g_s^2$ would give $g_s^3$ if that term stood alone.

\subsubsection{Any other coefficient fails}
\label{q:sec:orders:falsify}

Repeat \cref{q:sec:square:collect} with a general coefficient $c\,\frac{g_s^3}{5!}$ on the quintic coupling and expand around $\Psi^2$, $\Sigma_2=\Psi^2+X$.
The $\mc V_4$ cross terms cancel as in \eqref{q:eq:orders:cancel}, and
\begin{equation}
  \label{q:eq:orders:c}
  \widetilde S_{\le 5} - S_{\le 5}
    = - \frac{g_s^2}{4!} \, \mc V_4(X, X)
      + \frac{c \, g_s^3}{5!} \Big[ \mc V_5(X, X, \Psi) + 2 \mc V_5(X, \Psi^2, \Psi) \Big]
      + \frac{(c - 1) \, g_s^3}{5!} \, \mc V_5(\Psi^5).
\end{equation}
The last term is independent of $X$ and survives on shell at order $g_s^3$.
Hence closure forces
\begin{equation}
  \boxed{\;
  c = 1,
  \;}
\end{equation}
and any other value leaves the uncancelled mismatch $(c-1)\frac{g_s^3}{5!}\mc V_5(\Psi^5)$ in the effective action, a five-point contact term at order $g_s^3$; this has been checked symbolically by extracting the order-$g_s^3$ coefficient of the on-shell mismatch for $c\neq1$.
The naive rule, replace $\Psi^2$ by $\Sigma_2$ and keep the $1/5!$, is therefore derived and not assumed, and it is the $X$-dependence in \eqref{q:eq:orders:c} that carries all the dynamics.


The third block, from the note's section on the quotient, shows that the kernel condition~\eqref{p:eq:hs5:kernel-condition} is forced by the matching for every choice of representatives, and that keeping the kernel directions with the quadratic form the quintic coupling gives them is worse.

\subsubsection{Necessary as well as sufficient}
\label{q:sec:quotient:necessary}

Condition \eqref{p:eq:hs5:kernel-condition} might look stronger than what the matching needs, which only ever evaluates $\mc V_5$ on $\Psi^2$.
For a totally symmetric $\mc V_5$ it is not: the diagonal condition
\begin{equation}
  \label{q:eq:quotient:diagonal}
  \mc V_5(\Xi, \Psi^2, \Psi) = 0
  \quad\text{for all } \Xi \in \ker \mc V_4,\ \Psi \in \mc H_0
\end{equation}
is equivalent to it, by polarization.
Fix $\Xi\in N$ and let
\begin{equation}
  T_\Xi(\Psi_1, \Psi_2, \Psi_3)
    := \mc V_5(\Xi, \Psi_1, \Psi_2, \Psi_3),
\end{equation}
a symmetric trilinear form on $\mc H_0$ since $\mc V_5$ is symmetric.
Condition \eqref{q:eq:quotient:diagonal} says that its diagonal $T_\Xi(\Psi,\Psi,\Psi)$ vanishes identically, and a symmetric multilinear form is determined by its diagonal,
\begin{equation}
  T_\Xi(\Psi_1, \Psi_2, \Psi_3)
    = \frac{1}{3!} \sum_{S \subseteq \{1, 2, 3\}} (-1)^{3 - |S|} \, T_\Xi(\Psi_S, \Psi_S, \Psi_S),
  \qquad
  \Psi_S := \sum_{i \in S} \Psi_i ,
\end{equation}
so $T_\Xi\equiv0$.
Since $\Sym^2\mc H_0$ is spanned by the symmetrized products $\Psi_1\otimes\Psi_2+\Psi_2\otimes\Psi_1$, this is $\mc V_5(\Xi,A,\Psi)=0$ for all $A$ and $\Psi$, which is \eqref{p:eq:hs5:kernel-condition}.
The symmetry of $\mc V_5$ is what makes the step work; for a non-symmetric five-form the two conditions differ (this is the falsifier of the step in the symbolic check).

Necessity for the matching now follows.
If \eqref{p:eq:hs5:kernel-condition} fails, the quotient is not a theory and one must pick representatives, that is, a complement $C$ and $\Sigma_2\in C$.
On $C$ the action has exactly the form of \cref{q:sec:square} with $(\mc V_4,\mc V_5(\cdot,\cdot,\Psi),\Psi^2)$ replaced by their restrictions and projections $(\mc V_4|_C,\mc V_5(\cdot,\cdot,\Psi)|_C,\pi_C\Psi^2)$, and $\mc V_4|_C$ is non-degenerate for every complement.
Completing the square as before, the on-shell value is $\frac{g_s^2}{4!}\mc V_4(\Psi^4)+\frac{g_s^3}{5!}\mc V_5(\pi_C\Psi^2,\pi_C\Psi^2,\Psi)+O(g_s^4)$ (the quartic term is unchanged because $\mc V_4$ does not see $\pi_N\Psi^2$), so the mismatch is
\begin{equation}
  \label{q:eq:quotient:mismatch}
  \boxed{\;
  \Delta_C
    := \widetilde S_{\le 5}\big|_{C, \text{on shell}} - S_{\le 5}
    = - \frac{g_s^3}{5!} \Big[ \mc V_5(\Psi^2, \Psi^2, \Psi) - \mc V_5(\pi_C \Psi^2, \pi_C \Psi^2, \Psi) \Big] + O(g_s^4),
  \;}
\end{equation}
a five-point term at order $g_s^3$, a genuine contamination of the five-point tree and not a higher-point spill-over, and it depends on the complement through $\pi_N\Psi^2$.
Requiring $\Delta_C=0$ for every complement forces \eqref{q:eq:quotient:diagonal}, hence \eqref{p:eq:hs5:kernel-condition}: as $C$ varies, $\pi_N\Psi^2$ ranges over all of $N$, and rescaling it separates the terms linear and quadratic in it.
Requiring it for one fixed complement is a $C$-dependent statement and not a property of the theory; there is no distinguished complement, the pairing being indefinite.
The alternative of keeping all of $\Sym^2\mc H_0$ and relying on the dressed kernel, whose quintic part gives the kernel directions a quadratic form $W_{NN}$ proportional to $\Psi$, is worse: when $W_{NN}$ is invertible the kernel components of the saddle are not small, the mismatch is $-\frac{g_s^3}{5!}v_N(W_{NN}^{-1}v_N)$ with $v_N$ the restriction of $v$ to $N$, a rational function of $\Psi$ of degree five, the kernel modes have no propagator at $\Psi=0$, and the expansion around the vacuum does not exist; when $W_{NN}$ is singular the contamination starts one order earlier, at $g_s^2$.



\subsection{The two presentations of the sextic order}
\label{sec:order:sextic-presentations}

The next block is a section of an older note on the $3$-string field at the sextic order.
It compares two ways of introducing $\Sigma_3$: with the elementary sextic vertex in its kinetic term and its source, which fails, and with the dressed vertex, which closes and is the paper's presentation.
It also writes the complete cubic action through order $\gs^4$ in the presentation with $\Sigma_2$ only and in the dressed presentation, and the equations of motion.
It states two results:
\begin{itemize}
  \item within $(\Psi, \Sigma_2, \Sigma_3)$ with elementary vertices and generic tensors, the sextic order does not close: \status{Conditional(PH1, PH5, PH14 [$O(g_s^4)$]) -- Proven}\record{014-b4-sigma3-elementary (R20)};
  \item the gauge laws are identical in the dressed and the elementary presentations, and the failure of the elementary one is dynamical: \status{mixed}\record{015-b5-ogs4-gauge-eom-ibl (R21)}; the gauge laws are derived in \cref{sec:order:sextic:both-gauge}.
\end{itemize}
The note uses the older notation of the project, translated below.

\begin{quote}
  \small
  \textbf{Translation box: older note on the sextic presentations $\to$ paper.}
  \begin{itemize}
    \setlength{\itemsep}{0.3ex}
    \item $V_n$ $\to$ $\vertex{n}$; $\tfrac12 V_2(\Psi^2)$ is the kinetic term of $\Psi$; $\Hilb_0$ $\to$ $\Hs$, $\Sym^k \Hilb_0$ $\to$ $\Hn{k}$; $\Sig \equiv \Sig_2$; $\Sig_k = \Psi^{\otimes k}$ $\to$ $\Sigma_k = \Psi^k$.
    \item $V_4$ read as a bilinear form on $\Sym^2 \Hilb_0$, and its inverse $V_4^{-1}$, $\to$ $\vertex{4}(A, B) = \pair[2]{A}{\sprod{2}{2}(B)}$ and $\sprod{2}{2}^{-1}$; likewise $M = V_6 |_{\Sym^3}$ $\to$ $\sprod{3}{3}$ of~\eqref{p:eq:hs6:p33}.
    \item $S_{\rm HS}$, $S_3$ (the cubic actions through the quartic and the quintic order) $\to$ $\caction[\le 4]$, $\caction[\le 5]$; $S^{(6)}_{\rm eff}$ $\to$ $S_{\le 6}$.
    \item The composite $\mathcal C = V_5 V_4^{-1} V_5$ and the feedback $\tfrac{\gs^4}{600} \, v V_4^{-1} v$, with $v = V_5(\cdot, \Psi^{\otimes 2}, \Psi)$, $\to$ the composite vertex $\vertex{5 \cdot 5}$ of~\eqref{p:eq:hs5:exchange-vertex} and $\tfrac{\gs^4}{600} \vertex{5 \cdot 5}(\Psi^6)$.
    \item The dressed vertex $W_6 = V_6 - \tfrac65 \, \mathcal C$ $\to$ the reduced vertex $\dvertex{6}$ of~\eqref{p:eq:hs6:dressed-vertex}.
    \item The dressed kinetic form $\tilde M = M - \tfrac65 N$ of $\Sig_3$ and the dressed source $\Mean{\tilde u(\Sig_2), \Sig_3} = \tilde V_6(\Sig_3, \Sig_2, \Psi)$ $\to$ $\dvertex{6}(\Sigma_3, \Sigma_3) = \pair[3]{\Sigma_3}{\wprod{3}{3}(\Sigma_3)}$ and $\dvertex{6}(\Psi \otimes \Sigma_2, \Sigma_3)$ of~\eqref{p:eq:hs6:S6tilde}; compare the explicit form~\eqref{p:eq:hs6:S6tilde-explicit}. The map $\Phi$ has no name in the paper.
    \item \emph{Reality.} The note's conjugation is $\ddagger = (-1)^{\Ngh + 1} \star$ on each strand, so that its condition $\Sig_k^\ddagger = \Sig_k$ is the paper's $\Sigma_k^\star = (-1)^k \Sigma_k$: the two read the same for $\Sigma_2$, while for odd $k$ the sign is explicit, $\Sigma_3^\star = - \Sigma_3$, and $\Psi^\ddagger = \Psi$ is $\Psi^\star = - \Psi$ (\cref{ch:conventions}).
  \end{itemize}
\end{quote}

\subsubsection{The \texorpdfstring{$\Sig_3$}{Sigma3} presentations of the sextic: elementary
fails, dressed closes}
\label{ped:sec:milestones:sigma3}

The dressed vertex \eqref{p:eq:hs6:dressed-vertex} carries $V_4^{-1}$ \emph{inside} a three-$\Sig_2$
coupling --- an unusual object. (In this and the next subsection we write $\Sig_2\equiv\Sig$
for the bilocal auxiliary, to keep the tower index explicit.) It is natural to ask whether one can
avoid the composite altogether by introducing the \emph{next} field of the auxiliary
tower early: a three-string auxiliary
\begin{equation}
	\Sig_3\in\Sym^3\Hilb_0,\qquad \Ngh(\Sig_3)=6,
	\label{ped:eq:Sigma3def}
\end{equation}
(the tower assignment $\Ngh(\Sig_k)=2k$; $\Sig_3$ is first \emph{needed} only at
$V_7=(3{,}2{,}2)$, but nothing forbids introducing it at $O(\gs^4)$). The hope: replace the
composite $W_6(\Sig_2^3)$ by \emph{elementary} vertices only --- a bilinear kinetic term
$V_6(\Sig_3,\Sig_3)$ and a cubic source $V_6(\Sig_3,\Sig_2,\Psi)$ --- so that the
$V_4^{-1}$ inside $W_6$ is traded for a $\Sig_3$-propagator $V_6^{-1}$. Ghost saturation is
consistent: $V_6(\Sig_3,\Sig_3)$ saturates at $6{+}6=12=2\cdot6$, and
$V_6(\Sig_3,\Sig_2,\Psi)$ at $6{+}4{+}2=12$. Both presentations below are machine-checked
on generic (non-$\Loo$) tensors in a finite-dimensional model (the scripts are listed in \cref{ch:finite-models}).

\subparagraph{The elementary presentation fails.} Take
\begin{equation}
	S_{\rm elem}
	= S_3
	+ \tfrac{\gs^4}{6!}\big[-V_6(\Sig_3,\Sig_3)+2\,V_6(\Sig_3,\Sig_2,\Psi)\big],
	\label{ped:eq:Selem}
\end{equation}
with $S_3$ the quintic action, $\caction[\le 4]$ plus \eqref{p:eq:hs5:S5tilde} (the $-1{:}2$ kinetic-to-source ratio mirrors
\eqref{p:eq:hs4:S4tilde}). Since $\Sig_3$ enters only quadratically$+$linearly, it can be eliminated
exactly (Gaussian): with $M:=V_6|_{\Sym^3}$ and $u(\Sig_2):=V_6(\cdot,\Sig_2,\Psi)$, the
elimination produces the $\Sig_3$-channel exchange
$\mathcal D(\Sig_2,\Sig_2)=u^\top M^{-1}u$, \emph{quadratic} in $\Sig_2$. Three facts kill
the hope:
\begin{enumerate}[label=(\roman*),leftmargin=*]
	\setlength{\itemsep}{0.4ex}
	\item \emph{On-shell, $\mathcal D$ collapses to the bare contact.} By $\Sym^3$
	completeness and total symmetry of $V_6$, the source collapses on the Veronese diagonal:
	$u(P)=M\,p_3$ with $p_3:=\Psi^{\otimes3}$, hence
	$\mathcal D(P,P)=p_3^\top M\,p_3=V_6(\Psi^6)$ --- a telescoping $V_6V_6^{-1}V_6=V_6$,
	the exact analogue of the quartic collapse \eqref{p:eq:hs4:telescoping}. The $\Sig_3$ sector supplies precisely the
	\emph{target} sextic contact, nothing else.
	\item \emph{The feedback is a $\Sig_2$-channel object and survives.} The quintic
	overshoot $\tfrac{\gs^4}{600}vV_4^{-1}v$ is built from $V_4,V_5$ \emph{only} --- it is
	exactly $V_6$-independent, so no $V_6$-built exchange can cancel it. Structurally:
	$\mathcal D$ is quadratic in $\Sig_2$, while the composite it would have to reproduce,
	$\mathcal C=V_5V_4^{-1}V_5$, is cubic in $\Sig_2$ --- different field degrees, different
	invariants. The joint on-shell mismatch $\tfrac{\gs^4}{600}vV_4^{-1}v$ remains for
	generic tensors.
	\item \emph{No repair exists within the content.} A cubic-in-$\Sig_2$ $\Sig_3$-source
	would need $V_7(\Sig_3,\Sig_2,\Sig_2)$ --- septic, hence $O(\gs^5)$, forbidden at this
	order; and a $\Sig_3$-source at order $\gs^k$ feeds an exchange at $\gs^{2k-4}$, so only
	$k=4$ contributes at $O(\gs^4)$ at all.
\end{enumerate}
The conclusion: \emph{the composite $\mathcal C$ is
irreducible} --- no elementary $\Sig_2\oplus\Sig_3$ presentation closes the $O(\gs^4)$
matching; the $V_4^{-1}\to V_6^{-1}$ trade is not an identity. (A new auxiliarity condition
appears alongside the one for $\Sig_2$: $\Sig_3$ is purely auxiliary iff $V_6$ is non-degenerate
on $\Sym^3\Hilb_0$.)

\subparagraph{The dressed presentation closes.} The cure is not to reproduce the \emph{bare} $V_6(\Psi^6)$ with $\Sig_3$ but
to pair $\Sig_3$ with the \emph{dressed effective sextic} --- the vertex that already
contains the lower-order feedback. Define the linear map
\begin{equation}
	\Phi:\ \Sym^3\Hilb_0\to\big(\Sym^2\Hilb_0\big)^*,
	\qquad
	\Phi(X):=V_5(\,\cdot\,,X)\big|_{\Sym^2},
	\label{ped:eq:Phimap}
\end{equation}
i.e.\ contract $V_5$ on three of its five slots with $X\in\Sym^3$ and read the remaining
two slots as a covector on $\Sym^2$ (well defined by total symmetry of $V_5$). Its
key property is the \emph{Veronese collapse}
\begin{equation}
	\Phi(\Psi^{\otimes3})=v
	\qquad\big(v=V_5(\cdot,\Psi^{\otimes2},\Psi)\text{, the quintic source covector}\big):
	\label{ped:eq:Phicollapse}
\end{equation}
on the diagonal, splitting three $\Psi$-slots as $2{+}1$ is immaterial. Then set
\begin{equation}
	\boxed{\;
	\tilde M:=M-\tfrac65\,N,
	\qquad
	N:=\Phi^\top V_4^{-1}\Phi,
	\qquad
	\tilde u(\Sig_2):=u(\Sig_2)-\tfrac65\,\Phi^\top V_4^{-1}\,v(\Sig_2),
	\;}
	\label{ped:eq:tildeMdef}
\end{equation}
with $v(\Sig_2):=V_5(\cdot,\Sig_2,\Psi)$ the off-shell extension of $v$. Here $M$,
$\tilde M$ are bilinear forms on $\Sym^3\Hilb_0$, while $u(\Sig_2)$, $\tilde u(\Sig_2)$
are \emph{covectors} on $\Sym^3\Hilb_0$ (linear functionals of the would-be $\Sig_3$
argument), themselves linear in $\Sig_2$; we write $\Mean{\tilde u(\Sig_2),\Sig_3}$ for
the evaluation. Explicitly, in an orthonormal basis $\{\Lambda^2_I\}$ of $\Sym^2$:
\begin{equation}
	\begin{aligned}
	\Mean{\tilde u(\Sig_2),\Sig_3}
	&= V_6(\Sig_3,\Sig_2,\Psi)
	- \tfrac65\sum_{I,J} V_5(\Lambda^2_I,\Sig_3)\,(V_4^{-1})_{IJ}\,V_5(\Lambda^2_J,\Sig_2,\Psi)
	\\&=: \tilde V_6(\Sig_3,\Sig_2,\Psi),
	\end{aligned}
	\label{ped:eq:utildeexplicit}
\end{equation}
the first term the bare source, the second the $\Phi$-map bridge: one $V_5$ eats the
three strands of $\Sig_3$, the other eats $\Sig_2$ and $\Psi$, and a $\Sig_2$-propagator
$V_4^{-1}$ joins them. The action is
\begin{equation}
	S_{\rm dr}
	= S_3
	+ \tfrac{\gs^4}{6!}\big[-\tilde M(\Sig_3,\Sig_3)+2\,\Mean{\tilde u(\Sig_2),\Sig_3}\big].
	\label{ped:eq:Sdr}
\end{equation}
Both $\tilde M$ and $N$ are $\Psi$-\emph{independent} (built from $V_4,V_5,V_6$ and fixed
bases alone) and symmetric; this will matter for the gauge law. Two identities, both direct
consequences of \eqref{ped:eq:Phicollapse}, do all the work:
\begin{equation}
	\tilde M(p_3,p_3)
	= V_6(\Psi^6)-\tfrac65\,vV_4^{-1}v
	= W_6(\Psi^6),
	\qquad
	\tilde u(P)=\tilde M(p_3,\cdot).
	\label{ped:eq:tildecollapse}
\end{equation}
The second is an equality of covectors on $\Sym^3$: $\tilde u(P)$ means \enquote{the
source covector evaluated on the $\Sig_2$-saddle $\Sig_2=P=\Psi^{\otimes2}$}, i.e.\
$\Mean{\tilde u(P),X}=\tilde V_6(X,\Psi^{\otimes2},\Psi)$ for every $X\in\Sym^3$, and the
identity says this functional coincides with pairing $X$ against $\Psi^{\otimes3}$
through the dressed kinetic form, $\tilde M(p_3,X)$. It is proved in one line from
\eqref{ped:eq:Phicollapse}: $u(P)=M(p_3,\cdot)$ is the bare collapse of the elementary presentation (total symmetry of
$V_6$: splitting six $\Psi$-slots as $3{+}2{+}1$ or $3{+}3$ is immaterial), and
$\Phi^\top V_4^{-1}v\big|_{\Sig_2=P}=N(p_3,\cdot)$ by $\Phi(p_3)=v$. Consequently the
dressed saddle is \emph{exactly} $\Sig_3^*=\Psi^{\otimes3}$ on-shell, and the resulting
$\Sig_3$-contact is $\tfrac{\gs^4}{6!}\tilde M(p_3,p_3)=\tfrac{\gs^4}{6!}W_6(\Psi^6)$ ---
the geometric sextic \emph{and} the feedback subtraction, in one Gaussian step. The
matching closes exactly as in \cref{p:sec:hs:sextic}, with the composite
$-\tfrac65\mathcal C$ now living \emph{inside the $\Sig_3$-kinetic pairing} instead of
inside a three-$\Sig_2$ vertex: the map $\Phi$ is the relocation mechanism, converting the
$\Sig_2$-channel feedback $vV_4^{-1}v$ (a $\Sym^2$-channel object) into a $\Sym^3$-kinetic
correction $N$. The composite is not eliminated --- the elementary presentation shows it cannot be ---
it is \emph{repackaged} where it does the least harm: in a quadratic kernel, the natural
home of a propagator-like object, rather than in a cubic vertex.

\subparagraph{The complete cubic action through $O(\gs^4)$, both presentations.} For
reference:
\begin{equation}
	\boxed{
	\begin{aligned}
	S^{\rm min}_4
	&= \tfrac12V_2(\Psi^2)+\tfrac{\gs}{3!}V_3(\Psi^3)
	-\tfrac{\gs^2}{4!}V_4(\Sig_2,\Sig_2)+\tfrac{2\gs^2}{4!}V_4(\Sig_2,\Psi^{\otimes2})
	\\&\quad
	+\tfrac{\gs^3}{5!}V_5(\Sig_2,\Sig_2,\Psi)
	+\tfrac{\gs^4}{6!}\,W_6(\Sig_2,\Sig_2,\Sig_2),
	\\[1ex]
	S^{\rm dr}_4
	&= \tfrac12V_2(\Psi^2)+\tfrac{\gs}{3!}V_3(\Psi^3)
	-\tfrac{\gs^2}{4!}V_4(\Sig_2,\Sig_2)+\tfrac{2\gs^2}{4!}V_4(\Sig_2,\Psi^{\otimes2})
	\\&\quad
	+\tfrac{\gs^3}{5!}V_5(\Sig_2,\Sig_2,\Psi)
	-\tfrac{\gs^4}{6!}\tilde M(\Sig_3,\Sig_3)+\tfrac{2\gs^4}{6!}\Mean{\tilde u(\Sig_2),\Sig_3}.
	\end{aligned}
	}
	\label{ped:eq:S4both}
\end{equation}
Both are cubic in the tower fields $(\Psi,\Sig_2,\Sig_3)$, agree exactly with
$S^{(6)}_{\rm eff}$ on-shell up to $O(\gs^5)$ feedback, and carry the homogeneous gauge
laws of \cref{sec:order:sextic:both-gauge}. The minimal presentation is more economical (one field
fewer); the dressed presentation is structurally uniform --- every order adds one field and
one HS pair (kinetic $+$ source), which is what generalises
(\cref{ch:all-orders}).

\subparagraph{Equations of motion and the auxiliarity chain.} In both presentations the
$\Sig_3$-EOM is Gaussian --- $\Sig_3$ enters bilinearly $+$ linearly --- with exact saddle
$\Sig_3^*=\tilde M^{-1}\tilde u(\Sig_2)$ (dressed; bare with $M,u$), collapsing to
$\Psi^{\otimes3}$ on-shell by \eqref{ped:eq:tildecollapse}. The $\Sig_2$-EOM is the
\emph{implicit} saddle equation ($\tilde X_6=0$ with the $\Sig_3$-exchange
folded into the Hessian): the $\Sig_2$-saddle is a perturbative fixed point
$\Sig_2^*(\Psi)=\Psi^{\otimes2}+O(\gs)$, Neumann-solvable at small $\gs$. The auxiliarity
conditions chain up: $\Sig_2$ is purely auxiliary iff $V_4|_{\Sym^2}$ is non-degenerate
(\cref{p:sec:hs:quartic}); $\Sig_3$ iff the relevant kinetic kernel is --- $M=V_6|_{\Sym^3}$
(elementary) or $\tilde M$ (dressed; this needs \emph{both} $V_4$ and $V_6$ well behaved,
tying the two conditions together). Reality extends verbatim from
\cref{p:eq:hs4:reality} (in the sign convention of the translation box): $\Sig_3^\ddagger=\Sig_3$, inherited from
$(\Psi^{\otimes3})^\ddagger=\Psi^{\otimes3}$, makes every new term real.



\subsubsection{Gauge laws and products of the sextic order in both presentations}
\label{sec:order:sextic:both-gauge}

\status{mixed}\record{015-b5-ogs4-gauge-eom-ibl (R21)}
The gauge laws of both presentations and the identity of their structure \status{Conditional(PH1, PH2 (b), PH5) + R70 -- Proven}; the statement that the residual of the elementary presentation is dynamical and not removable at the string vertices \status{Conditional(PH1, PH2 (b), PH5, PH14 [$O(g_s^4)$]) + R70 -- Proven}; the equations of motion and the chain of auxiliarity conditions \status{Proven}; the identification of the products and their ghost numbers \status{Proven}.

The paper gives the gauge law of the sextic order in the dressed presentation (\cref{p:sec:hs:sextic:gauge}).
We derive it here for both presentations at once, in the notation of the block above, and show that the two have the same gauge structure: the failure of the elementary presentation lies in the matching, not in the gauge invariance.
The setting is the algebraic one of the block above: generic finite-dimensional tensors, tree level, order $\gs^4$.
We assume that $V_4$ is invertible on the quotient of $\Sym^2 \Hilb_0$ by its kernel, that $V_5$ and $V_6$ annihilate that kernel, and the condition on the null vectors of the dressed kinetic form of $\Sig_3$ among the kernel conditions of the sextic order (\cref{p:sec:hs:sextic:kernel}); the elementary presentation needs in addition $M = V_6 |_{\Sym^3}$ invertible, and the statement that its residual cannot be removed needs the feedback to be generic, not expressible through the vertices $V_n(\Psi^n)$ and boundary terms.
The analytic properties of the inverses $\tilde M^{-1}$, $M^{-1}$ and $V_4^{-1}$ at the level of the conformal field theory are not covered.

\paragraph{The two presentations.}
Both presentations share the action through the quintic order.
They differ in the sector of $\Sig_3$, which we write with a kinetic form $K_3$ and a source $s(\Sig_2)$, linear in $\Sig_2$:
\begin{equation}
  \label{eq:ch2:sext:both}
  S
    = S_3 + \frac{\gs^4}{6!} \Big[ - K_3(\Sig_3, \Sig_3) + 2 \, \Mean{s(\Sig_2), \Sig_3} \Big],
  \qquad
  (K_3, s)
    =
    \begin{cases}
      (\tilde M, \tilde u)
        & \text{dressed},
      \\
      (M, u)
        & \text{elementary}.
    \end{cases}
\end{equation}
In both cases $K_3$ does not depend on $\Psi$: $\tilde M$ and $M$ are built from $V_4$, $V_5$, $V_6$ alone.
The elementary presentation may carry in addition a coupling $\tfrac{\gs^4}{6!} c_6 \, V_6(\Sig_2, \Sig_2, \Sig_2)$, to test whether it can replace the missing subtraction; the block above takes $c_6 = 0$, and so does the dressed presentation, in which a term $c_6 V_6(\Sig_2^3)$ would count the contact term twice, once from the collapse of $\Sig_3$ and once from the coupling.
The saddle of $\Sig_3$ is
\begin{equation}
  \Sig_3^*
    = K_3^{-1} s(\Sig_2),
\end{equation}
and on-shell, $\Sig_2 = \Psi^{\otimes 2}$, the collapse $s(\Psi^{\otimes 2}) = K_3 \, \Psi^{\otimes 3}$ of the block above gives $\Sig_3^* = \Psi^{\otimes 3}$ and the contact term $\tfrac{\gs^4}{6!} K_3(\Psi^{\otimes 3}, \Psi^{\otimes 3})$, that is $\tfrac{\gs^4}{6!} W_6(\Psi^6)$ in the dressed presentation and $\tfrac{\gs^4}{6!} V_6(\Psi^6)$ in the elementary one.

\paragraph{The gauge transformation of $\Sig_3$.}
The string field transforms with the gauge transformation of closed string field theory truncated at this order,
\begin{equation}
  \delta \Psi
    = Q_B \Lambda
      + \gs \, \ell_2(\Psi, \Lambda)
      + \frac{\gs^2}{2} \, \ell_3(\Psi, \Psi, \Lambda)
      + \frac{\gs^3}{6} \, \ell_4(\Psi, \Psi, \Psi, \Lambda)
      + \frac{\gs^4}{24} \, \ell_5(\Psi^{\otimes 4}, \Lambda),
\end{equation}
under which $\delta S_{\le 6} = O(\gs^5)$.
We look for $\delta \Sig_2$ and $\delta \Sig_3$ such that
\begin{equation}
  \label{eq:ch2:sext:goal}
  \delta S - \delta S_{\le 6}
    = \delta \mathbb F
\end{equation}
exactly, off-shell and for arbitrary $\delta \Psi$, with $\mathbb F$ independent of the auxiliary fields.
Let $Y_3 := \Sig_3 - \Sig_3^*(\Psi, \Sig_2)$ be the deviation of $\Sig_3$ from its saddle.
Completing the square in~\eqref{eq:ch2:sext:both}, the action depends on $Y_3$ only through $- \tfrac{\gs^4}{6!} K_3(Y_3, Y_3)$, whose variation is
\begin{equation}
  \delta \Big[ - \frac{\gs^4}{6!} \, K_3(Y_3, Y_3) \Big]
    = - \frac{2 \gs^4}{6!} \, K_3(\delta Y_3, Y_3),
\end{equation}
since $K_3$ does not depend on $\Psi$.
For the part of $\delta S$ that depends on the auxiliary fields to cancel for every $Y_3$, we need $K_3(\delta Y_3, Y_3) = 0$ for all $Y_3$, hence $\delta Y_3 = 0$, $K_3$ being non-degenerate.
The $3$-string field thus transforms as its saddle,
\begin{equation}
  \label{eq:ch2:sext:gauge3}
  \delta Y_3
    = 0,
  \qquad
  \delta \Sig_3
    = \delta \Sig_3^*(\Psi, \Sig_2)
    = K_3^{-1} \Big[ (\partial_\Psi s) \, \delta \Psi + (\partial_{\Sig_2} s) \, \delta \Sig_2 \Big],
\end{equation}
in both presentations, with no rotation of the deviation and no inhomogeneous term.
This is checked numerically in a finite-dimensional model by three tests: the law agrees with a finite difference of the saddle map computed from the sources alone; the variations of the action of $\Sig_3$ at two different random $Y_3$ are equal; and replacing $K_3$ by a kernel that depends on $\Psi$ breaks the second test by exactly the predicted residual, which a rotation of $Y_3$ then removes.

\paragraph{The gauge transformation of $\Sig_2$.}
Eliminating $\Sig_3$ first, by its exact Gaussian, adds to the action through the quintic order a term quadratic in $\Sig_2$, the exchange of $\Sig_3$:
\begin{equation}
  S \big|_{\Sig_3^*}
    = S_3 + \frac{\gs^4}{6!} \, \mathcal E(\Sig_2, \Psi),
  \qquad
  \mathcal E
    = s^\top K_3^{-1} s,
\end{equation}
that is $\tilde u^\top \tilde M^{-1} \tilde u$ in the dressed presentation and $u^\top M^{-1} u = \mathcal D$ in the elementary one.
Around the joint saddle $\Sig_2^*(\Psi)$, with $Y_2 := \Sig_2 - \Sig_2^*(\Psi)$, the action takes the form
\begin{equation}
  S \big|_{\Sig_3^*}
    = S_{\le 6} + \Phi_2(Y_2; \Psi) + \mathbb F(\Psi),
  \qquad
  \Phi_2
    = - \frac{\gs^2}{4!} \, \mathcal K(Y_2, Y_2) + \frac{\gs^4}{6!} \, T(Y_2, Y_2, Y_2),
\end{equation}
with $\mathcal K$ the Hessian at the saddle, which includes the contribution of $\mathcal E$, and $T$ the cubic term.
In the dressed presentation $T = 0$ at this order: there is no coupling of three $\Sig_2$, and $\mathcal E$ is quadratic in $\Sig_2$, since $s$ is linear in it, hence quadratic in $Y_2$, so it contributes to $\mathcal K$ only.
In the elementary presentation $T = c_6 V_6$ when $c_6 \ne 0$, and $\mathcal E$ again contributes to $\mathcal K$ only.
Requiring $\delta \Phi_2 = 0$ for arbitrary $\delta \Psi$ gives the implicit law
\begin{equation}
  \label{eq:ch2:sext:gauge2}
  \delta Y_2
    = \mathcal K^{-1} \Big[ - \frac12 \, (\delta \mathcal K)(\cdot, Y_2) + \frac{\gs^2}{20} \, T(\cdot, Y_2, \delta Y_2) \Big],
  \qquad
  \delta \Sig_2
    = \delta \big[ \Sig_2^*(\Psi) \big] + \delta Y_2,
\end{equation}
which is solved by a Neumann series for small $\gs$ and is homogeneous, $\delta Y_2 |_{Y_2 = 0} = 0$.
In the dressed presentation it reduces to the invariance of the quadratic form $\mathcal K$, with no cubic correction.
With $\delta Y_3 = 0$, the equality~\eqref{eq:ch2:sext:goal} holds exactly, with
\begin{equation}
  \label{eq:ch2:sext:F}
  \mathbb F
    =
    \begin{cases}
      O(\gs^5)
        & \text{dressed},
      \\
      \dfrac{\gs^4}{600} \, v V_4^{-1} v + \dfrac{\gs^4}{6!} \, c_6 \, V_6(\Psi^6) + O(\gs^5)
        & \text{elementary}.
    \end{cases}
\end{equation}
The off-shell exactness of~\eqref{eq:ch2:sext:goal} and the law~\eqref{eq:ch2:sext:gauge2} are not checked numerically here.
They are inherited from the explicit expansion around the saddle in the presentation with $\Sig_2$ only, where the off-shell equality was verified by hand: the actions of both presentations differ from that one only by the exchange $\mathcal E$, which is quadratic in $Y_2$ and Gaussian in $\Sig_3$, so that the structure, a quadratic form in $Y_2$ plus a remainder independent of the auxiliary fields, is preserved term by term.
What is checked numerically is the on-shell value of $\mathbb F$ at order $\gs^4$, which extrapolates to zero in the dressed presentation and to the residual~\eqref{eq:ch2:sext:F} in the elementary one, and the stationarity of the saddle.

\paragraph{The comparison.}
The two presentations side by side:
\begin{center}
  \small
  \begin{tabular}{lll}
    \toprule
                                                & dressed                                          & elementary                               \\
    \midrule
    kinetic form $K_3$ of $\Sig_3$              & $\tilde M = M - \tfrac65 N$, composite           & $M$, bare                                \\
    $\delta Y_3$                                & $0$                                              & $0$                                      \\
    law of $\Sig_2$                             & implicit, $T = 0$                                & implicit, $T = c_6 V_6$                  \\
    $\delta S - \delta S_{\le 6}$               & $\delta \mathbb F$, $\mathbb F = O(\gs^5)$       & $\delta \mathbb F$, $\mathbb F = O(\gs^4) \ne 0$ \\
    independent $2$-string gauge parameter      & none                                             & none                                     \\
    \bottomrule
  \end{tabular}
\end{center}
The gauge structure is the same: $\Sig_3$ transforms as its saddle, $\Sig_2$ by an implicit homogeneous law, and neither presentation has an independent $2$-string gauge parameter.
The elementary presentation does not fail for lack of a gauge law: the law exists, is homogeneous and gives the exact equality~\eqref{eq:ch2:sext:goal}.
It fails because $\mathbb F$ is non-zero at order $\gs^4$: its on-shell contact term of $\Sig_3$ is the bare $V_6(\Psi^6)$ instead of $W_6(\Psi^6)$, so the feedback $\tfrac{\gs^4}{600} v V_4^{-1} v$ of the quintic order is not subtracted.
The gauge law does not distinguish the composite kinetic form from the elementary one; only the matching does, and the obstruction lies in the target action, not in the gauge invariance.

\paragraph{The products at order $\gs^4$.}
In the language of the products of \cref{ch:ibl}, the action at this order uses, at genus zero, the following maps with $k$ inputs and $\ell$ outputs, whose physical ghost number is $3 \ell - 2 k$:
\begin{center}
  \small
  \begin{tabular}{llll}
    \toprule
    product                          & built from                                  & $\Ngh$ & enters through                             \\
    \midrule
    $\ell_4$, $4$ in, $1$ out         & $\vertex{5}$                                & $-5$   & the quintic coupling and the dressing      \\
    $\ell_5$, $5$ in, $1$ out         & $\vertex{6}$                                & $-7$   & the sextic contact term                    \\
    $\sprod{3}{2}$, $3$ in, $2$ out   & $\vertex{5}$                                & $0$    & the two-output product of $\vertex{5}$     \\
    $\sprod{4}{2}$, $4$ in, $2$ out   & $\vertex{6}$                                & $-2$   & the two-output product of $\vertex{6}$     \\
    $\tilde p_{4, 2}$, $4$ in, $2$ out & the dressed vertex                          & $-2$   & the sector of $\Sig_3$ in the dressed form \\
    \bottomrule
  \end{tabular}
\end{center}
The dressed product is not the two-output product of a vertex but a composite of lower products: the composite $\mathcal C = V_5 V_4^{-1} V_5$ has the shape of two products $\sprod{3}{2}$ joined by the propagator $\sprod{2}{2}^{-1}$ of $\Sig_2$, so that
\begin{equation}
  \tilde p_{4, 2}
    = \sprod{4}{2} - \frac65 \, \big( \text{composition of } \sprod{3}{2}, \sprod{2}{2}^{-1}, \sprod{3}{2} \big).
\end{equation}
It is the first two-output product that is not obtained from a single vertex, and it is the counterpart among the products of the reduced vertex $W_6 = V_6 - \tfrac65 \mathcal C$.
Of the relations among these products, those with one output are the $L_\infty$ relations through $\ell_4$ and $\ell_5$:
\begin{equation}
  \com{Q, \ell_4} + \com{\ell_2, \ell_3}
    = 0,
  \qquad
  \com{Q, \ell_5} + \com{\ell_2, \ell_4} + \frac12 \, \com{\ell_3, \ell_3}
    = 0,
\end{equation}
of ghost numbers $-4$ and $-6$, which the vertices of closed string field theory satisfy.
The relations with two outputs, at signatures $(3, 2)$ and $(4, 2)$, do not follow from the $L_\infty$ relations in their augmented form: the augmented $(3, 2)$ relation fails on the physical sector for a cyclic $L_\infty$ algebra, and every augmented relation with at least two inputs and two outputs fails for some cyclic $L_\infty$ algebra.
What holds are the relations with one input or one output and the non-augmented family, which \cref{ch:ibl} states.

The content of this section is that the choice between the elementary and the dressed $3$-string field is invisible to the gauge symmetry and visible only in the matching: the gauge law follows the saddle in both cases, and only the on-shell value of the contact term of $\Sig_3$ decides whether the feedback of the quintic order is subtracted.
